% Dissertation : rédiger une dissertation — Français 1re Tech — Cours signé OnBuch+
% Programme officiel MINESEC. Compiler avec : tectonic N03-dissertation-rediger-dissertation.tex
\newcommand{\DOCMATIERE}{Français}
\newcommand{\DOCNIVEAU}{1re Tech}
\newcommand{\DOCMODULE}{Français – classe de Première technique}
\newcommand{\DOCLECON}{N03}
\newcommand{\DOCTITRE}{Dissertation : rédiger une dissertation}
% OnBuch+ — préambule « cours signés » ENRICHI (compilé avec tectonic / XeLaTeX)
% Système visuel « Pop » d'OnBuch+ : fond crème (jamais blanc), bordures encre
% épaisses, ombres « dures » décalées, titres Archivo Black en capitales.
% Version MATHÉMATIQUES : graphes (pgfplots/tikz), boîtes pédagogiques
% (définition, propriété, méthode, attention, le savais-tu, exercice résolu,
% exercice, TP, bilan), page de garde et signature OnBuch+ ancrée en bas.
%
% Macros attendues dans main.tex AVANT \input{preamble} :
%   \DOCMATIERE  \DOCNIVEAU  \DOCMODULE  \DOCLECON (numéro)  \DOCTITRE
\documentclass[11pt]{article}

\usepackage{fontspec}
\usepackage[french]{babel}
\usepackage[a4paper,margin=2cm,top=3.4cm,bottom=2.8cm,headheight=1.4cm,footskip=1.3cm]{geometry}
\usepackage{amsmath,amssymb,amsfonts}
\usepackage{mathtools} % \xmapsto, \coloneqq... souvent utilisés par les modèles
\usepackage{cancel} % simplifications barrées (bilans de forces)
% \abs{z} (module, valeur absolue) et \norm{u} (norme), utilisés par les modèles
\providecommand{\abs}{}\let\abs\relax\DeclarePairedDelimiter{\abs}{\lvert}{\rvert}
\providecommand{\norm}{}\let\norm\relax\DeclarePairedDelimiter{\norm}{\lVert}{\rVert}
\usepackage[table]{xcolor} % [table] : fournit aussi \rowcolors (alternance de lignes)
\usepackage{pagecolor}
\usepackage{tikz}
\usetikzlibrary{arrows.meta,positioning,calc,shapes.geometric,shapes.misc,shapes.arrows,decorations.pathmorphing,decorations.pathreplacing,decorations.markings,patterns,shadows,mindmap,trees,fit,backgrounds,matrix,intersections,angles,quotes}
\usepackage{pgfplots}
\usepackage[version=4]{mhchem} % équations nucléaires \ce{^{235}_{92}U}
\usepackage[european,siunitx]{circuitikz} % schémas électriques
\pgfplotsset{compat=1.18}
\usepgfplotslibrary{fillbetween,groupplots} % aires entre courbes (calcul intégral)
\usepackage{enumitem}
\usepackage{fancyhdr}
% [most] charge toutes les bibliothèques tcolorbox (listings, minted, raster,
% documentation...) et gonfle inutilement la mémoire TeX sur un document long
% (~300 boîtes) : on ne charge que ce qui est réellement utilisé.
\usepackage[skins,breakable]{tcolorbox}
\usepackage{graphicx}
\usepackage{colortbl}
\usepackage{booktabs}
\usepackage{tabularx}
\usepackage{multirow}
\usepackage{diagbox} % case d'en-tête diagonale des tableaux à double entrée
% Colonnes extensibles centrée / gauche / droite (C, L, R), souvent utilisées par les modèles
\newcolumntype{C}{>{\centering\arraybackslash}X}
\newcolumntype{L}{>{\raggedright\arraybackslash}X}
\newcolumntype{R}{>{\raggedleft\arraybackslash}X}
\usepackage{array}
\usepackage{multicol}
\usepackage{siunitx}
% parse-numbers=false : accepte aussi \SI{2,5 \times 10^{-3}}{...}
\sisetup{locale=FR,output-decimal-marker={,},group-minimum-digits=4,parse-numbers=false}
\DeclareSIUnit\ohm{\ensuremath{\symup{\Omega}}} % U+2126 absent de la police
\DeclareSIUnit\division{div} % réglages d'oscilloscope (ms/div, V/div)
\DeclareSIUnit\tr{tr} % tours (vitesse de rotation, ex. \SI{1500}{\tr\per\minute})
\DeclareSIUnit\molar{M} % concentration molaire (ex. \SI{3}{\molar}, \SI{1}{\micro\molar})
\DeclareSIUnit\an{an} % année (le modèle confond parfois avec \year, un compteur TeX) — ex. \SI{5}{\kilo\meter\per\an}
\DeclareSIUnit\FCFA{FCFA} % franc CFA (ex. \SI{500}{\FCFA\per\kilo\gram})
\DeclareSIUnit\degreeBrix{\textdegree Brix} % teneur en sucre d'un sirop/jus (réfractométrie)
\DeclareSIUnit\month{mois} % le modèle utilise parfois \month, un compteur TeX, comme unité de durée
\DeclareSIUnit\microliter{\micro\liter} % le modèle écrit parfois \microliter en un seul token
\DeclareSIUnit\million{millions} % dénombrement (ex. \SI{15}{\million\per\mL} spermatozoïdes)
\usepackage{qrcode}
% \degree : commande fréquemment utilisée par les modèles pour un angle ou une
% température en dehors de \SI{}{} (non fournie par les paquets chargés).
\providecommand{\degree}{^{\circ}}
% \qed (fin de démonstration), utilisé par les modèles sans amsthm
\providecommand{\qed}{\hfill$\square$}
% \barre : double barre des valeurs interdites dans un tableau de variations (tkz-tab)
\providecommand{\barre}{\ensuremath{\|}}
% environnement proof (amsthm non chargé) utilisé par les modèles
\ifdefined\proof\else\newenvironment{proof}[1][Démonstration]{\par\noindent\textit{#1.}\ }{\qed\par}\fi
\usepackage{lastpage}
\usepackage{hyperref}
\hypersetup{hidelinks}

% ============================================================
% Palette « Pop » OnBuch+ — mêmes hex que la classe P (plus_ui.dart)
% ============================================================
\definecolor{poporange}{HTML}{F59321}
\definecolor{popdark}{HTML}{E07A0C}
\definecolor{popcream}{HTML}{FFF6E8}
\definecolor{popink}{HTML}{1C1714}
\definecolor{poppurple}{HTML}{7A5AE0}
\definecolor{popgreen}{HTML}{2E9E6A}
\definecolor{poppink}{HTML}{E0457B}
\definecolor{popblue}{HTML}{2F7DE1}
\definecolor{popgold}{HTML}{C98A12}
\definecolor{popmuted}{HTML}{8A7F76}
\definecolor{popline}{HTML}{EADFD2}
\definecolor{popgray}{HTML}{8A7F76}
\colorlet{popgrey}{popgray}
% Alias tolérants (le modèle nomme parfois les couleurs différemment)
\colorlet{popwhite}{white}
\colorlet{popblack}{popink}
\colorlet{popred}{poppink}
\colorlet{popyellow}{popgold}
% Teintes claires (fonds de tableaux/encadrés) — le modèle nomme parfois
% les couleurs avec un suffixe L (ex. popblueL) sans les avoir définies.
\colorlet{poporangeL}{poporange!20}
\colorlet{popdarkL}{popdark!20}
\colorlet{poppurpleL}{poppurple!20}
\colorlet{popgreenL}{popgreen!20}
\colorlet{poppinkL}{poppink!20}
\colorlet{popblueL}{popblue!20}
\colorlet{popgoldL}{popgold!20}
\colorlet{popredL}{popred!20}
\colorlet{popgrayL}{popgray!20}
% Teintes claires pour fonds de tableaux / zones
\colorlet{popblueL}{popblue!10!white}
\colorlet{poporangeL}{poporange!14!white}
\colorlet{popgreenL}{popgreen!12!white}
\colorlet{poppurpleL}{poppurple!12!white}
\colorlet{poppinkL}{poppink!12!white}
\colorlet{popgoldL}{popgold!14!white}

\pagecolor{popcream}

% ============================================================
% Typographie
% ============================================================
\defaultfontfeatures{Ligatures=TeX}
\setmainfont{PlusJakartaSans-Medium.ttf}[Path=../../../fonts/,BoldFont=PlusJakartaSans-Bold.ttf]
% Maths en sans-serif (Fira Math) avec chiffres et lettres droites de Plus
% Jakarta, pour que les formules se fondent dans le texte.
\usepackage[math-style=ISO,bold-style=ISO]{unicode-math}
\setmathfont{FiraMath-Regular.otf}[Path=../../../fonts/]
\setmathfont{PlusJakartaSans-Medium.ttf}[Path=../../../fonts/,range={up/num,up/{Latin,latin},\mathup}]
\usepackage{icomma} % virgule décimale sans espace en mode math
% Caractères mathématiques Unicode absents des polices de texte (Plus Jakarta,
% Archivo Black) : rendus par leur équivalent mathématique, en texte comme en
% formule — sinon ils s'affichent en carré vide (ex. « Arithmétique dans ℤ »).
\usepackage{newunicodechar}
\newunicodechar{ℝ}{\ensuremath{\mathbb{R}}}
\newunicodechar{ℤ}{\ensuremath{\mathbb{Z}}}
\newunicodechar{ℂ}{\ensuremath{\mathbb{C}}}
\newunicodechar{ℕ}{\ensuremath{\mathbb{N}}}
\newunicodechar{ℚ}{\ensuremath{\mathbb{Q}}}
\newunicodechar{𝔻}{\ensuremath{\mathbb{D}}}
\newunicodechar{α}{\ensuremath{\alpha}}
\newunicodechar{β}{\ensuremath{\beta}}
\newunicodechar{γ}{\ensuremath{\gamma}}
\newunicodechar{δ}{\ensuremath{\delta}}
\newunicodechar{ε}{\ensuremath{\varepsilon}}
\newunicodechar{θ}{\ensuremath{\theta}}
\newunicodechar{λ}{\ensuremath{\lambda}}
\newunicodechar{μ}{\ensuremath{\mu}}
\newunicodechar{σ}{\ensuremath{\sigma}}
\newunicodechar{τ}{\ensuremath{\tau}}
\newunicodechar{φ}{\ensuremath{\varphi}}
\newunicodechar{ψ}{\ensuremath{\psi}}
\newunicodechar{ω}{\ensuremath{\omega}}
\newunicodechar{Σ}{\ensuremath{\Sigma}}
% majuscules produites par \MakeUppercase dans les titres (α-aminés -> Α)
\newunicodechar{Α}{\ensuremath{\alpha}}
\newunicodechar{Β}{\ensuremath{\beta}}
\newunicodechar{Γ}{\ensuremath{\Gamma}}
\newunicodechar{Δ}{\ensuremath{\Delta}}
\newunicodechar{Ε}{\ensuremath{\varepsilon}}
\newunicodechar{Θ}{\ensuremath{\Theta}}
\newunicodechar{Λ}{\ensuremath{\Lambda}}
\newunicodechar{Μ}{\ensuremath{\mu}}
\newunicodechar{Τ}{\ensuremath{\tau}}
\newunicodechar{Φ}{\ensuremath{\Phi}}
\newunicodechar{Ψ}{\ensuremath{\Psi}}
\newunicodechar{Ω}{\ensuremath{\Omega}}
\newunicodechar{↦}{\ensuremath{\mapsto}}
\newunicodechar{∈}{\ensuremath{\in}}
\newunicodechar{∉}{\ensuremath{\notin}}
\newunicodechar{∧}{\ensuremath{\wedge}}
\newunicodechar{⇔}{\ensuremath{\Leftrightarrow}}
\newunicodechar{⟺}{\ensuremath{\Longleftrightarrow}}
\newunicodechar{⇒}{\ensuremath{\Rightarrow}}
\newunicodechar{⟹}{\ensuremath{\Longrightarrow}}
\newunicodechar{∘}{\ensuremath{\circ}}
\newunicodechar{∪}{\ensuremath{\cup}}
\newunicodechar{∩}{\ensuremath{\cap}}
\newunicodechar{≡}{\ensuremath{\equiv}}
\newunicodechar{⊂}{\ensuremath{\subset}}
\newunicodechar{⊥}{\ensuremath{\perp}}
\newunicodechar{∃}{\ensuremath{\exists}}
\newunicodechar{∀}{\ensuremath{\forall}}
\newunicodechar{∛}{\ensuremath{\sqrt[3]{\,}}}
\newunicodechar{∜}{\ensuremath{\sqrt[4]{\,}}}
\newunicodechar{⃗}{\ensuremath{{}^{\to}}}
\newunicodechar{ⁿ}{\ifmmode^{n}\else\textsuperscript{\ensuremath{n}}\fi}
\newunicodechar{ˣ}{\ifmmode^{x}\else\textsuperscript{\ensuremath{x}}\fi}
\newunicodechar{ᵇ}{\ifmmode^{b}\else\textsuperscript{\ensuremath{b}}\fi}
\newunicodechar{ʳ}{\ifmmode^{r}\else\textsuperscript{\ensuremath{r}}\fi}
\newunicodechar{ᵘ}{\ifmmode^{u}\else\textsuperscript{\ensuremath{u}}\fi}
\newunicodechar{ʸ}{\ifmmode^{y}\else\textsuperscript{\ensuremath{y}}\fi}
\newunicodechar{ᵃ}{\ifmmode^{a}\else\textsuperscript{\ensuremath{a}}\fi}
\newunicodechar{ᵉ}{\ifmmode^{e}\else\textsuperscript{\ensuremath{e}}\fi}
\newunicodechar{ᵏ}{\ifmmode^{k}\else\textsuperscript{\ensuremath{k}}\fi}
\newunicodechar{ᵖ}{\ifmmode^{p}\else\textsuperscript{\ensuremath{p}}\fi}
\newunicodechar{ᴺ}{\ifmmode^{N}\else\textsuperscript{\ensuremath{N}}\fi}
\newunicodechar{ᶜ}{\ifmmode^{c}\else\textsuperscript{\ensuremath{c}}\fi}
\newunicodechar{ʰ}{\ifmmode^{h}\else\textsuperscript{\ensuremath{h}}\fi}
\newunicodechar{ᵠ}{\ifmmode^{q}\else\textsuperscript{\ensuremath{q}}\fi}
\newunicodechar{ₙ}{\ifmmode_{n}\else\textsubscript{\ensuremath{n}}\fi}
\newunicodechar{ₐ}{\ifmmode_{a}\else\textsubscript{\ensuremath{a}}\fi}
\newunicodechar{ᵢ}{\ifmmode_{i}\else\textsubscript{\ensuremath{i}}\fi}
\newunicodechar{ᵣ}{\ifmmode_{r}\else\textsubscript{\ensuremath{r}}\fi}
\newunicodechar{ₖ}{\ifmmode_{k}\else\textsubscript{\ensuremath{k}}\fi}
\newunicodechar{ᵦ}{\ifmmode_{\beta}\else\textsubscript{\ensuremath{\beta}}\fi}
\newunicodechar{ₓ}{\ifmmode_{x}\else\textsubscript{\ensuremath{x}}\fi}
\newfontfamily\popdisplay{ArchivoBlack-Regular.ttf}[Path=../../../fonts/]
\newfontfamily\poplogo{SpaceGrotesk-Bold.ttf}[Path=../../../fonts/]
\newfontfamily\popsemi{PlusJakartaSans-SemiBold.ttf}[Path=../../../fonts/]
\newfontfamily\popxbold{PlusJakartaSans-ExtraBold.ttf}[Path=../../../fonts/]
\newfontfamily\popxboldls{PlusJakartaSans-ExtraBold.ttf}[Path=../../../fonts/,LetterSpace=3.0]

\setlist[enumerate]{leftmargin=1.7em,itemsep=5pt,topsep=4pt}
\setlist[itemize]{leftmargin=1.7em,itemsep=4pt,topsep=4pt,label={\color{poporange}\textbullet}}
\setlength{\parindent}{0pt}
\setlength{\parskip}{5pt}
\renewcommand{\arraystretch}{1.25}

% pgfplots : style Pop par défaut
\pgfplotsset{
  every axis/.append style={axis line style={popink,line width=1pt},tick style={popink},
    grid style={popline,line width=0.5pt},label style={font=\small},tick label style={font=\footnotesize}},
  popcurve/.style={poporange,line width=1.8pt,no markers,smooth},
}
% Même style disponible dans un \draw TikZ simple (hors pgfplots)
\tikzset{popcurve/.style={poporange,line width=1.8pt}}
\tikzset{popgrid/.style={popline,very thin}}
\colorlet{popcurve}{poporange} % le nom du style est parfois utilisé comme couleur

% ============================================================
% Pill / squiggle
% ============================================================
\newcommand{\poppill}[3]{%
  \tikz[baseline=-0.55ex]\node[draw=popink,line width=1.2pt,rounded corners=2.6pt,%
    fill=#2,inner xsep=6.5pt,inner ysep=3pt]{%
    \popxboldls\fontsize{7}{7}\selectfont\color{#3}\MakeUppercase{#1}};%
}
\newcommand{\popsquiggle}[1]{%
  \tikz[baseline=-0.4ex]{
    \draw[#1,line width=2.6pt,line cap=round]
      plot [smooth] coordinates {(0,0.09) (0.34,0.01) (0.68,0.21) (1.02,0.01) (1.36,0.10)};
  }%
}
% Mot-clé surligné (marqueur orange sous le texte)
\newcommand{\cle}[1]{{\popxbold\color{popdark}#1}}

% ============================================================
% En-tête / pied
% ============================================================
\pagestyle{fancy}
\fancyhf{}
\renewcommand{\headrulewidth}{0pt}
\renewcommand{\footrulewidth}{0pt}
\lhead{\raisebox{-0.15ex}{\poplogo\fontsize{13}{13}\selectfont\color{popink}OnBuch\color{poporange}+}}
\rhead{\poppill{\DOCMATIERE\ \textbullet\ \DOCNIVEAU\ \textbullet\ Leçon \DOCLECON}{white}{popink}}
\lfoot{\popxbold\fontsize{7.6}{7.6}\selectfont\color{popmuted}\MakeUppercase{Cours signé OnBuch+}\ \textbullet\ {\popsemi\DOCTITRE}}
\rfoot{\tikz[baseline=-0.6ex]\node[rounded corners=8.5pt,fill=popink,minimum height=17pt,minimum width=17pt,inner xsep=5pt,inner sep=0]{\popxbold\fontsize{8}{8}\selectfont\color{popcream}\thepage\,/\,\pageref*{LastPage}};}

% ============================================================
% Titres
% ============================================================
\newcommand{\chaptitre}[2]{%
  \begin{tcolorbox}[enhanced,colback=poporange,colframe=popink,boxrule=2.6pt,arc=11pt,
      drop shadow=popink,left=15pt,right=15pt,top=12pt,bottom=11pt,before skip=3pt,after skip=18pt]
    \poppill{#2}{popink}{popcream}\par\vspace{9pt}
    {\raggedright\hyphenpenalty=10000\exhyphenpenalty=10000\popdisplay\fontsize{22}{24}\selectfont\color{popcream}\MakeUppercase{#1}\par}
  \end{tcolorbox}
}
% Section numérotée : pastille orange avec le numéro + titre Archivo
\newcounter{popsec}
\newcommand{\coursec}[1]{%
  \stepcounter{popsec}\setcounter{popsub}{0}%
  \par\vspace{16pt}\noindent
  \begin{minipage}{\linewidth}
  \tikz[baseline=-0.7ex]\node[draw=popink,line width=1.6pt,fill=poporange,rounded corners=4pt,minimum size=22pt,inner sep=0]{\popdisplay\fontsize{12}{12}\selectfont\color{popcream}\thepopsec};%
  \hspace{8pt}{\popdisplay\fontsize{15}{17}\selectfont\color{popink}\MakeUppercase{#1}}\par
  \vspace{4pt}\noindent{\color{poporange}\rule{36pt}{3.6pt}}
  \end{minipage}\par\nopagebreak\vspace{8pt}
}
\newcounter{popsub}
\newcommand{\courssub}[1]{%
  \stepcounter{popsub}%
  \par\vspace{9pt}\noindent{\popxbold\fontsize{11.5}{13}\selectfont\color{popink}{\color{poporange}\thepopsec.\thepopsub}\hspace{6pt}#1}\par\nopagebreak\vspace{4pt}
}

% ============================================================
% Boîtes pédagogiques — même recette : fond blanc, bordure épaisse colorée,
% ombre dure encre, badge plein. Titre optionnel [#1].
% ============================================================
\tcbset{popbase/.style={breakable,enhanced,colback=white,boxrule=2.3pt,arc=9pt,
  shadow={3pt}{-3pt}{0pt}{popink},left=14pt,right=14pt,top=11pt,bottom=11pt,before skip=11pt,after skip=15pt,
  fontupper=\popsemi\fontsize{10.6}{14.5}\selectfont\color{popink},
  fontlower=\popsemi\fontsize{10.6}{14.5}\selectfont\color{popink}}}
% Coche dessinée (le glyphe ✓ n'existe pas dans les polices du document)
\newcommand{\popcheck}[1]{\tikz[baseline=-0.5ex]\draw[#1,line width=1.6pt,line cap=round,line join=round](-0.13,0.0)--(-0.04,-0.09)--(0.13,0.1);}
\providecommand{\coche}{\popcheck{popgreen}}
\providecommand{\resultat}[1]{\par\smallskip\noindent\fcolorbox{popgreen}{popgreenL}{\parbox{\dimexpr\linewidth-2\fboxsep-2\fboxrule\relax}{\textbf{Résultat.} #1}}\par\smallskip}
\newcommand{\popboxhead}[3]{\poppill{#1}{#2}{white}\ifx\relax#3\relax\else\hspace{6pt}{\popxbold\fontsize{10.6}{12}\selectfont\color{popink}#3}\fi\par\vspace{7pt}}

\newtcolorbox{aretenir}[1][]{popbase,colframe=popgreen,before upper={\popboxhead{À retenir}{popgreen}{#1}}}
\newtcolorbox{exemplebox}[1][]{popbase,colframe=poporange,before upper={\popboxhead{Exemple}{poporange}{#1}}}
\newtcolorbox{definition}[1][]{popbase,colframe=popblue,before upper={\popboxhead{Définition}{popblue}{#1}}}
\newtcolorbox{propriete}[1][]{popbase,colframe=poppurple,before upper={\popboxhead{Propriété}{poppurple}{#1}}}
\newtcolorbox{methode}[1][]{popbase,colframe=popgold,before upper={\popboxhead{Méthode}{popgold}{#1}}}
\newtcolorbox{attention}[1][]{popbase,colframe=poppink,before upper={\popboxhead{Attention}{poppink}{#1}}}
\newtcolorbox{experience}[1][]{popbase,colframe=popblue,colback=popblueL,before upper={\popboxhead{Expérience}{popblue}{#1}}}
% « Le savais-tu ? » : carte violette pleine (contexte, culture, Cameroun)
\newtcolorbox{savaistu}[1][]{popbase,colback=poppurple,colframe=popink,
  fontupper=\popsemi\fontsize{10.4}{14.2}\selectfont\color{white},
  before upper={\poppill{Le savais-tu\,?}{white}{poppurple}\ifx\relax#1\relax\else\hspace{6pt}{\popxbold\color{white}#1}\fi\par\vspace{7pt}}}
% Exercice résolu : énoncé en haut, \tcblower puis la solution sur fond crème
\newcounter{popexo}\newcounter{popexr}
\newtcolorbox{exoresolu}[1][]{popbase,colframe=poporange,colbacklower=popcream,
  segmentation style={popink,line width=1.2pt,dashed},
  before upper={\stepcounter{popexr}\popboxhead{Exercice résolu \thepopexr}{poporange}{#1}},
  before lower={\poppill{Solution}{popink}{popcream}\par\vspace{6pt}}}
% Exercice (non résolu en place) — la correction est dans la partie Corrigés
\newtcolorbox{exercice}[1][]{popbase,colframe=popink,
  before upper={\stepcounter{popexo}\popboxhead{Exercice \thepopexo}{popink}{#1}}}
% Corrigé
\newtcolorbox{corrige}[1][]{popbase,colframe=popgreen,colback=popgreenL,
  before upper={\popboxhead{Corrigé}{popgreen}{#1}}}
% Objectifs / prérequis / bilan
\newtcolorbox{objectifs}{popbase,colframe=popink,colback=poporangeL,
  before upper={\popboxhead{Objectifs de la leçon}{popink}{}}}
\newtcolorbox{prerequis}{popbase,colframe=popink,colback=popblueL,
  before upper={\popboxhead{Prérequis}{popblue}{}}}
\newtcolorbox{bilan}{popbase,colframe=popink,colback=white,boxrule=2.8pt,
  before upper={\popboxhead{Fiche bilan}{poporange}{L'essentiel à connaître pour l'examen}}}
% Figure encadrée avec légende
\newtcolorbox{popfigure}[1][]{enhanced,breakable=false,colback=white,colframe=popline,boxrule=1.4pt,arc=8pt,
  left=10pt,right=10pt,top=10pt,bottom=8pt,before skip=10pt,after skip=12pt,halign=center,
  lower separated=false,halign lower=center,
  fontlower=\popsemi\fontsize{9}{11.5}\selectfont\color{popmuted},
  #1}
\newcommand{\legende}[1]{\par\vspace{6pt}{\popsemi\fontsize{9}{11.5}\selectfont\color{popmuted}#1}}

% ============================================================
% Signature OnBuch+
% ============================================================
\newcommand{\poplogoinline}[1]{{\poplogo\fontsize{#1}{#1}\selectfont\color{popink}OnBuch\color{poporange}+}}
\newcommand{\signatureonbuch}{%
  \begin{tcolorbox}[enhanced,colback=popink,colframe=popink,boxrule=2.2pt,arc=10pt,
      drop shadow=poporange,left=14pt,right=14pt,top=10pt,bottom=10pt,before skip=0pt,after skip=0pt]
    \tikz[baseline=-0.7ex]\node[circle,fill=popgreen,draw=popcream,line width=1.3pt,minimum size=22pt,inner sep=0]{\popcheck{white}};%
    \hspace{9pt}\begin{minipage}[c]{0.62\linewidth}
      {\popxbold\fontsize{12}{13}\selectfont\color{popcream}Signé\ \textbullet\ {\poplogo OnBuch\color{poporange}+}}\par\vspace{2pt}
      {\raggedright\popsemi\fontsize{8}{10.5}\selectfont\color{popline}\MakeUppercase{Équipe pédagogique OnBuch+}\\\MakeUppercase{Conforme au programme officiel MINESEC}\par}
    \end{minipage}\hfill
    {\poplogo\fontsize{22}{22}\selectfont\color{popcream}OnBuch\color{poporange}+}
  \end{tcolorbox}%
}

% ============================================================
% Page de garde
% #1 titre · #2 matière · #3 niveau · #4 module · #5 n° leçon · #6 durée ·
% #7 description / objectifs (texte ou itemize)
% ============================================================
\newcommand{\pagegarde}[7]{%
  \thispagestyle{empty}
  \newgeometry{a4paper,margin=1.7cm,top=1.6cm,bottom=1.4cm}%
  \begin{tikzpicture}[remember picture,overlay]
    \fill[poporange!18] ([xshift=-3.2cm,yshift=-3.6cm]current page.north east) circle (4.6cm);
    \fill[poppurple!14] ([xshift=2.4cm,yshift=7.6cm]current page.south west) circle (3.8cm);
  \end{tikzpicture}%
  {\poplogo\fontsize{30}{30}\selectfont\color{popink}OnBuch\color{poporange}+}\hfill
  \raisebox{0.6ex}{\poppill{Cours signé \textbullet\ Premium}{popink}{popcream}}\par
  \vspace{1pt}\popsquiggle{poppurple}\par
  \vspace{8pt}
  \begin{tcolorbox}[enhanced,colback=poporange,colframe=popink,boxrule=2.8pt,arc=13pt,
      drop shadow=popink,left=18pt,right=18pt,top=12pt,bottom=13pt,after skip=10pt]
    \poppill{#2 \textbullet\ #3}{popink}{popcream}\hspace{5pt}\poppill{#4}{white}{popink}\par\vspace{8pt}
    {\popdisplay\fontsize{14}{14}\selectfont\color{popink}LEÇON #5}\par\vspace{4pt}
    {\raggedright\hyphenpenalty=10000\exhyphenpenalty=10000\popdisplay\fontsize{25}{27}\selectfont\color{popcream}\MakeUppercase{#1}\par}
    \vspace{8pt}
    {\popxbold\fontsize{9.5}{9.5}\selectfont\color{popink}\MakeUppercase{Durée conseillée : #6}}
  \end{tcolorbox}
  \begin{tcolorbox}[enhanced,colback=white,colframe=popink,boxrule=2.2pt,arc=10pt,
      drop shadow=popink,left=16pt,right=16pt,top=10pt,bottom=10pt,after skip=8pt]
    \poppill{Dans cette leçon}{poppurple}{white}\par\vspace{5pt}
    \popsemi\fontsize{9.3}{12.6}\selectfont\color{popink}#7
  \end{tcolorbox}
  \noindent\begin{minipage}[t]{0.487\linewidth}
    \begin{tcolorbox}[enhanced,colback=popgreen,colframe=popink,boxrule=2.2pt,arc=10pt,
        drop shadow=popink,left=11pt,right=11pt,top=10pt,bottom=10pt,height=3.1cm,valign=center]
      \tcbox[colback=white,colframe=popink,boxrule=1.3pt,arc=5pt,boxsep=0pt,nobeforeafter,
        left=3pt,right=3pt,top=3pt,bottom=3pt,valign=center]{\qrcode[height=1.9cm]{https://play.google.com/store/apps/details?id=com.luvvix.onbuch&hl=fr}}%
      \hspace{8pt}\begin{minipage}[c]{0.5\linewidth}
        {\popdisplay\fontsize{11}{12.5}\selectfont\color{white}\MakeUppercase{Télécharge OnBuch}}\par\vspace{4pt}
        {\raggedright\popsemi\fontsize{8.4}{11}\selectfont\color{white}Gratuit \textbullet\ Google Play\\Tuteur IA, annales, concours, résultats\par}
      \end{minipage}
    \end{tcolorbox}
  \end{minipage}\hfill
  \begin{minipage}[t]{0.487\linewidth}
    \begin{tcolorbox}[enhanced,colback=poppurple,colframe=popink,boxrule=2.2pt,arc=10pt,
        drop shadow=popink,left=11pt,right=11pt,top=10pt,bottom=10pt,height=3.1cm,valign=center]
      \tikz[baseline=-0.7ex]\node[circle,fill=white,draw=popink,line width=1.3pt,minimum size=1.6cm,inner sep=0]{\popdisplay\fontsize{24}{24}\selectfont\color{poppurple}+};%
      \hspace{8pt}\begin{minipage}[c]{0.6\linewidth}
        {\popdisplay\fontsize{11}{12.5}\selectfont\color{white}\MakeUppercase{Passe à OnBuch+}}\par\vspace{4pt}
        {\raggedright\popsemi\fontsize{8.4}{11}\selectfont\color{white}Tous les cours signés, Tuteur IA illimité, fiches PDF — onglet OnBuch+ de l'app\par}
      \end{minipage}
    \end{tcolorbox}
  \end{minipage}
  \vspace{8pt}
  \popcontenu
  \vfill
  \signatureonbuch
  \restoregeometry
  \newpage
}

% Rangée « Ce cours contient » (page de garde)
\newcommand{\poptile}[3]{% #1 couleur #2 gros texte #3 légende
  \begin{tcolorbox}[enhanced,colback=#1,colframe=popink,boxrule=2pt,arc=9pt,shadow={2.5pt}{-2.5pt}{0pt}{popink},
      left=6pt,right=6pt,top=9pt,bottom=9pt,halign=center,valign=center,height=1.75cm,width=\linewidth,nobeforeafter]
    {\popdisplay\fontsize{12.5}{13}\selectfont\color{white}\MakeUppercase{#2}}\par\vspace{3pt}
    {\centering\popsemi\fontsize{7.8}{9.5}\selectfont\color{white}#3\par}
  \end{tcolorbox}}
\newcommand{\popcontenu}{%
  \par\noindent{\popxboldls\fontsize{7.5}{7.5}\selectfont\color{popmuted}CE COURS SIGNÉ CONTIENT}\par\vspace{7pt}
  \noindent\begin{minipage}[t]{0.235\linewidth}\poptile{popblue}{Cours}{complet, illustré, pas à pas}\end{minipage}\hfill
  \begin{minipage}[t]{0.235\linewidth}\poptile{poporange}{Méthodes}{et exercices résolus}\end{minipage}\hfill
  \begin{minipage}[t]{0.235\linewidth}\poptile{poppink}{Exercices}{type examen + corrigés}\end{minipage}\hfill
  \begin{minipage}[t]{0.235\linewidth}\poptile{popgreen}{Fiche bilan}{l'essentiel à retenir}\end{minipage}\par}

% Sommaire
\newtcolorbox{sommaire}{enhanced,colback=white,colframe=popink,boxrule=2.2pt,arc=10pt,
  drop shadow=popink,left=18pt,right=18pt,top=13pt,bottom=13pt,before skip=6pt,after skip=20pt,
  before upper={\poppill{Sommaire}{poppurple}{white}\par\vspace{9pt}\popsemi\fontsize{10.6}{14.5}\selectfont\color{popink}}}

% Signature de fin de cours — poussée tout en bas de la dernière page
\newcommand{\findecours}{%
  \par\vfill\nopagebreak
  \begin{center}\popsquiggle{poporange}\end{center}\vspace{4pt}
  \signatureonbuch
}
\AtBeginDocument{\def\llbracket{\mathopen{[\mkern-3.5mu[}}\def\rrbracket{\mathclose{]\mkern-3.5mu]}}} % intervalles d'entiers (glyphe ⟦ absent de la police)
% Glyphes absents de Fira Math : redéfinis à partir de glyphes disponibles
\AtBeginDocument{%
  \def\setminus{\mathbin{\backslash}}\let\smallsetminus\setminus
  \def\vdots{\mathord{\vbox{\baselineskip3.2pt\lineskiplimit0pt\kern4pt\hbox{.}\hbox{.}\hbox{.}}}}%
  \def\ddots{\mathinner{\mkern1mu\raise6.5pt\hbox{.}\mkern2mu\raise3.6pt\hbox{.}\mkern2mu\raise0.7pt\hbox{.}\mkern1mu}}%
  \def\triangle{\mathord{\scalebox{0.9}{$\symup{\Delta}$}}}%
  \def\nabla{\mathord{\raisebox{\depth}{\scalebox{1}[-1]{$\symup{\Delta}$}}}}%
  \def\checkmark{\popcheck{popdark}}%
  \def\star{\mathbin{\ast}}\def\bigstar{\mathord{\ast}}%
  \def\diamond{\mathbin{\scalebox{0.6}{$\Diamond$}}}%
  \def\bigcup{\mathop{\mathchoice{\raisebox{-0.3em}{\scalebox{1.7}{$\cup$}}}{\scalebox{1.25}{$\cup$}}{\cup}{\cup}}\displaylimits}%
  \def\bigcap{\mathop{\mathchoice{\raisebox{-0.3em}{\scalebox{1.7}{$\cap$}}}{\scalebox{1.25}{$\cap$}}{\cap}{\cap}}\displaylimits}%
  \def\lesseqqgtr{\mathrel{\lessgtr}}\def\bigoplus{\mathop{\oplus}\displaylimits}%
}
\providecommand{\ssi}{si et seulement si} % abréviation fréquente
\AtBeginDocument{\providecommand{\wideparen}{\overparen}} % arc de cercle

\begin{document}

\pagegarde{Dissertation : rédiger une dissertation}{Français}{1re Tech}{Français – classe de Première technique}{N03}{19 séances (fiche de progression harmonisée)}{Cette leçon outille l'élève de Première technique pour analyser un sujet, construire un plan détaillé et rédiger une dissertation complète (introduction, développement, conclusion). Elle s'appuie sur la lecture méthodique du Lion et la perle de Wole Soyinka et mobilise les outils de la langue (types de phrases, synonymie, texte argumentatif) pour produire un texte cohérent, clair et convaincant.
\vspace{4pt}
\begin{multicols}{2}\small
\begin{itemize}
\item À la fin de cette leçon, je sais identifier les caractéristiques d'un texte argumentatif (thèse, arguments, exemples, connecteurs).
\item À la fin de cette leçon, je sais analyser un sujet de dissertation : dégager le thème, la problématique et les présupposés.
\item À la fin de cette leçon, je sais construire un plan détaillé cohérent (dialectique, thématique ou analytique).
\item À la fin de cette leçon, je sais rédiger une introduction complète (amorce, présentation du sujet, problématique, annonce du plan).
\item À la fin de cette leçon, je sais rédiger un paragraphe de développement (idée directrice, argument, exemple, lien logique).
\item À la fin de cette leçon, je sais rédiger une conclusion (bilan, réponse à la problématique, ouverture).
\end{itemize}\end{multicols}}

\begin{sommaire}
\begin{multicols}{2}
\begin{enumerate}
\item Le texte argumentatif et la dissertation : repères
\item Analyse du sujet et construction du plan détaillé
\item Rédiger l'introduction
\item Rédiger le développement : le paragraphe argumentatif
\item Rédiger la conclusion
\item Confrontation, évaluation et amélioration des productions
\item Activité d'intégration
\item Méthodes et exercices résolus
\item Exercices
\item Corrigés des exercices
\item Fiche bilan
\end{enumerate}
\end{multicols}
\end{sommaire}

\begin{objectifs}
\begin{itemize}
\item À la fin de cette leçon, je sais identifier les caractéristiques d'un texte argumentatif (thèse, arguments, exemples, connecteurs).
\item À la fin de cette leçon, je sais analyser un sujet de dissertation : dégager le thème, la problématique et les présupposés.
\item À la fin de cette leçon, je sais construire un plan détaillé cohérent (dialectique, thématique ou analytique).
\item À la fin de cette leçon, je sais rédiger une introduction complète (amorce, présentation du sujet, problématique, annonce du plan).
\item À la fin de cette leçon, je sais rédiger un paragraphe de développement (idée directrice, argument, exemple, lien logique).
\item À la fin de cette leçon, je sais rédiger une conclusion (bilan, réponse à la problématique, ouverture).
\item À la fin de cette leçon, je sais employer correctement les types de phrases obligatoires et facultatifs pour nuancer et rythmer mon argumentation.
\item À la fin de cette leçon, je sais utiliser la synonymie et l'antonymie pour éviter les répétitions et opposer les idées.
\item À la fin de cette leçon, je sais confronter ma production à celle d'autrui, l'évaluer selon une grille et l'améliorer.
\end{itemize}
\end{objectifs}

\begin{prerequis}
\begin{itemize}
\item La lecture méthodique d'une œuvre complète (notions de Seconde : incipit, personnages, thèmes).
\item Les classes grammaticales et la phrase simple (notions de 4e/3e).
\item Le paragraphe argumentatif simple rédigé en Seconde.
\item Les connecteurs logiques de base (d'abord, ensuite, cependant, donc).
\item La présentation générale du Lion et la perle (auteur, contexte nigérian, résumé).
\end{itemize}
\end{prerequis}

\newpage

\coursec{Le texte argumentatif et la dissertation : repères}

\courssub{Situation d'accroche et problématique}

Imaginez la scène : nous sommes au lycée technique de Bafoussam, en pleine finale du concours d'éloquence. Un élève, solide dans ses convictions, lance : \og La formation technique est plus utile au Cameroun que la formation générale \fg{}. Son adversaire a des idées justes, il connaît les chiffres de l'insertion professionnelle, il sent la force de l'argument économique. Mais il s'embrouille : pas d'introduction qui capte l'attention, des arguments qui se chevauchent sans connecteurs logiques, pas de conclusion qui cloue le bec. Le jury, des enseignants expérimentés, sanctionne : zéro point à l'argumentation. L'idée était là, la \cle{structure} manquait.

\medskip
\noindent\textbf{Problématique :} Comment transformer une opinion personnelle en un raisonnement rigoureux, structuré et persuasif ? Comment maîtriser les codes de la dissertation pour convaincre un jury d'examinateurs au Baccalauréat technique ?

\medskip
Cette première section pose les fondations : comprendre ce qu'est un texte argumentatif, identifier les attendus précis de la dissertation technique, et découvrir comment la littérature (avec \emph{Le Lion et la perle}) nourrit notre culture argumentative.

\courssub{Le texte argumentatif : visée, structure et marques}

\begin{definition}[Le texte argumentatif]
Un \textbf{texte argumentatif} est un discours (oral ou écrit) dont la \textbf{visée} principale est d'\textbf{influencer} un destinataire (lecteur ou auditeur) pour qu'il adhère à une \textbf{thèse} (un point de vue, une opinion défendue). Pour y parvenir, l'auteur construit un \textbf{raisonnement} articulé autour d'\textbf{arguments} (idées forces) illustrés par des \textbf{exemples} (faits précis, citations, statistiques, anecdotes).
\end{definition}

\medskip
\noindent Contrairement au texte explicatif qui \textbf{informe} (répondre à \og comment ? \fg{} ou \og pourquoi ? \fg{} de manière neutre), le texte argumentatif \textbf{prend parti}. Il ne dit pas \og La formation technique existe \fg{}, il dit \og La formation technique \textbf{doit être privilégiée} car... \fg{}.

\begin{propriete}[Les trois piliers de l'argumentation]
Toute argumentation solide repose sur un triptyque indissociable :
\begin{enumerate}
    \item \textbf{La Thèse :} La réponse affirmative ou négative à la question posée. C'est le \og fil conducteur \fg{}.
    \item \textbf{Les Arguments :} Les raisons logiques qui soutiennent la thèse. Ils répondent à \og Pourquoi ? \fg{}.
    \item \textbf{Les Exemples (ou preuves) :} Les illustrations concrètes qui ancrent l'argument dans le réel. Ils répondent à \og Preuve ? \fg{}.
\end{enumerate}
\end{propriete}

\medskip
\noindent\underline{Distinction cruciale : Convaincre vs Persuader.}
\begin{itemize}
    \item \textbf{Convaincre} (visée \cle{logique} / \cle{raison} / \cle{logos}) : S'adresse à l'intelligence. Utilise des faits, des statistiques, des raisonnements déductifs/inductifs. \og Le taux d'insertion des techniciens est de 72\% selon l'INS 2023. \fg{}
    \item \textbf{Persuader} (visée \cle{émotionnelle} / \cle{pathos}) : S'adresse au cœur, aux valeurs, à la sensibilité. Utilise l'ironie, l'éloquence, les figures de style, l'appel aux valeurs patriotiques. \og Nos villages meurent parce que nos jeunes fuient vers des diplômes qui ne nourrissent pas leur famille. \fg{}
\end{itemize}

\begin{attention}[Piège majeur : Confondre persuader et convaincre]
Au Baccalauréat technique, la dissertation est un exercice de \textbf{conviction rationnelle}. Un texte qui ne fait qu'émouvoir (pathos) sans démontrer (logos) sera sanctionné. \textbf{La dissertation convainc d'abord ; elle peut persuader en appui, jamais à la place de la raison.}
\end{attention}

\medskip
\noindent Comment repérer un texte argumentatif ? Par ses \textbf{marques linguistiques} (programme : Stylistique/Rhétorique et Morphosyntaxe -- Types de phrase) :

\begin{center}
\begin{tabularx}{\linewidth}{|l|X|X|}\hline
\rowcolor{popblueL}
\textbf{Marqueur} & \textbf{Fonction} & \textbf{Exemples} \\ \hline
\textbf{Connecteurs logiques} & Organisent le raisonnement (enchaînement, opposition, conséquence, concession) & \emph{En effet, par conséquent, toutefois, néanmoins, de surcroît, en outre, dès lors, c'est pourquoi} \\ \hline
\textbf{Modalisateurs} & Marquent l'engagement du locuteur (certitude, doute, volonté, obligation) & \emph{Il est certain que, sans doute, il faut, il est indispensable que, assurément, probablement} \\ \hline
\textbf{Champs lexicaux de l'opinion} & Révèlent la subjectivité assumée (jugement de valeur) & \emph{Bénéfique, nuisible, indispensable, catastrophique, progrès, régression, avantage, inconvénient} \\ \hline
\textbf{Types de phrases (Obligatoires)} & \emph{Déclarative} (affirmation), \emph{Interrogative} (question rhétorique), \emph{Impérative} (injonction), \emph{Exclamative} (émotion forte) & \emph{Modernisons nos campagnes ! (Imp.) / Qui oserait nier l'évidence ? (Int. rhét.) / Quel gâchis ! (Excl.)} \\ \hline
\end{tabularx}
\end{center}

\begin{savaistu}[Lexique et argumentation : Synonymie / Antonymie]
Le programme (Sémantique/Lexicologie) insiste sur la \textbf{synonymie} (nuances) et l'\textbf{antonymie} (opposition). En dissertation :
\begin{itemize}
    \item \textbf{Synonymes :} Éviter les répétitions (\og atout / avantage / force / mérite \fg{}) et nuancer (\og problème / difficulté / obstacle / défi \fg{}).
    \item \textbf{Antonymes :} Structurer l'opposition (Thèse/Antithèse) : \og modernité / tradition \fg{}, \og progrès / régression \fg{}, \og atout / limite \fg{}.
\end{itemize}
Maîtriser ces relations lexicales enrichit le \textbf{critère Langue} (variété, précision) et le \textbf{critère Plan} (articulation logique).
\end{savaistu}

\begin{exemplebox}[Analyse rapide d'un extrait argumentatif]
\textbf{Extrait :} \og \textbf{Il est urgent} (modalisateur) \textbf{d'investir} (impératif) dans l'agro-industrie. \textbf{En effet} (connecteur explication), le Cameroun importe 600 milliards de FCFA de produits alimentaires par an (argument chiffré). \textbf{Néanmoins} (connecteur concession), nos terres sont fertiles. \textbf{Par conséquent} (connecteur conséquence), former des techniciens agricoles \textbf{est une nécessité vitale} (jugement de valeur). \fg{}
\par\medskip
\noindent\textbf{Analyse :} Thèse = Investir dans l'agro-industrie. Argument = Dépendance alimentaire coûteuse. Exemple = Chiffre 600 milliards. Contre-argument anticipé (concession) = Terres fertiles sous-exploitées. Conclusion partielle = Nécessité vitale.
\end{exemplebox}

\courssub{La dissertation : définition, enjeux et critères d'évaluation au Baccalauréat technique}

\begin{definition}[La dissertation]
La \textbf{dissertation} est un \textbf{exercice écrit d'argumentation} qui consiste à développer de manière \textbf{organisée}, \textbf{personnelle} et \textbf{rigoureuse} une \textbf{réflexion} à partir d'un \textbf{sujet donné}. Elle mobilise les connaissances du candidat, sa culture générale et sa maîtrise de la langue pour \textbf{convaincre} un lecteur fictif (le correcteur) de la validité d'une thèse.
\end{definition}

\medskip
\noindent Au Baccalauréat technique (Probatoire / Baccalauréat), la dissertation est l'une des épreuves reines de l'expression écrite (coefficient variable selon les séries, souvent 2 à 4). Elle sanctionne la capacité à \textbf{structurer sa pensée}.

\begin{savaistu}[Critères officiels d'évaluation au Bac Technique]
La grille de correction officielle (MINESEC) repose sur quatre piliers pondérés (indicatifs) :
\begin{enumerate}
    \item \textbf{Pertinence du contenu (environ 40\%) :} Adéquation au sujet, qualité de l'analyse, justesse des arguments, richesse des exemples (ancrage camerounais valorisé).
    \item \textbf{Cohérence du plan et de la structure (environ 30\%) :} Respect de la structure Introduction / Développement / Conclusion, logique interne du plan (dialectique, thématique, analytique), transitions fluides.
    \item \textbf{Correction de la langue (environ 20\%) :} Syntaxe, orthographe, grammaire, vocabulaire précis, variété des connecteurs, absence de fautes graves (accords, temps).
    \item \textbf{Présentation et soin (environ 10\%) :} Lisibilité, paragraphes distincts, propreté, gestion du temps (copie complète).
\end{enumerate}
\emph{Astuce :} Une copie au plan parfait mais truffée de fautes perd les points de langue. Une copie sans fautes mais hors sujet perd les points de contenu. L'équilibre est la clé.
\end{savaistu}

\medskip
\noindent La dissertation technique se distingue de la dissertation littéraire (Première générale) par son ancrage \textbf{professionnel et concret}. Les sujets touchent souvent à l'environnement technique, économique, social ou culturel du Cameroun (ex: \og L'entrepreneuriat jeune au Cameroun : chance ou illusion ? \fg{}, \og La digitalisation de l'administration : atouts et limites \fg{}).

\courssub{Les types de sujets et le choix du plan adapté}

Avant d'écrire, il faut \textbf{décortiquer le sujet}. Trois grandes familles de formulations existent :

\begin{center}
\begin{tabularx}{\linewidth}{|l|X|l|}\hline
\rowcolor{popblueL}
\textbf{Type de sujet} & \textbf{Formulation type} & \textbf{Plan souvent privilégié} \\ \hline
\textbf{Question fermée (Oui/Non)} & \og Faut-il... ? \fg{}, \og Est-il souhaitable que... ? \fg{}, \og La technique tue-t-elle l'artisanat ? \fg{} & \textbf{Dialectique} (Thèse / Antithèse / Synthèse) \\ \hline
\textbf{Citation à discuter} & \og "La tradition est le feu qu'on entretient, pas les cendres qu'on vénère." (Discuter) \fg{} & \textbf{Dialectique} ou \textbf{Thématique} (selon la complexité) \\ \hline
\textbf{Sujet de réflexion / Ouvert} & \og Les défis de l'industrialisation au Cameroun. \fg{}, \og L'impact des réseaux sociaux sur la jeunesse. \fg{} & \textbf{Thématique} ou \textbf{Analytique} (Constat/Causes/Conseqs/Solutions) \\ \hline
\end{tabularx}
\end{center}

\medskip
\noindent Le choix du plan n'est pas arbitraire : il \textbf{découle de la formulation du sujet}.

\begin{methode}[Reconnaître le type de plan adapté à la formulation du sujet]
\begin{enumerate}
    \item \textbf{Analyser les mots-clés :} \og Faut-il \fg{}, \og Pour ou contre \fg{} $\rightarrow$ Opposition $\rightarrow$ \textbf{Plan Dialectique}.
    \item \textbf{Identifier la dimension :} \og Causes \fg{}, \og Conséquences \fg{}, \og Solutions \fg{}, \og Aspects \fg{} $\rightarrow$ Énumération/Analyse $\rightarrow$ \textbf{Plan Thématique} ou \textbf{Analytique}.
    \item \textbf{Tester le plan :} Écrire les grandes parties (I, II, III) en une phrase. Si elles s'opposent naturellement (Pour / Contre / Dépassement) $\rightarrow$ Dialectique. Si elles explorent des facettes complémentaires (Économique / Social / Culturel) $\rightarrow$ Thématique. Si elles suivent une logique causale (Réalité / Origines / Issus) $\rightarrow$ Analytique.
\end{enumerate}
\end{methode}

\medskip
\noindent Voici les trois architectures majeures :

\begin{popfigure}
\centering
\begin{tikzpicture}[
    node distance=1.2cm and 1.5cm,
    box/.style={draw=popblue, thick, rounded corners, fill=popblueL, text width=4.5cm, align=center, minimum height=2.8cm},
    arrow/.style={-Stealth, thick, popdark},
    title/.style={font=\bfseries\large, text=popblue},
    subtitle/.style={font=\small, text=popmuted}
]
% Plan Dialectique
\node[box, align=center] (dial){
    \textbf{PLAN DIALECTIQUE}\\[0.2cm]
    \textit{(Sujet polémique : "Faut-il... ?")}\\[0.3cm]
    \textbf{I. Thèse} (Arguments POUR)\\
    \textbf{II. Antithèse} (Arguments CONTRE)\\
    \textbf{III. Synthèse} (Dépassement / Nuance)
};
\node[title, above=0.2cm of dial.north] {A. Dialectique};

% Plan Thématique
\node[box, right=of dial, align=center] (them){
    \textbf{PLAN THÉMATIQUE}\\[0.2cm]
    \textit{(Sujet large : "Les aspects de...")}\\[0.3cm]
    \textbf{I. Dimension 1} (ex: Économique)\\
    \textbf{II. Dimension 2} (ex: Social)\\
    \textbf{III. Dimension 3} (ex: Culturel/Écologique)
};
\node[title, above=0.2cm of them.north] {B. Thématique};

% Plan Analytique
\node[box, right=of them, align=center] (ana){
    \textbf{PLAN ANALYTIQUE}\\[0.2cm]
    \textit{(Sujet processus/problème : "Causes/Effets")}\\[0.3cm]
    \textbf{I. Constat / État des lieux}\\
    \textbf{II. Causes / Origines}\\
    \textbf{III. Conséquences / Solutions}
};
\node[title, above=0.2cm of ana.north] {C. Analytique};
\end{tikzpicture}
\legende{Les trois architectures types de la dissertation technique. Le choix dépend de la formulation du sujet.}
\end{popfigure}

\begin{exemplebox}[Application : Choix du plan pour un sujet camerounais]
\textbf{Sujet :} \og Faut-il moderniser les villages camerounais au prix de leurs traditions ? \fg{}
\par\medskip
\noindent\textbf{Analyse :} Formulation \og Faut-il \fg{} = Question fermée, opposition nette (Modernité vs Tradition).
\par\medskip
\noindent\textbf{Plan Dialectique imposé :}
\begin{itemize}
    \item \textbf{I. Thèse :} La modernisation est un impératif de développement (santé, éducation, revenus, désenclavement).
    \item \textbf{II. Antithèse :} Elle menace l'identité, la cohésion sociale, l'écologie, le patrimoine immatériel (rites, langue, chefferies).
    \item \textbf{III. Synthèse :} Une \textbf{modernité endogène} : adapter les techniques (solaire, eau, formation) aux valeurs traditionnelles (concertation, respect des aînés, gestion communautaire). Ni muséification, ni déracinement.
\end{itemize}
\end{exemplebox}

\begin{attention}[Pièges classiques dans le choix du plan]
\begin{itemize}
    \item \textbf{Forcer un plan dialectique} sur un sujet ouvert (ex: \og L'agriculture au Cameroun \fg{}). On obtient un \og Pour/Contre \fg{} artificiel et vide.
    \item \textbf{Mélanger les plans :} Faire I. Causes, II. Avantages, III. Inconvénients. Incohérence logique fatale.
    \item \textbf{Oublier la Synthèse} en dialectique : Laisser le lecteur sur une opposition stérile. La synthèse n'est pas une conclusion, c'est une \textbf{troisième partie argumentée} qui résout la tension.
\end{itemize}
\end{attention}

\courssub{Le Lion et la perle : l'affrontement tradition/modernité comme matériau argumentatif (Lecture méthodique n°1)}

La littérature n'est pas un décor : c'est un \textbf{réservoir d'arguments}, de personnages-types et de situations emblématiques. La Lecture méthodique n°1 de \emph{Le Lion et la perle} (\emph{The Lion and the Jewel}) de Wole Soyinka (programme officiel) nous offre une mise en scène théâtrale parfaite du conflit qui structure nombre de sujets de dissertation technique au Cameroun.

\begin{experience}[Le village d'Ilujinlé : laboratoire du débat Tradition vs Modernité]
\textbf{Contexte :} Pièce en trois parties (Matin, Midi, Soir). Trois personnages centraux incarnent trois positions argumentatives.
\begin{center}
\begin{tabularx}{\linewidth}{|l|X|X|X|}\hline
\rowcolor{popblueL}
\textbf{Personnage} & \textbf{Position (Thèse)} & \textbf{Arguments types} & \textbf{Registre / Style} \\ \hline
\textbf{Lakunle} (Instituteur) & \textbf{Modernité radicale / Mimétisme occidental} & Progrès technique, éducation formelle, monogamie, rejet de la dot (\og barbarie \fg{}), émancipation de la femme par l'école. & Langage ampoulé, citations mal digérées, ironie, \textbf{persuasion ratée} (ridicule). \\ \hline
\textbf{Baroka} (Bale / Chef) & \textbf{Tradition pragmatique / Modernité contrôlée} & Sagesse des aînés, pouvoir politique, ruse, adaptation (machine à timbres, projet usine), polygamie comme alliance sociale. & Proverbes, paraboles, humour, \textbf{conviction efficace} (réussite). \\ \hline
\textbf{Sidi} (La Belle) & \textbf{Le peuple / L'enjeu} & Beauté = capital, fierté, désir de reconnaissance, refus d'être \og achetée \fg{} (dot), séduction par la modernité (magazine) mais attachement aux rites. & Chant, danse, images, \textbf{volonté de choisir}. \\ \hline
\end{tabularx}
\end{center}
\end{experience}

\medskip
\noindent Pourquoi est-ce un \og matériau argumentatif \fg{} pour vous ?
\begin{itemize}
    \item \textbf{Lakunle} incarne l'argument \og moderne \fg{} mal maîtrisé : il \textbf{persuade} (ou croit le faire) par le mépris, sans comprendre le réel. \textbf{Contre-exemple} de dissertation : thèse sans ancrage réel, vocabulaire creux, connecteurs maladroits.
    \item \textbf{Baroka} incarne l'argument \og traditionnel \fg{} intelligent : il \textbf{convainc} par la ruse, la connaissance de l'adversaire, l'adaptation. \textbf{Modèle} : thèse claire, arguments concrets (machine, timbres), exemples vécus, concession stratégique (il veut le progrès \og à sa manière \fg{}).
    \item \textbf{Sidi} représente le \textbf{sujet même} : la jeunesse camerounaise tiraillée. Son choix final (Baroka) est une \textbf{synthèse} implicite : elle accepte la tradition (polygamie, dot) mais a failli basculer par vanité moderne (photos).
\end{itemize}

\begin{exoresolu}[Exploiter Le Lion et la perle dans une dissertation]
\textbf{Sujet d'entraînement :} \og La modernisation passe-t-elle nécessairement par le rejet des traditions ? \fg{}
\par\medskip
\noindent\textbf{Intégration littéraire (Partie I ou II ou Exemple) :}
\par\noindent\og Dans \emph{Le Lion et la perle}, Wole Soyinka met en scène l'échec de Lakunle, qui veut imposer la modernité par le déni de la culture yoruba (refus de la dot, mépris des rites). Son argumentaire, purement livresque et déconnecté du vécu villageois, ne \textbf{convainc} personne, pas même Sidi. À l'inverse, Baroka, figure de la tradition, intègre la modernité (machine à timbres, projet d'usine) pour mieux asseoir son pouvoir. La pièce démontre que \textbf{la modernité rejetée par orgueil (Lakunle) stérilise, tandis que la tradition qui s'adapte (Baroka) dure}. \fg{}
\par\medskip
\noindent\textbf{Analyse de l'exemple :} Thèse illustrée (Modernité brutale = échec), Argument (Déconnexion du réel), Exemple précis (Scène de la dot / Scène de la machine), Vocabulaire argumentatif (à l'inverse, démontre que, tandis que).
\end{exoresolu}

\courssub{Schéma global de la dissertation : l'architecture en trois temps}

Pour clore cette section de repères, visualisons la structure cible. Chaque bloc a ses \textbf{règles strictes} (détail dans les sections 3, 4, 5).

\begin{popfigure}
\centering
\begin{tikzpicture}[
    node distance=1.5cm and 1cm,
    mainblock/.style={draw=popdark, very thick, rounded corners=8pt, fill=popcream, text width=5.5cm, align=center, minimum height=4.5cm, inner sep=5pt},
    subblock/.style={draw=popblue, thick, rounded corners=4pt, fill=popblueL, text width=5cm, align=left, font=\small},
    arrow/.style={-Stealth, very thick, poporange},
    title/.style={font=\bfseries\large\color{popblue}, align=center},
    step/.style={font=\small, text=popdark}
]
% Introduction
\node[mainblock, align=center] (intro){
    \textbf{INTRODUCTION} \\[0.3cm]
    \begin{minipage}{5cm}
    \begin{enumerate}
        \item[\textbf{1.}] \textbf{Amorce / Accroche} (Citation, chiffre, actualité, fait de société)
        \item[\textbf{2.}] \textbf{Définition / Délimitation} des termes clés du sujet
        \item[\textbf{3.}] \textbf{Problématique} (Question unique, issue de la tension du sujet)
        \item[\textbf{4.}] \textbf{Annonce du plan} (I, II, III formulés en phrases nominales)
    \end{enumerate}
    \end{minipage}
};

% Développement
\node[mainblock, right=of intro, align=center] (dev){
    \textbf{DÉVELOPPEMENT} \\[0.3cm]
    \begin{minipage}{5cm}
    \textbf{Partie I :} Idée directrice 1\\
    \hspace{0.5cm} Arg 1 + Exemple\\
    \hspace{0.5cm} Arg 2 + Exemple\\
    \hspace{0.5cm} \textit{Mini-transition}\\[0.2cm]
    \textbf{Partie II :} Idée directrice 2\\
    \hspace{0.5cm} Arg 1 + Exemple\\
    \hspace{0.5cm} Arg 2 + Exemple\\
    \hspace{0.5cm} \textit{Mini-transition}\\[0.2cm]
    \textbf{Partie III :} Idée directrice 3 (Synthèse ou 3ème dimension)\\
    \hspace{0.5cm} Arg 1 + Exemple\\
    \hspace{0.5cm} Arg 2 + Exemple
    \end{minipage}
};

% Conclusion
\node[mainblock, right=of dev, align=center] (conc){
    \textbf{CONCLUSION} \\[0.3cm]
    \begin{minipage}{5cm}
    \begin{enumerate}
        \item[\textbf{1.}] \textbf{Bilan / Récapitulation} synthétique des 3 parties (réponse à la problématique)
        \item[\textbf{2.}] \textbf{Ouverture} (Perspective, autre sujet lié, citation finale, projection avenir)
    \end{enumerate}
    \end{minipage}
};

\draw[arrow] (intro.east) -- node[above, font=\small\bfseries\color{poporange}] {Annonce du plan} (dev.west);
\draw[arrow] (dev.east) -- node[above, font=\small\bfseries\color{poporange}] {Réponse apportée} (conc.west);

% Annotations verticales
\node[title, above=0.8cm of intro.north] {\textbf{1 Paragraphe unique}};
\node[title, above=0.8cm of dev.north] {\textbf{3 Parties = 3 Paragraphes majeurs}};
\node[title, above=0.8cm of conc.north] {\textbf{1 Paragraphe unique}};
\end{tikzpicture}
\legende{Architecture type de la dissertation au Baccalauréat technique. L'introduction et la conclusion font un paragraphe chacune. Le développement compte 3 parties (donc 3 grands paragraphes), chacune subdivisée en 2 arguments + exemples.}
\end{popfigure}

\begin{aretenir}[Repères essentiels de la section 1]
\begin{itemize}
    \item Le texte argumentatif vise à \textbf{convaincre} (raison) plus qu'à persuader (émotion). Structure : Thèse $\rightarrow$ Arguments $\rightarrow$ Exemples.
    \item La dissertation est un \textbf{exercice codifié} (Intro / Développement 3 parties / Conclusion) évalué sur 4 critères : Contenu, Plan, Langue, Présentation.
    \item \textbf{Le sujet dicte le plan :} Question fermée $\rightarrow$ Dialectique ; Sujet large $\rightarrow$ Thématique ; Problème processuel $\rightarrow$ Analytique.
    \item \emph{Le Lion et la perle} (Lecture n°1) fournit des \textbf{archétypes argumentatifs} : Lakunle (modernité stérile), Baroka (tradition adaptative), Sidi (enjeu). Utilisez-les comme exemples culturels forts.
    \item \textbf{Maîtriser les connecteurs logiques}, les \textbf{modalisateurs}, les \textbf{types de phrases} et le \textbf{lexique (synonymie/antonymie)} est la condition sine qua non de la cohérence (critères Plan + Langue).
\end{itemize}
\end{aretenir}

\medskip
\noindent Nous avons maintenant la carte et la boussole. La section suivante nous plonge dans l'atelier : \textbf{Analyser un sujet réel et construire son plan détaillé} (Section 2/6).

\coursec{Analyse du sujet et construction du plan détaillé}

Dans la section précédente, nous avons identifié les repères communs au texte argumentatif et à la dissertation : thèse, arguments, exemples, connecteurs logiques, structure en introduction-développement-conclusion. Nous savons désormais \emph{quoi} écrire. La question cruciale devient : \emph{comment} passer d'un sujet brut à un plan solide, cohérent et pertinent ? C'est tout l'objet de cette section. L'analyse du sujet n'est pas une formalité administrative : c'est le moment où se joue la réussite ou l'échec de la dissertation. Un sujet mal analysé conduit immanquablement au hors-sujet, à la répétition ou à l'incohérence. À l'inverse, une analyse rigoureuse produit un plan détaillé qui, comme nous le verrons, représente déjà plus de la moitié du travail de rédaction.

\courssub{Étape 1 : Lire le sujet plusieurs fois et souligner les mots-clés}

La première lecture doit être lente, active, crayon en main. On ne lit pas un sujet de dissertation comme on lit un roman. On le \emph{disseque}.

\begin{methode}[Lecture active du sujet]
\begin{enumerate}
\item \textbf{Première lecture :} compréhension globale. De quoi parle-t-on ? Quel domaine (social, culturel, politique, technique, philosophique) ?
\item \textbf{Deuxième lecture :} repérage des \cle{mots-clés} (noms, verbes, adjectifs porteurs de sens). Soulignez-les.
\item \textbf{Troisième lecture :} repérage des \cle{mots-outils} (conjonctions, prépositions, adverbes) qui indiquent la relation logique entre les termes : \og et \fg, \og ou \fg, \og mais \fg, \og parce que \fg, \og si \fg, \og est-ce que \fg, \og dans quelle mesure \fg.
\item \textbf{Quatrième lecture :} vérification. Chaque mot souligné a-t-il été interrogé ? Y a-t-il des termes ambigus, polisémiques ?
\end{enumerate}
\end{methode}

\begin{exemplebox}[Application sur un sujet type]
Sujet : \og \emph{La tradition est-elle un obstacle au développement ?} \fg

\begin{itemize}
\item \textbf{Mots-clés (soulignés) :} \textbf{tradition}, \textbf{obstacle}, \textbf{développement}.
\item \textbf{Mots-outils :} \og est-ce que \fg (question fermée, appelle oui/non mais exige nuance), \og un \fg (article indéfini $\rightarrow$ un obstacle parmi d'autres, pas \emph{le} seul).
\item \textbf{Termes à définir :} \emph{tradition} (coutumes, héritage, valeurs, rites ?), \emph{obstacle} (frein, blocage, ralentissement, remise en cause ?), \emph{développement} (économique, humain, social, durable, technique ?).
\end{itemize}
\end{exemplebox}

\begin{popfigure}
\begin{tikzpicture}[
  scale=0.9,
  every node/.style={font=\small},
  keyword/.style={draw=popblue, thick, rounded corners, fill=popblueL, inner sep=4pt},
  tool/.style={draw=poporange, thick, rounded corners, fill=poporangeL, inner sep=4pt},
  arrow/.style={->, thick, popdark}
]
\node[keyword] (trad) at (0,3) {tradition};
\node[keyword] (obs) at (4,3) {obstacle};
\node[keyword] (dev) at (8,3) {développement};
\node[tool] (estce) at (4,1) {est-ce que};
\node[tool] (un) at (6,1) {un (indéfini)};
\draw[arrow] (trad) -- node[above, sloped, pos=0.5] {relation ?} (obs);
\draw[arrow] (obs) -- node[above, sloped, pos=0.5] {vers} (dev);
\draw[arrow, dashed] (estce) -- (obs);
\draw[arrow, dashed] (un) -- (obs);
\node[align=center, text width=10cm, popmuted] at (4,-1.5) {
\textbf{Lecture 1 :} Sujet sur les rapports tradition / développement.\\
\textbf{Lecture 2 :} Trois concepts centraux à définir.\\
\textbf{Lecture 3 :} Question fermée + article indéfini $\rightarrow$ nuance obligatoire.\\
\textbf{Lecture 4 :} Polysémie de \emph{développement} (éco / humain / durable) à lever.
};
\end{tikzpicture}
\legende{Analyse visuelle du sujet : mots-clés (bleu) et mots-outils (orange). Les flèches indiquent les relations logiques à interroger.}
\end{popfigure}

\courssub{Étape 2 : Définir les termes du sujet et repérer les présupposés}

Chaque mot-clé porte une \cle{polysémie} (plusieurs sens). Le candidat doit choisir le ou les sens pertinents pour le sujet, les expliciter brièvement au brouillon, et les garder en tête pour la rédaction. Parallèlement, tout sujet contient des \cle{présupposés} : des affirmations implicites que le sujet tient pour acquises et qui orientent la réflexion.

\begin{definition}[Présupposé]
Un présupposé est une idée que le sujet suppose vraie sans la démontrer. L'identifier permet de savoir ce qui est \emph{donné} et ce qui reste à \emph{prouver} ou à \emph{discuter}.
\end{definition}

\begin{exemplebox}[Définitions et présupposés pour « La tradition est-elle un obstacle au développement ? »]
\begin{center}
\begin{tabularx}{\linewidth}{|l|X|X|}\hline
\textbf{Terme} & \textbf{Définition retenue (brouillon)} & \textbf{Présupposés détectés} \\ \hline
Tradition & Ensemble de coutumes, valeurs, savoir-faire transmis de génération en génération ; peut être vivant ou figé. & La tradition existe, elle est identifiable, elle a une consistance. \\ \hline
Obstacle & Ce qui freine, ralentit, empêche ou détourne une dynamique ; pas nécessairement un blocage total. & Il y a une dynamique (le développement) que la tradition pourrait entraver. \\ \hline
Développement & Processus d'amélioration durable des conditions de vie (économiques, sociales, culturelles, techniques) d'une société. & Le développement est désirable, mesurable, et il y a une direction (le « progrès »). \\ \hline
\end{tabularx}
\end{center}
\textbf{Présupposé majeur du sujet :} \og La tradition et le développement sont deux forces distinctes qui s'affrontent. \fg $\rightarrow$ À discuter : ne peuvent-elles pas se nourrir mutuellement ?
\end{exemplebox}

\begin{attention}[Piège : définir \emph{pendant} la rédaction]
Ne définissez \textbf{pas} les termes dans l'introduction de la dissertation finale (sauf si le sujet l'exige explicitement). Les définitions servent \emph{uniquement} au brouillon pour clarifier votre pensée. L'introduction montrera la maîtrise des concepts par l'usage précis que vous en ferez, pas par une liste de définitions scolaires.
\end{attention}

\courssub{Étape 3 : Formuler la problématique sous forme de question(s)}

La \cle{problématique} est la question centrale, précise et débattue, à laquelle la dissertation entière devra répondre. Elle transforme le sujet (souvent large, parfois caricatural) en une question de recherche intellectuelle. Une bonne problématique :
\begin{itemize}
\item reprend les termes du sujet (ou leurs synonymes précis) ;
\item manifeste une \textbf{tension}, une \textbf{complexité} (d'où le \og dans quelle mesure \fg, \og comment \fg, \og pourquoi \fg) ;
\item n'est \textbf{ni trop large} (risque de dispersion), \textbf{ni trop étroite} (risque de closure) ;
\item guide la construction du plan : chaque grande partie apporte un élément de réponse.
\end{itemize}

\begin{methode}[De la question du sujet à la problématique]
\begin{enumerate}
\item Repérer le type de sujet : \og Question fermée \fg (Oui/Non) $\rightarrow$ transformer en \og Dans quelle mesure... ? \fg / \og Comment... ? \fg ; \og Question ouverte \fg $\rightarrow$ recentrer ; \og Citation \fg $\rightarrow$ expliciter la thèse de l'auteur et la confronter.
\item Écrire 2 à 3 brouillons de problématique.
\item Choisir celle qui permet le meilleur plan dialectique (thèse / antithèse / synthèse ou causes / conséquences / perspectives).
\item La formuler \textbf{une seule fois}, clairement, à la fin de l'introduction.
\end{enumerate}
\end{methode}

\begin{exemplebox}[Problématiques possibles pour le sujet modèle]
\begin{itemize}
\item \textbf{Trop large :} \og Quels sont les rapports entre tradition et développement ? \fg (pas de tension, catalogue).
\item \textbf{Trop étroite :} \og La tradition empêche-t-elle l'industrialisation au Cameroun ? \fg (réduit le développement à l'industriel, la tradition au Cameroun).
\item \textbf{Équilibrée (retenue) :} \og \textbf{Dans quelle mesure la tradition, loin d'être un simple frein, peut-elle constituer un levier ou un cadre nécessaire pour un développement endogène et durable ?} \fg
\end{itemize}
Cette problématique :
\begin{itemize}
\item garde \og obstacle \fg (frein) mais ouvre sur \og levier \fg et \og cadre nécessaire \fg ;
\item précise \og développement \fg en \og endogène et durable \fg ;
\item appelle un plan en 3 parties : 1. Tradition comme frein (thèse) — 2. Tradition comme ressource (antithèse) — 3. Dépassement : articulation tradition / modernité pour un développement choisi (synthèse).
\end{itemize}
\end{exemplebox}

\courssub{Étape 4 : Mobiliser ses connaissances (brainstorming)}

Une fois la problématique fixée, on \emph{vidange} sa mémoire sur le brouillon. Aucune idée n'est triée à ce stade. On note tout : arguments théoriques, exemples littéraires (surtout \emph{Le Lion et la perle}), faits de société camerounais, actualité, connaissances techniques de sa spécialité.

\begin{experience}[Banque d'exemples depuis \emph{Le Lion et la perle} (Lecture méthodique n°2)]
Lors de la lecture méthodique n°2 de \emph{Le Lion et la perle} (Wole Soyinka), vous avez relevé les arguments des personnages. Voici comment les transformer en \cle{banque d'exemples} utilisable pour tout sujet touchant tradition / modernité / développement :
\begin{center}
\begin{tabularx}{\linewidth}{|l|X|X|}\hline
\textbf{Personnage} & \textbf{Position / Argument} & \textbf{Exemple exploitable en dissertation} \\ \hline
Baroka (le Bale) & Tradition comme pouvoir, ruse, adaptation ; il utilise la modernité (timbre-poste) pour asseoir son autorité traditionnelle. & La tradition n'est pas figée : elle \emph{absorbe} la modernité pour se perpétuer. Exemple d'articulation féconde. \\ \hline
Lakunle (instituteur) & Modernité comme rupture, rejet de la tradition (dot, polygamie), discours progressiste mais stérile, déconnecté du réel. & Modernité \emph{importée}, dogmatique $\rightarrow$ échec. Le développement sans ancrage culturel est illusoire. \\ \hline
Sidi (la perle) & Beauté, fierté, manipulation ; elle choisit in fine le tradition (Baroka) après avoir flerté avec la modernité (photo, magazine). & Le peuple / la jeunesse : attiré par les apparences de la modernité, mais ancré dans les valeurs traditionnelles. \\ \hline
Sadiku (épouse aînée) & Fidélité à la tradition, mais complice de la ruse de Baroka ; incarne la transmission orale, la rumeur. & La tradition comme réseau de communication, de régulation sociale. \\ \hline
\end{tabularx}
\end{center}
\textbf{Autres réservoirs :} vie camerounaise (chefferies, ngondo, tontines, agriculture familiale, startups tech à Yaoundé/Douala, diaspora), actualité (ODD, transition énergétique, IA, urbanisation), spécialité technique (normes, innovation, matériaux locaux, formation duale).
\end{experience}

\courssub{Étape 5 : Trier, regrouper et hiérarchiser les idées en parties et sous-parties}

C'est l'étape de l'\cle{architecture}. On passe du chaos créatif à la structure logique. Trois opérations successives :

\begin{enumerate}
\item \textbf{Trier :} éliminer le hors-sujet, les doublons, les exemples trop faibles, les idées qui ne répondent pas à la problématique.
\item \textbf{Regrouper :} rassembler les idées par \emph{affinité thématique} (ex. : tous les arguments « tradition = frein » ensemble ; tous les « tradition = ressource » ensemble).
\item \textbf{Hiérarchiser :} ordonner les groupes selon une \textbf{logique dialectique} (thèse / antithèse / synthèse) ou \textbf{analytique} (causes / mécanismes / conséquences / perspectives) ou \textbf{thématique} (économique / social / culturel / politique) — le choix dépend de la problématique.
\end{enumerate}

\begin{propriete}[Critères d'un bon regroupement]
\begin{itemize}
\item \textbf{Exhaustivité relative :} les 2 ou 3 parties couvrent les dimensions essentielles de la problématique.
\item \textbf{Mutualité exclusive :} pas de chevauchement entre parties (chaque argument a sa place unique).
\item \textbf{Progressivité :} l'ordre des parties fait avancer la pensée (du plus évident au plus complexe, du constat à la proposition).
\item \textbf{Équilibre :} parties de poids comparable (nombre d'arguments, densité d'exemples).
\end{itemize}
\end{propriete}

\courssub{Étape 6 : Rédiger le plan détaillé}

Le \cle{plan détaillé} est le document de travail final au brouillon. Il ne se contente pas de titres de parties : il explicite, pour chaque sous-partie, l'\textbf{idée directrice} (phrase complète), les \textbf{arguments} (raisonnements) et les \textbf{exemples} précis (avec références). C'est la \og partition \fg que le rédacteur suivra.

\begin{methode}[Structure type d'une entrée de plan détaillé]
Pour chaque sous-partie (ex. I.A, I.B, II.A, II.B, III.A, III.B) :
\begin{enumerate}
\item \textbf{Idée directrice} (1 phrase affirmative, répondant à un pan de la problématique).
\item \textbf{Argument(s)} (1 à 2 phrases : raisonnement, mécanisme, cause/effet).
\item \textbf{Exemple(s) précis} (référence : \emph{Le Lion et la perle}, Cameroun, actualité, spécialité technique).
\item \textbf{Transition vers la sous-partie suivante} (connecteur logique + annonce).
\end{enumerate}
\end{methode}

\begin{exoresolu}[Plan détaillé complet : « La tradition est-elle un obstacle au développement ? »]
\textbf{Problématique :} Dans quelle mesure la tradition, loin d'être un simple frein, peut-elle constituer un levier ou un cadre nécessaire pour un développement endogène et durable ?

\bigskip
\textbf{INTRODUCTION} (à rédiger in fine) : Accroche (proverbe / actualité) — Définitions implicites — Présupposés — Problématique — Annonce du plan.

\bigskip
\textbf{I. La tradition peut apparaître comme un obstacle au développement lorsqu'elle se fige en conservatisme ou qu'elle est instrumentalisée.}
\begin{itemize}
\item \textbf{A. Le poids des coutumes figées freine l'innovation technique et l'émancipation individuelle.}
      \begin{itemize}
      \item \textit{Idée directrice :} Certaines pratiques traditionnelles (dot excessive, interdits alimentaires, exclusion des femmes de la gestion foncière) limitent la mobilité sociale, l'investissement productif et l'égalité des chances.
      \item \textit{Argument :} La fixation rituelle transforme des règles de cohésion en barrières à l'entrée pour les jeunes et les femmes, ressources humaines majeures du développement.
      \item \textit{Exemples :} \emph{Le Lion et la perle} : Lakunle refuse de payer la dot $\rightarrow$ blocage du mariage ; Cameroun : accès à la terre pour les jeunes agriculteurs freiné par la gestion coutumière ; actualité : mariages précoces et scolarisation des filles.
      \item \textit{Transition :} Mais cette vision partielle occulte la dimension vivante, adaptative de la tradition.
      \end{itemize}
\item \textbf{B. L'instrumentalisation politique de la tradition entretient l'immobilisme au profit d'élites conservatrices.}
      \begin{itemize}
      \item \textit{Idée directrice :} Le pouvoir traditionnel ou étatique peut brandir la \og coutume \fg pour légitimer la prédation, la corruption ou le refus de réformes (décentralisation, parité, transparence).
      \item \textit{Argument :} La tradition devient idéologie de conservation des privilèges, non culture vivante.
      \item \textit{Exemples :} \emph{Le Lion et la perle} : Baroka utilise la tradition pour asseoir son pouvoir personnel (timbre, ruse) ; Cameroun : certains chefs traditionnels perçoivent des redevances foncières opaques ; actualité : débats sur la chefferie et la décentralisation.
      \item \textit{Transition :} Pourtant, réduite à ses excès, la tradition ignore sa fonction de régulation et d'innovation endogène.
      \end{itemize}
\end{itemize}

\bigskip
\textbf{II. La tradition vivante est un réservoir de valeurs, de savoir-faire et d'organisation sociale propices au développement endogène.}
\begin{itemize}
\item \textbf{A. Les valeurs traditionnelles (solidarité, responsabilité, respect de l'environnement) fondent un capital social indispensable.}
      \begin{itemize}
      \item \textit{Idée directrice :} L'entraide (tontines, travail communautaire), l'éthique de la parole donnée, la gestion commune des ressources (forêts sacrées, parcours) sont des actifs immatériels du développement.
      \item \textit{Argument :} Le capital social réduit les coûts de transaction, sécurise l'investissement local, préserve les communs.
      \item \textit{Exemples :} \emph{Le Lion et la perle} : le village se réunit, décide collectivement (place du marché) ; Cameroun : tontines (njangi) financent PME, agriculture ; forêts sacrées = biodiversité préservée ; spécialité technique : gestion communautaire de l'eau / électricité.
      \item \textit{Transition :} Au-delà des valeurs, les savoir-faire techniques traditionnels sont source d'innovation frugale.
      \end{itemize}
\item \textbf{B. Les savoir-faire endogènes (agriculture, architecture, pharmacopée, métallurgie) offrent des réponses adaptées aux défis locaux.}
      \begin{itemize}
      \item \textit{Idée directrice :} L'agroforesterie paysanne, la construction en banco, la pharmacopée, la transformation artisanale sont des technologies éprouvées, résilientes, peu carbonées.
      \item \textit{Argument :} L'innovation frugale (jugaad) part du besoin local et des ressources locales $\rightarrow$ développement soutenable.
      \item \textit{Exemples :} Cameroun : système \og ankara \fg (jachère améliorée), maisons en terre stabilisée (coût, confort thermique), plantes médicinales (Prunus africana) ; spécialité technique : matériaux locaux (banco, bambou) en BTP, biogaz à partir de déchets agricoles.
      \item \textit{Transition :} Tradition et modernité ne s'opposent pas : elles peuvent s'hybrider dans un projet de développement choisi.
      \end{itemize}
\end{itemize}

\bigskip
\textbf{III. Le développement durable suppose une articulation féconde entre tradition et modernité : ni rejet, ni mimétisme, mais réappropriation critique.}
\begin{itemize}
\item \textbf{A. La modernité technique et institutionnelle a besoin de l'ancrage culturel pour être légitime et efficace.}
      \begin{itemize}
      \item \textit{Idée directrice :} Les politiques publiques (éducation, santé, aménagement) qui intègrent les langues nationales, les autorités coutumières, les calendriers agricoles ont de meilleurs taux d'adoption.
      \item \textit{Argument :} L'appropriation sociale (ownership) conditionne la durabilité des projets techniques.
      \item \textit{Exemples :} Cameroun : écoles bilingues / programmes en langues nationales ; santé : collaboration tradipraticiens / médecins modernes ; spécialité technique : formation duale (entreprise + centre de formation) inspirée de l'apprentissage traditionnel.
      \item \textit{Transition :} Réciproquement, la tradition gagne à être interrogée, reformulée par les lumières de la science et du droit.
      \end{itemize}
\item \textbf{B. Une tradition critique, ouverte au dialogue des savoirs, devient le ferment d'une modernité africaine singulière.}
      \begin{itemize}
      \item \textit{Idée directrice :} La tradition n'est pas un musée : elle se réinvente par le droit (code de la famille), l'art (littérature, musique), la technique (startups agritech/fintech ancrées localement).
      \item \textit{Argument :} La créolisation culturelle (Glissant) / l'afropolitanisme (Mbembe) : modernité \emph{africaine} = hybridation assumée.
      \item \textit{Exemples :} \emph{Le Lion et la perle} : Soyinka écrit en anglais, forme occidentale, mais thème/yoruba/structure orale $\rightarrow$ œuvre hybride majeure ; Cameroun : startups (GiftedMom, Agro-Hub) = tech + savoir paysan ; littérature : Mongo Beti, Calixthe Beyala, Patrice Nganang.
      \item \textit{Transition vers conclusion :} La tradition n'est obstacle que si on la fige ; elle est levier si on la fait dialoguer avec le présent.
      \end{itemize}
\end{itemize}

\bigskip
\textbf{CONCLUSION} : Bilan synthétique (réponse à la problématique) — Ouverture (autre échelle, autre domaine, perspective prospective).

\end{exoresolu}

\begin{popfigure}
\begin{tikzpicture}[
  scale=0.75,
  every node/.style={font=\footnotesize, align=center},
  part/.style={draw=popblue, thick, fill=popblueL, rounded corners, inner sep=6pt, text width=3.2cm},
  sub/.style={draw=popgreen, thick, fill=popgreenL, rounded corners, inner sep=4pt, text width=2.8cm},
  ex/.style={draw=poporange, thick, fill=poporangeL, rounded corners, inner sep=3pt, text width=2.5cm},
  arrow/.style={->, thick, popdark}
]
% Sujet
\node[part, text width=5cm, fill=poppurpleL, draw=poppurple] (sujet) at (0,10) {Sujet : La tradition est-elle un obstacle au développement ?\\(Problématique : articulation tradition / développement endogène)};
% Parties
\node[part, align=center] (I) at (-7,6){I. Tradition comme obstacle\\ (conservatisme, instrumentalisation)};
\node[part, align=center] (II) at (0,6){II. Tradition comme ressource\\ (valeurs, savoir-faire endogènes)};
\node[part, align=center] (III) at (7,6){III. Articulation féconde\\ (hybridation, réappropriation critique)};
\draw[arrow] (sujet) -- (I);
\draw[arrow] (sujet) -- (II);
\draw[arrow] (sujet) -- (III);
% Sous-parties I
\node[sub, align=center] (IA) at (-9,3){A. Coutumes figées\\freinent innovation};
\node[sub, align=center] (IB) at (-5,3){B. Instrumentalisation\\politique de la coutume};
\draw[arrow] (I) -- (IA);
\draw[arrow] (I) -- (IB);
% Exemples I.A
\node[ex, align=center] (IA1) at (-10.5,0.5){Lakunle / dot\\(Cameroun: terre, mariages précoces)};
\node[ex, align=center] (IA2) at (-7.5,0.5){Exclusion femmes\\gestion foncière};
\draw[arrow] (IA) -- (IA1);
\draw[arrow] (IA) -- (IA2);
% Exemples I.B
\node[ex, align=center] (IB1) at (-6,0.5){Baroka / timbre-poste\\(pouvoir personnel)};
\node[ex] (IB2) at (-3,0.5) {Chefferie / redevances opacques (décentralisation)};
\draw[arrow] (IB) -- (IB1);
\draw[arrow] (IB) -- (IB2);
% Sous-parties II
\node[sub, align=center] (IIA) at (-2,3) {A. Valeurs : solidarité,\\ parole, communs};
\node[sub, align=center] (IIB) at (2,3) {B. Savoir-faire : agroforesterie,\\ banco, pharmacopée};
\draw[arrow] (II) -- (IIA);
\draw[arrow] (II) -- (IIB);
% Exemples II.A
\node[ex, align=center] (IIA1) at (-3.5,0.5) {Tontines (njangi)\\ financement PME};
\node[ex, align=center] (IIA2) at (-0.5,0.5) {Forêts sacrées =\\ biodiversité};
\draw[arrow] (IIA) -- (IIA1);
\draw[arrow] (IIA) -- (IIA2);
% Exemples II.B
\node[ex, align=center] (IIB1) at (0.5,0.5) {Ankara, banco,\\ biogaz (spéc. tech.)};
\node[ex, align=center] (IIB2) at (3.5,0.5) {Prunus africana,\\ innovation frugale};
\draw[arrow] (IIB) -- (IIB1);
\draw[arrow] (IIB) -- (IIB2);
% Sous-parties III
\node[sub, align=center] (IIIA) at (5,3){A. Modernité a besoin\\d'ancrage culturel};
\node[sub, align=center] (IIIB) at (9,3) {B. Tradition critique =\\ modernité africaine};
\draw[arrow] (III) -- (IIIA);
\draw[arrow] (III) -- (IIIB);
% Exemples III.A
\node[ex, align=center] (IIIA1) at (3.5,0.5) {Écoles langues\\ nationales};
\node[ex, align=center] (IIIA2) at (6.5,0.5){Santé : tradipraticiens\\+ médecins};
\draw[arrow] (IIIA) -- (IIIA1);
\draw[arrow] (IIIA) -- (IIIA2);
% Exemples III.B
\node[ex, align=center] (IIIB1) at (7.5,0.5) {Soyinka : hybride\\ anglais/yoruba};
\node[ex, align=center] (IIIB2) at (10.5,0.5){Startups agritech\\(GiftedMom, Agro-Hub)};
\draw[arrow] (IIIB) -- (IIIB1);
\draw[arrow] (IIIB) -- (IIIB2);
\end{tikzpicture}
\legende{Arbre du plan détaillé : du sujet (sommet) aux exemples (feuilles) en passant par les parties (bleu), sous-parties (vert) et arguments. Chaque feuille est un exemple précis, référencé.}
\end{popfigure}

\begin{popfigure}
\begin{tikzpicture}[
  scale=1.1,
  every node/.style={font=\small, align=center},
  funnel/.style={draw=popblue, thick, fill=popblue!10},
  label/.style={popdark, font=\small\bfseries}
]
% Entonnoir : trapèzes empilés (méthode robuste sans clip)
\foreach \y/\h/\wbase/\wtop/\txt/\col in {
  7/1.2/10/9.4/{Sujet brut : « La tradition est-elle un obstacle au développement ? »}/popblueL,
  5.5/1/7.5/6.9/{Mots-clés définis + présupposés identifiés}/popgreenL,
  4/1/5/4.4/{Problématique : « Dans quelle mesure la tradition peut-elle être levier d'un développement endogène ? »}/poporangeL,
  2.5/1/3.5/2.9/{Idées regroupées en 3 parties dialectiques}/poppurpleL,
  1/1/2/1.4/{Plan détaillé : idées directrices, arguments, exemples, transitions}/popgoldL}{
  \fill[\col] (-\wbase/2,\y-\h/2) -- (\wbase/2,\y-\h/2) -- (\wtop/2,\y+\h/2) -- (-\wtop/2,\y+\h/2) -- cycle;
  \draw[popblue, thick] (-\wbase/2,\y-\h/2) -- (\wbase/2,\y-\h/2) -- (\wtop/2,\y+\h/2) -- (-\wtop/2,\y+\h/2) -- cycle;
  \node[label] at (0,\y) {\txt};
}
% Flèches descendantes
\foreach \y in {6.4, 4.75, 3.25, 1.75} {
  \draw[->, thick, popdark] (0,\y+0.1) -- (0,\y-0.1);
}
% Légende latérale
\node[align=left, text width=5cm, popmuted] at (8,4) {
\textbf{Schéma en entonnoir} :\\
1. Sujet large $\rightarrow$ lecture active\\
2. Définitions / présupposés $\rightarrow$ clarification\\
3. Problématique $\rightarrow$ question précise\\
4. Regroupement $\rightarrow$ architecture\\
5. Plan détaillé $\rightarrow$ partition de rédaction
};
\end{tikzpicture}
\legende{Schéma en entonnoir : du sujet large au plan détaillé. Chaque étape rétrécit le champ et augmente la précision.}
\end{popfigure}

\begin{savaistu}[Le plan détaillé : la moitié de la bataille]
Au Probatoire / Baccalauréat technique, les correcteurs le confirment : \textbf{un plan détaillé rigoureux, équilibré, avec idées directrices affirmatives, arguments explicites et exemples précis, vaut à lui seul 6 à 8 points sur 20} (sur la note de dissertation). Pourquoi ? Parce qu'il prouve que le candidat a \emph{pensé} avant d'écrire. La rédaction n'est alors plus qu'un exercice de mise en forme (connecteurs, syntaxe, orthographe). \textbf{Conseil :} passez 45 à 60 minutes sur l'analyse + plan détaillé (sur 4h d'épreuve). Ne commencez la rédaction qu'une fois le plan \emph{validé} par vous-même (relecture critique : chaque partie répond-elle à la problématique ?).
\end{savaistu}

\begin{attention}[Hors-sujet : la faute capitale]
\begin{itemize}
\item \textbf{Répondre à côté de la problématique} (ex. : faire l'histoire de la tradition au lieu d'analyser ses rapports au développement) $\rightarrow$ note \emph{très faible}, même si le français est parfait.
\item \textbf{Confondre « tradition » et « coutumes néfastes »} sans nuance $\rightarrow$ plan manichéen, note moyenne.
\item \textbf{Oublier le « développement »} pour ne parler que de modernité $\rightarrow$ hors-sujet partiel.
\item \textbf{Vérification systématique :} à la fin du plan, relisez la problématique et demandez-vous : \og Ma partie I y répond-elle ? Ma partie II ? Ma partie III ? \fg Si non $\rightarrow$ revoir le plan \emph{avant} d'écrire.
\end{itemize}
\end{attention}

\begin{aretenir}[Check-list de validation du plan détaillé (brouillon)]
\begin{enumerate}
\item Le plan comporte-t-il 2 ou 3 parties équilibrées ?
\item Chaque partie a-t-elle 2 sous-parties (au minimum) ?
\item Chaque sous-partie a-t-elle : idée directrice (phrase) + argument(s) + exemple(s) précis + transition ?
\item Les exemples sont-ils variés (littérature \emph{Le Lion et la perle} + Cameroun + actualité + spécialité technique) ?
\item L'ordre des parties suit-il une logique (dialectique / analytique / thématique) explicite ?
\item \textbf{Test ultime :} si je ne lis que les idées directrices des 6 sous-parties, est-ce que j'ai la réponse complète à la problématique ?
\end{enumerate}
\end{aretenir}

\courssub{Lecture méthodique n°2 du \emph{Lion et la perle} : alimenter la banque d'exemples}

Dans le cadre de la réception des textes, la \textbf{lecture méthodique n°2} de \emph{Le Lion et la perle} (Wole Soyinka) vise précisément à \textbf{relever les arguments des personnages} pour constituer une \cle{banque d'exemples} réutilisable en dissertation. Voici la méthode opératoire :

\begin{methode}[Lecture méthodique n°2 : chasse aux arguments]
\begin{enumerate}
\item \textbf{Relire la pièce} (ou les extraits au programme) avec une grille de lecture : \emph{Personnage / Thème / Argument / Citation courte / Utilisation possible en dissertation}.
\item \textbf{Repérer les conflits d'idées :} Baroka vs Lakunle (tradition vs modernité), Sidi vs Sadiku (jeunesse vs ancienneté), individuel vs collectif.
\item \textbf{Noter les arguments \emph{pour} et \emph{contre} chaque thèse :} la pièce ne tranche pas, elle \emph{met en scène} le débat. C'est une mine pour la dialectique.
\item \textbf{Classer par thématiques transversales :} pouvoir, genre, éducation, économie, langage, rite, technique, identité.
\item \textbf{Rédiger des « fiches-exemples » prêtes à l'emploi :} \og Exemple \emph{Le Lion et la perle} : Baroka utilise le timbre-poste (modernité technique) pour asseoir son autorité traditionnelle $\rightarrow$ illustration de l'hybridation féconde / instrumentalisation. \fg
\end{enumerate}
\end{methode}

\begin{exemplebox}[Extrait de grille de lecture méthodique n°2 (à compléter par l'élève)]
\begin{center}
\begin{tabularx}{\linewidth}{|l|l|X|l|X|}\hline
\textbf{Personnage} & \textbf{Thème} & \textbf{Argument (résumé)} & \textbf{Citation / Scène} & \textbf{Utilisation dissertation} \\ \hline
Baroka & Poder / ruse & La tradition sait absorber la modernité pour se renforcer. & Scène du timbre-poste (acte 2) & I.B (instrumentalisation) / III.A (hybridation) \\ \hline
Lakunle & Progrès / éducation & La modernité libère l'individu des chaînes coutumières. & Refus de la dot, discours sur la machine (acte 1) & I.A (figement) / III.B (modernité dogmatique) \\ \hline
Sidi & Identité / image & La fierté culturelle vaut plus que la séduction moderne. & Choix final pour Baroka (acte 3) & II.A (valeurs) / III.B (réappropriation) \\ \hline
Sadiku & Transmission / rumeur & La parole traditionnelle circule, régule, manipule. & Annonce de l'impuissance de Baroka (acte 2) & II.A (communication, régulation sociale) \\ \hline
\end{tabularx}
\end{center}
\textbf{Consigne :} Complétez cette grille sur l'ensemble de la pièce. Elle deviendra votre \og boîte à outils \fg pour toute dissertation sur tradition / modernité / développement / identité.
\end{exemplebox}

\begin{exercice}[Entraînement : analyse de sujet et plan détaillé]
Traitez le sujet suivant en appliquant les 6 étapes : \og \textbf{L'école technique peut-elle former des citoyens ?} \fg
\begin{enumerate}
\item Soulignez mots-clés et mots-outils. Définissez les termes. Repérez les présupposés.
\item Formulez 3 problématiques possibles ; choisissez la meilleure ; justifiez.
\item Faites un brainstorming de 15 idées / exemples (dont 3 tirés de votre spécialité technique, 2 de \emph{Le Lion et la perle}, 2 de la vie camerounaise).
\item Proposez un plan détaillé (parties, sous-parties, idées directrices, arguments, exemples, transitions).
\end{enumerate}
\end{exercice}

\begin{corrige}[Exercice n°1 — Éléments de correction]
\textbf{1. Mots-clés :} école technique, former, citoyens. \textbf{Mots-outils :} peut-elle (modalité, capacité, pas obligation). \textbf{Définitions :} école technique = formation professionnelle / technologique ; former = instruire, socialiser, qualifier ; citoyen = sujet de droits/devoirs, acteur de la cité, conscience critique. \textbf{Présupposés :} l'école technique forme (déjà) des techniciens ; la citoyenneté s'apprend ; technique et citoyenneté ne vont pas de soi ensemble.

\textbf{2. Problématiques :} a) L'école technique forme-t-elle des citoyens ? (trop fermée) b) Quel rôle pour l'école technique dans la formation du citoyen ? (trop large) c) \textbf{Retenue :} \og Dans quelle mesure la formation technique, par sa pédagogie et ses valeurs, contribue-t-elle à l'émergence d'une citoyenneté active et responsable ? \fg (tension technique/citoyenneté, place pour nuance).

\textbf{3. Brainstorming (échantillon) :} Spécialité : travail en équipe (chantier), respect normes (sécurité), éthique pro ; \emph{Lion et la perle} : Lakunle instituteur = figure de l'éducateur moderne mais déconnecté ; Baroka = formation par l'apprentissage traditionnel ; Cameroun : insertion pro des jeunes, civisme fiscal, bénévolat technique (électrification village) ; Actualité : service national, volontariat, ODD 4/8/11.

\textbf{4. Plan détaillé (schéma) :} I. L'école technique : un vivier de vertus civiques (A. Pédagogie de projet = responsabilité collective ; B. Éthique professionnelle = respect règles / bien commun). II. Les limites : risque de technicisme sans conscience critique (A. Spécialisation étroite = déconnexion sociale ; B. Hiérarchie technique = soumission à l'autorité). III. Vers une ingénierie citoyenne : réformer les programmes (A. Modules citoyenneté / droit / environnement ; B. Projets d'utilité publique intégrés au cursus ; C. Partenariats collectivités / entreprises / ONG).
\end{corrige}

\coursec{Rédiger l'introduction}

Après avoir analysé le sujet et construit un plan détaillé — étape indispensable que nous avons travaillée dans la section précédente —, vous disposez désormais de la « carte » de votre dissertation. Il reste à en écrire la « porte d'entrée » : l'introduction. C'est la première chose que lit le correcteur ; elle conditionne son attention et sa bienveillance. Une introduction réussie n'est pas un exercice de style gratuit, c'est un contrat de lecture : vous annoncez ce que vous allez démontrer et comment vous allez le faire.

\courssub{La fonction de l'introduction : capter, poser, orienter}

\begin{definition}[Rôle de l'introduction]
L'introduction est le premier paragraphe de la dissertation. Elle remplit trois fonctions indissociables :
\begin{enumerate}
    \item \textbf{Capter l'attention} (l'amorce) : montrer que le sujet intéresse le vivant, le concret, l'actualité.
    \item \textbf{Poser le problème} (définition du sujet + problématique) : délimiter le champ de réflexion et formuler la question précise à laquelle la dissertation répondra.
    \item \textbf{Orienter la lecture} (annonce du plan) : donner au correcteur la structure logique du développement pour qu'il suive sans effort le fil de la démonstration.
\end{enumerate}
\end{definition}

\begin{aretenir}[L'image de l'entonnoir]
L'introduction suit un mouvement \textbf{du général vers le particulier}. On part d'un fait large (actualité, citation, constat) pour resserrer progressivement jusqu'à la problématique précise et au plan détaillé. C'est la forme d'un entonnoir.
\end{aretenir}

\begin{popfigure}
\begin{tikzpicture}[scale=0.9, every node/.style={font=\small}]
    % Couleurs
    \colorlet{ent1}{popblue!80}
    \colorlet{ent2}{poppurple!80}
    \colorlet{ent3}{poporange!80}
    \colorlet{ent4}{popgreen!80}
    % Forme entonnoir (trapèzes superposés)
    \def\topwidth{12}
    \def\botwidth{2}
    \def\layerh{2.2}
    \def\gap{0.15}
    % Couche 1 : Amorce
    \fill[ent1] (-\topwidth/2, 3*\layerh+2*\gap) -- (\topwidth/2, 3*\layerh+2*\gap) 
        -- (\topwidth/2*0.7, 2*\layerh+2*\gap) -- (-\topwidth/2*0.7, 2*\layerh+2*\gap) -- cycle;
    \node[white, align=center, font=\bfseries] at (0, 2.5*\layerh+2*\gap) {1. AMORCE (Accroche)\\\textit{Fait d'actualité, citation, constat,}\\\textit{ anecdote...}\\\textit{\textbf{But :} Captiver}};
    % Couche 2 : Présentation du sujet
    \fill[ent2] (-\topwidth/2*0.7, 2*\layerh+2*\gap) -- (\topwidth/2*0.7, 2*\layerh+2*\gap) 
        -- (\topwidth/2*0.4, \layerh+\gap) -- (-\topwidth/2*0.4, \layerh+\gap) -- cycle;
    \node[white, align=center, font=\bfseries] at (0, 1.5*\layerh+1.5*\gap) {2. PRÉSENTATION /\\DÉFINITION DU SUJET\\\textit{Délimiter, définir les termes,}\\\textit{ reformuler le sujet}\\\textit{\textbf{But :} Situer le débat}};
    % Couche 3 : Problématique
    \fill[ent3] (-\topwidth/2*0.4, \layerh+\gap) -- (\topwidth/2*0.4, \layerh+\gap) 
        -- (\topwidth/2*0.15, \gap) -- (-\topwidth/2*0.15, \gap) -- cycle;
    \node[white, align=center, font=\bfseries] at (0, 0.5*\layerh+0.5*\gap) {3. PROBLÉMATIQUE\\\textit{Question directe ou indirecte}\\\textit{Sans y répondre !}\\\textit{\textbf{But :} Poser la question centrale}};
    % Couche 4 : Annonce du plan
    \fill[ent4] (-\topwidth/2*0.15, \gap) -- (\topwidth/2*0.15, \gap) 
        -- (0, -\layerh) -- (0, -\layerh) -- cycle; % pointe
    \node[white, align=center, font=\bfseries] at (0, -0.5*\layerh) {4. ANNONCE DU PLAN\\\textit{« Nous verrons d'abord...}\\\textit{ ensuite... enfin... »}\\\textit{\textbf{But :} Guider le lecteur}};
    % Flèches latérales "Général -> Particulier"
    \draw[popmuted, thick, <->] (-7, 3*\layerh+2*\gap) -- (-7, -\layerh) node[midway, left, rotate=90, popmuted] {Du général $\rightarrow$ Au particulier};
\end{tikzpicture}
\legende{Schéma de l'introduction en entonnoir : les 4 étapes obligatoires, du large (amorce) au précis (plan).}
\end{popfigure}

\courssub{Étape 1 : L'amorce (ou accroche) — ouvrir sur le concret}

L'amorce est la première phrase (ou les deux premières phrases). Elle doit être \textbf{brève} (2 à 3 lignes maximum), \textbf{en lien direct avec le sujet} et \textbf{originale}. Interdiction formelle des lieux communs vides.

\begin{methode}[Les quatre procédés d'amorce efficaces]
\begin{enumerate}
    \item \textbf{Le fait d'actualité camerounais (recommandé au Probatoire) :} Ancre le sujet dans la réalité locale, montre votre culture générale.
    \item \textbf{La citation d'auteur :} Brève, exacte, référencée (auteur, œuvre). Elle doit éclairer le sujet, pas le paraphraser.
    \item \textbf{Le constat général / fait de société :} Observation partagée, vérité d'expérience, phénomène observable par tous.
    \item \textbf{L'anecdote ou l'exemple concret :} Micro-récit, situation vécue, exemple précis qui incarne l'abstraction du sujet.
\end{enumerate}
\end{methode}

\begin{exemplebox}[Amorces variées sur le sujet : « La tradition est-elle un frein au progrès ? »]
\begin{itemize}
    \item \textbf{Actualité CM :} « Au Cameroun, la digitalisation de l'état civil (e-EtatCivil) heurte parfois la gestion coutumière des noms et des successions, illustrant la tension entre modernité administrative et traditions orales. »
    \item \textbf{Citation :} « La tradition, ce n'est pas le culte des cendres, mais la conservation du feu », écrivait Gustav Mahler, suggérant que le passé alimente l'avenir plutôt qu'il ne l'entrave.
    \item \textbf{Constat :} « Toute société humaine se construit sur un héritage — langue, rites, techniques — qu'elle transmet ; pourtant, l'histoire montre que s'y enfermer peut mener à l'immobilisme. »
    \item \textbf{Anecdote :} « Dans un village du Nord, un jeune ingénieur agronome a dû adapter la technique du goutte-à-goutte au calendrier des ancêtres pour que les paysans l'adoptent : la tradition a guidé le progrès au lieu de le bloquer. »
\end{itemize}
\end{exemplebox}

\begin{attention}[Amorces interdites — Pièges à éviter]
\begin{itemize}
    \item \textbf{« Depuis la nuit des temps... »} / « Depuis l'aube de l'humanité... » : vide, banal, hors-sujet immédiat.
    \item \textbf{« Le dictionnaire Larousse définit... »} : ce n'est pas une amorce, c'est une définition (étape 2). De plus, le correcteur n'a pas besoin de votre dictionnaire.
    \item \textbf{Recopier le sujet :} « La tradition est-elle un frein au progrès ? C'est la question que nous allons traiter. » $\rightarrow$ \textbf{Zéro point} pour l'amorce.
    \item \textbf{Amorce hors-sujet :} parler de la mode vestimentaire pour un sujet sur la tradition juridique.
\end{itemize}
\end{attention}

\courssub{Étape 2 : Présentation et définition du sujet — délimiter le terrain}

Après l'amorce, vous \textbf{reprenez le sujet}, vous en \textbf{définez les termes clés} (un ou deux maximum) et vous en \textbf{reformulez la portée}. C'est le moment de montrer que vous avez compris les nuances.

\begin{methode}[Définir et reformuler]
\begin{enumerate}
    \item Identifiez les \textbf{mots-clés} du sujet (ex : « tradition », « progrès », « frein »).
    \item Donnez une \textbf{définition opérationnelle} (adaptée au sujet, pas la définition dictionnaire brute).
    \item \textbf{Reformulez le sujet} en une phrase affirmative qui en montre l'enjeu : « Il s'agit de savoir si... », « La question posée est de déterminer dans quelle mesure... »
\end{enumerate}
\end{methode}

\begin{exemplebox}[Définition / Reformulation sur « La tradition est-elle un frein au progrès ? »]
« La \textbf{tradition} désigne ici l'ensemble des coutumes, savoirs et valeurs transmis de génération en génération qui structurent l'identité collective. Le \textbf{progrès} s'entend comme l'amélioration matérielle, technique ou morale des conditions de vie. Le terme \textbf{frein} suppose un obstacle, un ralentissement, voire un blocage. \textbf{Il s'agit donc d'examiner si l'héritage culturel constitue nécessairement un obstacle à l'innovation, ou s'il peut au contraire en être le socle ou le régulateur.} »
\end{exemplebox}

\courssub{Étape 3 : La problématique — la question moteur}

C'est le \textbf{cœur stratégique} de l'introduction. La problématique n'est \textbf{PAS} le sujet recopié sous forme de question. C'est une \textbf{question précise, débattue, issue de votre analyse}, qui guide tout le développement.

\begin{propriete}[Caractéristiques d'une bonne problématique]
\begin{itemize}
    \item Elle est \textbf{unique} (une seule question principale).
    \item Elle est \textbf{formulée en question} (directe : « Comment... ? » / « Pourquoi... ? » ; ou indirecte : « Nous nous demanderons dans quelle mesure... »).
    \item Elle \textbf{ne contient pas la réponse} (pas de « ne pensez-vous pas que... », pas de « n'est-il pas vrai que... »).
    \item Elle \textbf{appelle une démonstration en plusieurs parties} (votre plan y répond point par point).
    \item Elle \textbf{contient une tension} (opposition, paradoxe, nuances).
\end{itemize}
\end{propriete}

\begin{exemplebox}[Sujet vs Problématique]
\begin{center}
\begin{tabularx}{\linewidth}{|l|X|}\hline
\textbf{Sujet} & \textbf{Problématique (exemple)} \\ \hline
La tradition est-elle un frein au progrès ? & \textit{Dans quelle mesure la tradition, loin de s'opposer au progrès, peut-elle en orienter le sens et en garantir l'ancrage social ?} \\ \hline
Le progrès technique suffit-il à faire le bonheur d'une société ? & \textit{Le progrès technique est-il une condition suffisante du bonheur collectif, ou le bonheur dépend-il prioritairement de choix éthiques, politiques et culturels que la technique ne saurait remplacer ?} \\ \hline
Faut-il tout dire pour être honnête ? & \textit{La franchise absolue est-elle la marque de l'honnêteté, ou l'honnêteté exige-t-elle parfois le silence, la retenue ou le mensonge par humanité ?} \\ \hline
\end{tabularx}
\end{center}
\end{exemplebox}

\begin{attention}[Erreurs fatales sur la problématique]
\begin{itemize}
    \item \textbf{Absence de problématique :} passer de la définition à l'annonce du plan. $\rightarrow$ Dissertation sans direction, note plafond 6/20.
    \item \textbf{Problématique = Sujet :} « La tradition est-elle un frein au progrès ? » recopiée. $\rightarrow$ Montre que vous n'avez pas analysé.
    \item \textbf{Problématique fermée / binaire :} « La tradition freine-t-elle le progrès ? » (Réponse : Oui/Non). $\rightarrow$ Plan binaire pauvre (Pour/Contre).
    \item \textbf{Problématique qui répond à côté :} sujet sur « tradition/progrès », problématique sur « tradition/modernité » (concepts différents).
\end{itemize}
\end{attention}

\courssub{Étape 4 : L'annonce du plan — la carte du développement}

L'annonce du plan doit être \textbf{claire, parallèle, fidèle au développement}. Elle utilise des connecteurs logiques et des formules de transition standardisées.

\begin{methode}[Formules d'annonce du plan — Modèles prêts à l'emploi]
\begin{enumerate}
    \item \textbf{Classique (3 parties) :} « Nous analyserons d'abord \textit{[idée partie 1]}, avant d'examiner \textit{[idée partie 2]}, pour enfin montrer \textit{[idée partie 3]}. »
    \item \textbf{Dialectique (Thèse/Antithèse/Synthèse) :} « Nous verrons dans un premier temps que \textit{[thèse]}, puis que \textit{[antithèse]}, ce qui nous conduira à soutenir que \textit{[synthèse]}. »
    \item \textbf{Analytique (Causes/Conséquences/Solutions) :} « Il conviendra d'identifier les causes de \textit{[phénomène]}, d'en mesurer les conséquences, puis d'esquisser des pistes de résolution. »
\end{enumerate}
\emph{Règle d'or :} Les libellés annoncés doivent reprendre \textbf{exactement} les idées directrices de vos parties (pas de surprise au développement).
\end{methode}

\begin{exemplebox}[Annonce du plan pour la problématique : « Dans quelle mesure la tradition... peut-elle orienter le sens du progrès ? »]
« \textbf{Nous verrons d'abord} que la tradition peut apparaître comme un frein lorsque elle fige les mentalités et rejette l'innovation (I). \textbf{Nous montrerons ensuite} qu'elle offre au contraire des repères éthiques et une cohésion sociale indispensables à tout progrès durable (II). \textbf{Nous démontrerons enfin} que les sociétés les plus dynamiques sont celles qui opèrent une \textbf{relecture critique} de leur tradition pour l'adapter aux défis contemporains (III). »
\end{exemplebox}

\courssub{Entraînement comparatif : introduction ratée vs introduction réussie}

Voici deux introductions complètes sur le même sujet type Probatoire, inspirées des thèmes du \emph{Lion et la Perle} (tradition vs modernité, pouvoir, séduction, changement). Lisez-les, comparez, repérez les critères vus plus haut.

\begin{popfigure}
\begin{tikzpicture}[scale=0.85, every node/.style={font=\small}]
    % Styles
    \tikzset{
        boxrat/.style={draw=poppink, thick, fill=poppink!10, rounded corners, text width=7.5cm, minimum height=6.5cm},
        boxok/.style={draw=popgreen, thick, fill=popgreen!10, rounded corners, text width=7.5cm, minimum height=6.5cm},
        titlebad/.style={fill=poppink, text=white, font=\bfseries, rounded corners=2pt, inner sep=2pt},
        titlegood/.style={fill=popgreen, text=white, font=\bfseries, rounded corners=2pt, inner sep=2pt},
        annot/.style={font=\tiny, text=popdark, align=left},
        goodannot/.style={font=\tiny, text=popgreen!80!black, align=left},
        badannot/.style={font=\tiny, text=poppink!80!black, align=left},
    }
    % BOÎTE GAUCHE : RATÉE
    \node[boxrat] (bad) at (-8.5, 0) {};
    \node[titlebad, anchor=north west] at (bad.north west) { INTRODUCTION RATÉE (Note ~4/20) };
    \node[annot, anchor=north west, text width=7cm] at ([xshift=0.3cm, yshift=-0.8cm]bad.north west) {
        \textbf{« Depuis la nuit des temps, les hommes ont des traditions.} \textcolor{poppink}{[Amorce vide, cliché, hors-sujet]}\\
        \textbf{Le sujet demande si la tradition est un frein au progrès.} \textcolor{poppink}{[Recopie du sujet = pas d'amorce]}\\
        \textbf{La tradition, c'est ce qu'on transmet. Le progrès, c'est l'évolution.} \textcolor{poppink}{[Définitions approximatives, pas de reformulation]}\\
        \textbf{Est-ce que la tradition empêche le progrès ?} \textcolor{poppink}{[Problématique = sujet, binaire, pas de tension]}\\
        \textbf{Nous allons voir les arguments pour et contre.} \textcolor{poppink}{[Plan vague, « pour/contre », pas d'idées directrices]}\\
        \textbf{En conclusion, on verra.} \textcolor{poppink}{[Conclusion annoncée dans l'intro, maladroit]} » 
    };
    % Annotations spécifiques (flèches)
    \draw[poppink, thick, ->] (-12.5, 1.8) -- (-11.2, 1.5) node[badannot, right, text width=3cm] {Amorce « nuit des temps» : \textbf{interdite}};
    \draw[poppink, thick, ->] (-12.5, 0.5) -- (-11.2, 0.2) node[badannot, right, text width=3cm] {Sujet recopié : \textbf{zéro point}};
    \draw[poppink, thick, ->] (-12.5, -0.8) -- (-11.2, -1.1) node[badannot, right, text width=3cm] {Problématique absente / binaire};
    \draw[poppink, thick, ->] (-12.5, -2.1) -- (-11.2, -2.4) node[badannot, right, text width=3cm] {Plan « pour/contre » : \textbf{plan au niveau 6e}};
    % BOÎTE DROITE : RÉUSSIE
    \node[boxok] (good) at (8.5, 0) {};
    \node[titlegood, anchor=north west] at (good.north west) { INTRODUCTION RÉUSSIE (Note ~16/20) };
    \node[annot, anchor=north west, text width=7cm] at ([xshift=0.3cm, yshift=-0.8cm]good.north west) {
        \textbf{« Au Cameroun, l'instauration du Code des personnes et de la famille en 1981 a suscité de vifs débats :} \textcolor{popgreen}{[Amorce actualité CM, concrète]}\\
        \textbf{certains y voyaient une remise en cause des coutumes successorales, d'autres une protection des femmes et des enfants.} \textcolor{popgreen}{[Mise en tension immédiate]}\\
        \textbf{La \textit{tradition}, ensemble de normes coutumières transmises oralement, structure l'identité collective ; le \textit{progrès}, ici juridique, vise l'égalité et la justice sociale.} \textcolor{popgreen}{[Définitions opérationnelles, termes clés]}\\
        \textbf{\textit{Dans quelle mesure la tradition, loin de s'opposer au progrès, peut-elle en orienter le sens et en garantir l'ancrage social ?}} \textcolor{popgreen}{[Problématique nuancée, tension, appelle 3 parties]}\\
        \textbf{Nous analyserons d'abord les résistances traditionnelles au changement législatif (I),} \textcolor{popgreen}{[Annonce I : idée directrice claire]}\\
        \textbf{puis nous montrerons comment le droit moderne s'inspire des valeurs coutumières pour être légitime (II),} \textcolor{popgreen}{[Annonce II : nuance, lien]}\\
        \textbf{enfin nous défendrons l'idée d'un dialogue fécond entre coutume et loi pour un progrès partagé (III).} \textcolor{popgreen}{[Annonce III : synthèse, ouverture]}\\
        » 
    };
    % Annotations positives
    \draw[popgreen!70!black, thick, ->] (12.5, 1.8) -- (11.2, 1.5) node[goodannot, left, text width=3cm] {Amorce ancrée, fait daté, tension};
    \draw[popgreen!70!black, thick, ->] (12.5, 0.2) -- (11.2, -0.1) node[goodannot, left, text width=3cm] {Définitions utiles, reformulation};
    \draw[popgreen!70!black, thick, ->] (12.5, -1.1) -- (11.2, -1.4) node[goodannot, left, text width=3cm] {Problématique \textbf{unique, ouverte, nuancée}};
    \draw[popgreen!70!black, thick, ->] (12.5, -2.4) -- (11.2, -2.7) node[goodannot, left, text width=3cm] {Plan \textbf{parallèle, fidèle, 3 idées fortes}};
\end{tikzpicture}
\legende{Comparatif annoté : introduction ratée (gauche) vs introduction réussie (droite) sur un sujet « tradition / progrès » à coloration camerounaise.}
\end{popfigure}

\courssub{Exemple complet commenté : sujet d'entraînement}

Sujet : \textbf{« Le progrès technique suffit-il à faire le bonheur d'une société ? »}

\begin{exoresolu}[Introduction rédigée et commentée]
\textbf{Énoncé :} Rédigez l'introduction de la dissertation sur ce sujet.

\textbf{Production de l'élève (version finale) :}
« \textit{En 2023, le Cameroun a lancé le projet « Smart Villages » pour connecter les zones rurales à l'internet haut débit : une prouesse technique saluée comme vecteur de développement. Pourtant, dans plusieurs localités, les équipements restent inutilisés par manque de formation, de maintenance ou d'adaptation aux besoins réels (éducation, santé, marchés). Le \textbf{progrès technique} désigne l'ensemble des innovations scientifiques et industrielles qui transforment les modes de production et de vie. Le \textbf{bonheur d'une société} ne saurait se réduire au confort matériel ; il implique la justice sociale, la cohésion, la liberté et le sens donné à l'existence. \textbf{La maîtrise technique est-elle une condition suffisante du bonheur collectif, ou le bonheur dépend-il de choix politiques, éthiques et culturels que la technique ne saurait remplacer ?} Nous examinerons d'abord les apports indéniables de la technique au bien-être matériel (I), avant d'analyser ses limites et ses dérives lorsqu'elle s'affranchit du projet humain (II), pour enfin soutenir que le bonheur social naît d'une \textbf{technique encadrée par une finalité éthique et démocratique} (III). »}

\tcblower

\textbf{Commentaire détaillé (grille d'évaluation) :}
\begin{itemize}
    \item \textbf{Amorce (lignes 1-2) :} Fait d'actualité camerounais précis (Smart Villages 2023), concret, qui \textbf{incarne} le sujet (technique vs usage réel). $\rightarrow$ \textbf{Réussi.}
    \item \textbf{Définitions / Reformulation (lignes 3-4) :} Mots-clés définis (\textit{progrès technique}, \textit{bonheur d'une société}) avec élargissement du second (pas que matériel). Reformulation implicite dans la phrase « ne saurait se réduire... ». $\rightarrow$ \textbf{Réussi.}
    \item \textbf{Problématique (ligne 5) :} Question indirecte, unique, nuancée (« condition suffisante » vs « choix politiques/éthiques »), tension claire. Elle \textbf{ne répond pas}. $\rightarrow$ \textbf{Réussi.}
    \item \textbf{Annonce du plan (lignes 6-7) :} Trois idées directrices distinctes, progressives (apports $\rightarrow$ limites $\rightarrow$ synthèse encadrée), vocabulaire parallèle (« examiner / analyser / soutenir »). Fidèle à un plan dialectique. $\rightarrow$ \textbf{Réussi.}
    \item \textbf{Langue / Présentation :} Paragraphe unique, connecteurs logiques (\textit{Pourtant, Or, avant de, pour enfin}), zéro faute, ton objectif. $\rightarrow$ \textbf{Réussi.}
\end{itemize}
\end{exoresolu}

\courssub{Exercices d'entraînement progressif}

\begin{exercice}[Exercice 1 — Identifier les 4 étapes]
Repérez dans l'introduction ci-dessous les 4 étapes (Amorce, Définition/Reformulation, Problématique, Annonce du plan) en les surlignant de 4 couleurs différentes.
« \textit{Lors de la CAN 2021 au Cameroun, l'arbitrage vidéo (VAR) a suscité autant d'espoirs que de controverses : si la technique a corrigé des erreurs flagrantes, elle a aussi ralenti le jeu et suscité de nouvelles incompréhensions. La \textbf{technique} est l'application rationnelle de connaissances scientifiques à des fins pratiques ; la \textbf{justice sportive} vise l'équité du résultat par le respect des règles. \textbf{L'introduction de la technologie dans le sport garantit-elle une justice plus parfaite, ou déplace-t-elle simplement le problème de l'erreur humaine vers l'interprétation technique ?} Nous étudierons d'abord les apports de la technique à la vérité du jeu (I), puis nous questionnerons les nouvelles zones d'ombre qu'elle crée (II), pour enfin réfléchir aux conditions d'une technologie au service de l'esprit sportif (III). } »
\end{exercice}

\begin{corrige}[Exercice 1]
\begin{itemize}
    \item \textcolor{popblue}{\textbf{Amorce :}} « Lors de la CAN 2021 au Cameroun, l'arbitrage vidéo (VAR) a suscité autant d'espoirs que de controverses... »
    \item \textcolor{poppurple}{\textbf{Définition/Reformulation :}} « La \textbf{technique} est l'application rationnelle... ; la \textbf{justice sportive} vise l'équité... »
    \item \textcolor{poporange}{\textbf{Problématique :}} « \textbf{L'introduction de la technologie dans le sport garantit-elle une justice plus parfaite, ou déplace-t-elle simplement le problème... ?} »
    \item \textcolor{popgreen}{\textbf{Annonce du plan :}} « Nous étudierons d'abord... (I), puis nous questionnerons... (II), pour enfin réfléchir... (III). »
\end{itemize}
\end{corrige}

\begin{exercice}[Exercice 2 — Rédiger votre introduction]
À partir du plan détaillé suivant (issu de la section précédente), rédigez l'introduction complète sur le sujet : \textbf{« Le Lion et la Perle : la tradition peut-elle s'adapter sans se trahir ? »}

\textbf{Plan détaillé fourni :}
\begin{itemize}
    \item \textbf{I. L'affrontement :} Baroka (tradition vivante, ruse) vs Lakunle (modernité mimétique, stérile) $\rightarrow$ la tradition n'est pas figée, elle sait user de la modernité.
    \item \textbf{II. La femme (Sidi) comme enjeu :} ni objet de tradition, ni trophée de modernité ; son choix révèle la nécessité d'une \textbf{synthèse personnelle}.
    \item \textbf{III. Le village d'Ilujinle :} espace de \textbf{négociation continue} où la tradition intègre le nouveau (danse du développement, timbre-poste) pour survivre.
\end{itemize}

\textbf{Consignes :} Amorce camerounaise (fait culturel/villageois), définitions de « tradition » et « adaptation/trahison », problématique nuancée, annonce du plan parallèle.
\end{exercice}

\begin{corrige}[Exercice 2 — Proposition de correction]
« \textit{Dans le Grassland camerounais, la chefferie traditionnelle gère aujourd'hui des coopératives agricoles et des micro-finances : la coutume gère le moderne sans s'y dissoudre. La \textbf{tradition}, chez Soyinka comme chez nous, n'est pas un musée mais un \textbf{corps vivant} de rites, de langage et d'autorité ; \textbf{s'adapter} c'est intégrer du nouveau pour durer, \textbf{se trahir} c'est perdre son sens profond au profit de l'apparence. \textbf{La tradition doit-elle se figer pour survivre, ou sa vitalité réside-t-elle précisément dans sa capacité à métaboliser le changement sans perdre son âme ?} Nous montrerons d'abord que la tradition, incarnée par Baroka, manie la modernité comme un outil sans en adopter l'idéologie (I). Nous analyserons ensuite comment le personnage de Sidi refuse les deux modèles purs pour inventer une voie personnelle (II). Nous conclurons qu'Ilujinle illustre une tradition \textbf{performative} : elle se réinvente par la danse, le rire et le symbole pour rester elle-même (III). »}
\end{corrige}

\courssub{Bilan et passage au développement}

\begin{aretenir}[Check-list de l'introduction — À cocher avant de passer au développement]
\begin{itemize}
    \item $\square$ \textbf{Amorce :} 2-3 lignes, ancrée (actualité CM, citation, constat, anecdote), \textbf{pas de cliché}.
    \item $\square$ \textbf{Définitions :} 1 à 2 termes clés définis \textbf{opératoirement} (adaptés au sujet).
    \item $\square$ \textbf{Reformulation :} Une phrase qui montre l'enjeu, la tension du sujet.
    \item $\square$ \textbf{Problématique :} \textbf{Une seule question}, ouverte, nuancée, \textbf{sans réponse}, issue de votre analyse.
    \item $\square$ \textbf{Annonce du plan :} 3 parties, libellés \textbf{parallèles}, idées directrices \textbf{claires}, connecteurs (\textit{d'abord, ensuite, enfin} / \textit{nous verrons que..., puis que..., pour montrer que...}).
    \item $\square$ \textbf{Forme :} \textbf{Un seul paragraphe}, saut de ligne après, langue soignée (pas de familier, pas de « je »).
\end{itemize}
\end{aretenir}

\begin{savaistu}[Le correcteur et l'introduction]
Les correcteurs du Probatoire le confirment : \textbf{l'introduction se lit en 30 secondes} et détermine souvent la note finale. Une introduction soignée (propre, structurée, sans faute, avec vraie problématique) place le correcteur en \textbf{disposition favorable} : il lira le développement en cherchant \textbf{la confirmation} de vos annonces, pas les erreurs. À l'inverse, une introduction bâclée crée un \textbf{à priori négatif} difficile à rattraper. Soignez-la comme une vitrine.
\end{savaistu}

\begin{methode}[Rituel d'examen : 10 minutes pour l'introduction]
\begin{enumerate}
    \item \textbf{2 min :} Relisez votre plan détaillé (idées directrices I, II, III).
    \item \textbf{2 min :} Choisissez et rédigez l'amorce (préférez actualité CM ou citation courte).
    \item \textbf{2 min :} Écrivez définitions + reformulation + problématique \textbf{d'un trait} (ne cherchez pas le mot parfait, visez la justesse logique).
    \item \textbf{2 min :} Recopiez l'annonce du plan \textbf{en reprenant exactement les libellés de votre plan}.
    \item \textbf{2 min :} Relisez à voix basse : fluidité ? cohérence ? pas de répétition ? saut de ligne après ?
\end{enumerate}
\end{methode}

\courssub{Transition vers la section 4}

Votre introduction est écrite, votre plan est annoncé. Le correcteur sait \textbf{où} vous allez. Il attend maintenant de voir \textbf{comment} vous y allez. La section 4 — « Rédiger le développement : le paragraphe argumentatif » — vous donne la méthode pas à pas pour transformer chaque idée directrice en un paragraphe solide : idée directrice $\rightarrow$ argument(s) $\rightarrow$ exemple(s) $\rightarrow$ mini-conclusion. C'est le muscle de la dissertation.

\coursec{Rédiger le développement : le paragraphe argumentatif}

Après avoir maîtrisé l'art de l'introduction — qui pose le sujet, le problématise et annonce le plan — nous entrons dans le cœur de la dissertation : le développement. C'est ici que se joue la convaincre. Chaque paragraphe n'est pas une simple suite de phrases ; c'est une \cle{unité de sens} architecturée pour faire avancer la démonstration. Un paragraphe mal construit, c'est une brique fissurée dans l'édifice : l'ensemble risque de s'effondrer sous le feu de la critique. Dans cette section, nous allons disséquer l'anatomie du paragraphe type, outiller votre syntaxe avec les types de phrases, enrichir votre vocabulaire par la synonymie et l'antonymie, et enfin analyser un modèle littéraire d'excellence : le duel verbal entre Lakunle et Baroka dans \emph{Le Lion et la perle}.

\courssub{La structure du paragraphe argumentatif : le modèle I-A-E-L}

Tout paragraphe argumentatif efficace répond à une logique implacable en quatre temps. Oubliez l'idée reçue selon laquelle « écrire, c'est laisser venir les idées ». En dissertation, \textbf{on ne subit pas sa pensée, on la construit}.

\begin{definition}[L'idée directrice (ou phrase d'accroche)]
L'\cle{idée directrice} est la phrase qui ouvre le paragraphe et en annonce l'argument principal. Elle fonctionne comme une \cle{mini-thèse} locale : elle dit « voici ce que je vais prouver dans ce paragraphe ». Elle doit être claire, affirmative et directement reliée à l'axe du développement (I, II ou III du plan). Tout le reste du paragraphe ne sert qu'à la justifier, l'illustrer, la nuancer.
\end{definition}

\begin{methode}[Le schéma I-A-E-L : construire un paragraphe en quatre étapes]
\begin{enumerate}
    \item \textbf{I — Idée directrice} : Une phrase affirmative qui pose l'argument central du paragraphe (ex. : « La modernité prônée par Lakunle est une modernité de surface, déconnectée des réalités sociales »).
    \item \textbf{A — Argument} : L'explication logique. Pourquoi cette idée est-elle vraie ? On mobilise ici des concepts, des mécanismes, des causalités (ex. : « Il confond modernité technique et rupture culturelle, refusant la synthèse entre tradition et progrès »).
    \item \textbf{E — Exemple} : L'illustration concrète. Elle ancre l'abstrait dans le réel. Elle puise dans l'œuvre au programme (\emph{Le Lion et la perle}), l'histoire, l'actualité camerounaise, ou une observation personnelle validée (ex. : « Dans la scène de l'école, Lakunle fait réciter aux enfants des leçons sur la « machine » sans lien avec leur environnement agricole »).
    \item \textbf{L — Lien (ou mini-conclusion)} : La phrase qui referme le paragraphe en résumant ce qui a été prouvé \textbf{et} qui ouvre sur le paragraphe suivant (transition). Elle répond à la question : « Qu'est-ce que cela change pour la suite ? » (ex. : « Cette modernité stérile explique l'échec de sa séduction face à Sidi, mais elle masque aussi la ruse d'un Baroka bien plus lucide »).
\end{enumerate}
\end{methode}

\begin{popfigure}
\begin{tikzpicture}[
    block/.style={rectangle, draw=popdark, thick, rounded corners=6pt, minimum height=2.2cm, minimum width=3.8cm, align=center, font=\small, fill=popcream},
    arrow/.style={->, >=Stealth, thick, popblue, shorten >=2pt, shorten <=2pt},
    labelstyle/.style={font=\tiny\bfseries, text=popdark, anchor=south, yshift=-2pt}
]
% Nodes
\node[block, fill=popblueL, align=center] (I) at (0,0){\textbf{I — IDÉE DIRECTRICE}\\\textit{« La thèse du paragraphe »}\\\textit{Annonce l'argument}};
\node[block, fill=popgreenL, align=center] (A) at (5.5,0){\textbf{A — ARGUMENT}\\\textit{« Le raisonnement »}\\\textit{Explique, analyse, prouve}};
\node[block, fill=poporangeL, align=center] (E) at (11,0){\textbf{E — EXEMPLE}\\\textit{« La preuve concrète »}\\\textit{Œuvre, histoire, actualité, vécu}};
\node[block, fill=poppinkL, align=center] (L) at (16.5,0){\textbf{L — LIEN / TRANSITION}\\\textit{« Le bilan + l'annonce »}\\\textit{Referme et ouvre}};

% Arrows
\draw[arrow] (I.east) -- node[labelstyle] {justifie} (A.west);
\draw[arrow] (A.east) -- node[labelstyle] {illustre} (E.west);
\draw[arrow] (E.east) -- node[labelstyle] {conclut / relie} (L.west);

% Global label
\node[above=0.3cm of I.north, font=\bfseries\large, text=popdark] {STRUCTURE TYPE D'UN PARAGRAPHE ARGUMENTATIF (I-A-E-L)};
\end{tikzpicture}
\legende{Schéma de la structure I-A-E-L : chaque bloc coloré représente une fonction rhétorique précise. L'enchaînement est logique et non arbitraire.}
\end{popfigure}

\courssub{Les types de phrases au service de l'argumentation}

La dissertation n'est pas un roman. La syntaxe y est une arme rhétorique. Le programme de Première technique exige la maîtrise des \cle{types de phrases obligatoires} (déclarative, interrogative, impérative, exclamative) et des \cle{types de phrases facultatifs} (négative, passive, emphatique). Choisir le bon type au bon moment, c'est guider le lecteur, marquer l'évidence, provoquer l'adhésion ou souligner l'absurde.

\begin{popfigure}
\begin{center}
\begin{tabularx}{\linewidth}{|>{\bfseries}l|X|X|X|}
\hline
\rowcolor{popblueL}
\textbf{Type de phrase} & \textbf{Forme canonique} & \textbf{Exemple argumentatif (dissertation)} & \textbf{Effet de sens / Fonction rhétorique} \\
\hline
\textbf{Déclarative} & Sujet + Verbe + Complément. Point final. & « La tradition n'est pas un frein, mais un socle. » & \textbf{Affirmation, énonciation de la thèse.} Pose l'évidence, ancre le propos. Base de l'argumentation. \\
\hline
\textbf{Interrogative} & Inversion, « est-ce que », intonation. Point d'interrogation. & « Peut-on bâtir un avenir en reniant ses racines ? » & \textbf{Question rhétorique.} Implique le lecteur, oriente sa pensée, feint le dialogue pour mieux imposer la réponse (souvent évidente). \\
\hline
\textbf{Impérative} & Verbe à l'impératif (2e pers.). Point ou point d'exclamation. & « Observons le sort de Lakunle : il finit isolé. » & \textbf{Injonction, guide de lecture.} Ordonne au lecteur d'effectuer une opération mentale (regarder, comparer, conclure). Autorité du scripteur. \\
\hline
\textbf{Exclamative} & « Que / Quel + ... » ou intonation forte. Point d'exclamation. & « Quelle ironie : le « moderne » est le plus archaïque ! » & \textbf{Marque l'affect, l'indignation, l'admiration.} Relance l'attention, souligne un paradoxe, un scandale, une évidence choc. À doser (voir \textit{Le savais-tu}). \\
\hline
\textbf{Négative} & « Ne ... pas / plus / jamais / rien / personne... » & « Lakunle ne comprend \textbf{pas} que le progrès sans héritage est vide. » & \textbf{Refus, opposition, correction.} Déconstruit une idée reçue, marque la limite, définit par l'exclusion. \\
\hline
\textbf{Passive} & Sujet + être + participe passé + (par agent). & « La perle est \textbf{convoitée} par le Lion ; elle n'est pas \textbf{choisie} par la jeune fille. » & \textbf{Décentrement, objectivation.} Masque l'agent (si inconnu ou secondaire), met l'accent sur l'action subie ou le processus. Style plus formel, distancié. \\
\hline
\textbf{Emphatique} & C'est ... qui / que ; Ce n'est pas ... que ; Le fait que ... & « \textbf{C'est bien} l'intelligence de Baroka \textbf{qui} triomphe, \textbf{et non} sa force. » & \textbf{Mise en relief (focalisation).} Isole un élément (sujet, COD, circonstantiel) pour en faire le point crucial. Force l'attention sur l'essentiel. \\
\hline
\end{tabularx}
\end{center}
\legende{Tableau récapitulatif des types de phrases et de leur exploitation rhétorique en dissertation. Variez-les pour donner du rythme et de la puissance à votre développement.}
\end{popfigure}

\begin{attention}[Piège syntaxique : la phrase unique ou la parataxe sauvage]
Ne tombez pas dans deux écueils opposés :
\begin{itemize}
    \item \textbf{La phrase unique géante :} un paragraphe d'une seule phrase de 12 lignes, truffée de subordonnées. Le lecteur s'y perd, la logique se noie.
    \item \textbf{La parataxe (juxtaposition) :} « Argument 1. Argument 2. Exemple 1. Exemple 2. » sans connecteurs logiques (\emph{donc, car, cependant, par conséquent}). Cela fait une liste de courses, pas une démonstration.
\end{itemize}
\textbf{Règle d'or :} Une idée principale = une phrase principale (souvent déclarative ou emphatique). Les nuances, concessions, illustrations = subordonnées ou phrases courtes interrogatives/impératives pour relancer.
\end{attention}

\courssub{Choisir et intégrer des exemples pertinents : l'art de l'illustration}

L'exemple (le \textbf{E} de I-A-E-L) n'est pas une décoration. C'est la \cle{preuve par le concret}. Un argument sans exemple est une assertion gratuite ; un exemple sans argument est une anecdote stérile.

\begin{aretenir}[Les trois réservoirs d'exemples au Probatoire]
\begin{enumerate}
    \item \textbf{L'œuvre au programme :} \emph{Le Lion et la perle} (Wole Soyinka). C'est votre réservoir \textbf{prioritaire}. Citations courtes (intégrées dans la phrase), références précises de scènes (l'école, la place du marché, la chambre de Baroka, la danse du voyageur). \emph{Exemple :} « Le mime de la voiture (scène 1) montre Lakunle mimant le moteur avec son corps : la modernité est une performance ridicule. »
    \item \textbf{Les réalités camerounaises :} Coutumes (chefferies, dot, sociétés secrètes), bilinguisme, urbanisation (Yaoundé, Douala), agriculture, diaspora, actualité politique/sociale récente. \emph{Exemple :} « Le « ngondo » chez les Sawa illustre la vitalité d'une tradition qui s'adapte aux médias modernes. »
    \item \textbf{Histoire et actualité mondiale :} Décolonisation, figures panafricaines (Nkrumah, Mandela), révolutions technologiques, débats éthiques (IA, bioéthique), grands textes (Déclaration des droits de l'homme). \emph{Exemple :} « Le destin de Lumumba rappelle que la lucidité politique ne suffit pas sans ancrage populaire. »
\end{enumerate}
\end{aretenir}

\begin{attention}[L'exemple ne prouve pas à lui seul — La confusion « Exemple = Preuve »]
\textbf{Erreur majeure :} Empiler trois exemples en pensant avoir prouvé son point. \emph{« Lakunle est ridicule : il mime la voiture. Il refuse la dot. Il parle anglais. Donc c'est un faux moderne. »} $\rightarrow$ C'est une \textbf{énumération}, pas un argument.
\textbf{Correction :} L'argument (le \textbf{A}) établit le \textbf{mécanisme} (le \emph{pourquoi}). L'exemple (le \textbf{E}) \textbf{incarne} ce mécanisme.
\begin{quote}
\textit{Argument :} « Lakunle réduit la modernité à une mimique technique, ignorant sa dimension éthique et sociale. \\
\textit{Exemple :} « Lors de la leçon sur la « machine » (scène 1), il impose un vocabulaire industriel à des enfants paysans, créant un fossé entre le savoir scolaire et le vécu. »
\end{quote}
L'exemple \emph{illustre} le mécanisme exposé dans l'argument. Il le rend visible, tangible, incontestable.
\end{attention}

\courssub{Enchaîner les paragraphes : l'art de la transition (Bilan + Annonce)}

Vos paragraphes ne sont pas des îles. Ils forment un archipel relié par des ponts : les \cle{transitions}. Une transition réussie fait deux choses simultanément :
\begin{enumerate}
    \item \textbf{Le bilan (rétrospectif) :} Elle résume en une phrase ce que le paragraphe précédent a \textbf{établi} (la thèse locale validée).
    \item \textbf{L'annonce (prospectif) :} Elle introduit le \textbf{nouvel angle} du paragraphe qui suit, en montrant le lien logique (cause/conséquence, opposition, approfondissement, nuance).
\end{enumerate}

\begin{exemplebox}[Transitions types selon la logique du plan]
\textbf{Plan dialectique (Thèse / Antithèse / Synthèse)}
\begin{itemize}
    \item \textit{Fin Thèse (I) $\rightarrow$ Début Antithèse (II) :} « \textbf{Si} la tradition apparaît comme un frein au regard des critères occidentaux [Bilan I], \textbf{elle révèle pourtant}, sous l'autorité de Baroka, une capacité d'adaptation insoupçonnée [Annonce II]. »
    \item \textit{Fin Antithèse (II) $\rightarrow$ Début Synthèse (III) :} « \textbf{Toutefois}, l'opposition binaire tradition/modernité masque une réalité plus complexe : la synthèse opérée par Sidi [Bilan II + Annonce III]. »
\end{itemize}
\textbf{Plan thématique (Causes / Conséquences / Solutions)}
\begin{itemize}
    \item \textit{Fin Causes $\rightarrow$ Début Conséquences :} « \textbf{Ces racines profondes du malentendu} [Bilan Causes] \textbf{engendrent inévitablement} une crise de l'identité chez la jeunesse [Annonce Conséquences]. »
\end{itemize}
\textbf{Mots-charnières indispensables :} \emph{Ainsi, par conséquent, en revanche, de surcroît, plus fondamentalement, c'est pourquoi, dès lors, au total.}
\end{exemplebox}

\courssub{Enrichir son lexique : synonymie et antonymie au service de la nuance}

Répéter « moderne », « tradition », « Lakunle », « Baroka », « prouve », « montre » dix fois par page, c'est l'assurance d'une note médiocre. La \cle{synonymie} et l'\cle{antonymie} sont vos outils pour \textbf{varier, nuancer, opposer, préciser}.

\begin{propriete}[Synonymie : éviter la répétition, graduer l'intensité]
Ne cherchez pas le « synonyme parfait » (il n'existe pas), mais le \textbf{quasi-synonyme contextuel} qui apporte une nuance.
\begin{center}
\begin{tabularx}{\linewidth}{|l|X|X|}
\hline
\rowcolor{popblueL}
\textbf{Mot de base} & \textbf{Quasi-synonymes (nuances)} & \textbf{Contexte d'emploi en dissertation} \\
\hline
\textbf{Moderne} & \emph{Contemporain, novateur, progressiste, occidentalisé, « à la page »} & « Lakunle se veut \emph{progressiste} » (volonté) vs « Il est seulement \emph{occidentalisé} » (constat critique). \\
\hline
\textbf{Tradition} & \emph{Coutume, héritage, legs ancestral, patrimoine, mœurs, coutumier} & « Le \emph{patrimoine} culturel » (valeur positive, muséale) vs « Les \emph{mœurs} rigides » (connotation pesante). \\
\hline
\textbf{Prouver / Montrer} & \emph{Démontrer, établir, attester, illustrer, révéler, souligner, mettre en lumière} & « L'exemple \emph{illustre} » (image) ; « Le texte \emph{révèle} » (dévoile un sens caché) ; « Le raisonnement \emph{établit} » (rigueur logique). \\
\hline
\textbf{Penser / Croire} & \emph{Estimer, considérer, juger, postuler, supposer, présumer} & « Baroka \emph{postule} que le pouvoir se négocie » (hypothèse stratégique) vs « Lakunle \emph{juge} la dot barbare » (jugement de valeur). \\
\hline
\end{tabularx}
\end{center}
\end{propriete}

\begin{propriete}[Antonymie : structurer l'opposition, créer le conflit]
L'antonymie n'est pas seulement « le contraire de ». Elle permet de \textbf{construire des chiasmes, des parallélismes, des contrastes saisissants}.
\begin{itemize}
    \item \textbf{Antonymie complémentaire (exclusif) :} \textbf{Acheter / Vendre, Vivant / Mort, Présent / Absent.} $\rightarrow$ Pour des choix binaires radicaux.
    \item \textbf{Antonymie gradable (degré) :} \textbf{Chaud / Froid, Riche / Pauvre, Moderne / Traditionnel.} $\rightarrow$ Pour montrer un spectre, une évolution. « Ni totalement moderne, ni archaïque : \emph{hybride}. »
    \item \textbf{Antonymie inverse (relation) :} \textbf{Acheter / Vendre, Donner / Recevoir, Dominer / Soumettre.} $\rightarrow$ Pour analyser les rapports de force. « Baroka \textbf{donne} pour mieux \textbf{recevoir} (le pouvoir). »
\end{itemize}
\textbf{Astuce de style :} Utilisez l'antithèse syntaxique (parallélisme + antonymie) : « Lakunle \textbf{parle} le changement mais \textbf{fige} les esprits ; Baroka \textbf{semble} immobile mais \textbf{fait} bouger les lignes. »
\end{propriete}

\courssub{Lecture méthodique n°3 : \emph{Le Lion et la perle} — Modèle de confrontation d'idées}

Le programme prévoit une troisième lecture méthodique de la pièce de Wole Soyinka, centrée sur \textbf{le texte argumentatif}. La confrontation entre Lakunle (l'instituteur « moderne ») et Baroka (le Bale « traditionnel ») dans les scènes 1 et 2 est un \textbf{modèle de débat argumentatif} à étudier pour votre propre dissertation.

\begin{experience}[Analyse guidée : La joute oratoire Lakunle vs Baroka (Scène 1 — Extrait de la discussion sur la dot / Scène 2 — La séduction de Sidi)]
\textbf{Matériel :} Texte de la pièce (éd. CLE ou Hatier). \textbf{Protocole :}
\begin{enumerate}
    \item \textbf{Identifiez les thèses :} Lakunle = « La dot est une vente d'être humain, barbare. » Baroka = « La dot scelle l'alliance, honore la femme et sa famille. »
    \item \textbf{Repérez les arguments (A) :} Lakunle argue du droit universel, de l'égalité, de l'économie (gaspillage). Baroka argue de la coutume, de la solidarité lignagère, de la valeur symbolique.
    \item \textbf{Analysez les exemples (E) :} Lakunle cite les « pays civilisés », la Bible (parfois à contre-sens). Baroka cite les ancêtres, la nature (le lion, la perle), son propre pouvoir réel.
    \item \textbf{Observez les types de phrases :}
    \begin{itemize}
        \item Lakunle : \textbf{Déclaratives} dogmatiques (« C'est une honte »), \textbf{Interrogatives} rhétoriques agressives (« Est-ce que tu paies pour ta mère ? »), \textbf{Impératives} (« Refuse ! »).
        \item Baroka : \textbf{Déclaratives} imagées/proverbiales (« La femme est la perle... »), \textbf{Emphatiques} (« C'est \emph{moi} qui décide »), \textbf{Interrogatives} feignant l'étonnement (« Tu crois que je ne sais pas ? »).
    \end{itemize}
    \item \textbf{Évaluez l'efficacité :} Qui « gagne » le débat rhétorique ? Pourquoi ? (Baroka : ancrage concret, humour, adaptation, éthos d'autorité. Lakunle : abstraction, rigidité, éthos de donneur de leçons).
\end{enumerate}
\textbf{Interprétation :} La dissertation, c'est \textbf{jouer Baroka} : ancrer ses arguments dans du concret, varier ses types de phrases, respecter l'intelligence du lecteur (le jury), user d'images et de proverbes (culture générale) plutôt que de concepts creux.
\end{experience}

\begin{savaistu}[La phrase exclamative : l'épice, pas le plat principal]
La phrase exclamative (\emph{« Quelle injustice ! », « Comme il a raison ! »}) est un \textbf{accélérateur de rythme}. Bien placée (fin de paragraphe, révélation d'un paradoxe), elle « réveille » le correcteur. Mais \textbf{abusée} (une par paragraphe), elle devient \textbf{pathétique} (au sens péjoratif : théâtrale, fausse). Elle trahit l'émotion là où la dissertation exige la \textbf{raison maîtrisée}. Règle : \textbf{0 ou 1 par grande partie (I, II, III)}, jamais dans l'introduction ni la conclusion. Préférez l'emphase (\emph{C'est... qui}) pour la force sans l'excès.
\end{savaistu}

\begin{exemplebox}[Paragraphe type analysé : « Lakunle incarne une modernité mal comprise »]
\textbf{Paragraphe rédigé :}
\begin{quote}
\textbf{[I]} \textbf{Lakunle incarne une modernité de façade, purement mimétique, qui refuse la synthèse culturelle.} \textbf{[A]} En effet, il confond l'adoption de techniques occidentales avec l'accession à la modernité politique et anthropologique ; il singe les formes (langue, vêtement, mobilier) sans en comprendre les fondements éthiques — la démocratie, l'esprit critique, l'égalité réelle. \textbf{[E]} La scène de l'école (Acte I) en offre la caricature saisissante : face à des élèves paysans, il fait réciter une leçon sur « la machine à vapeur » et mime le piston de son propre corps, transformant le savoir en spectacle grotesque, déconnecté de la réalité agraire d'Ilujinle. \textbf{[L]} \textbf{C'est donc une modernité stérile qu'il propose,} qui ne séduit pas Sidi et laisse le champ libre à la ruse modernisatrice de Baroka.
\end{quote}
\textbf{Analyse I-A-E-L :}
\begin{itemize}
    \item \textbf{I (Idée) :} Phrase déclarative emphatique (\emph{incarne une modernité...}). Thèse locale claire.
    \item \textbf{A (Argument) :} Explication causale (\emph{En effet, il confond...}). Concepts : mimétisme, fondements éthiques. Phrase négative (\emph{sans en comprendre}).
    \item \textbf{E (Exemple) :} Scène précise (Acte I), description concrète (\emph{mime le piston}), qualification (\emph{caricature saisissante, grotesque, déconnecté}). Intégration fluide (pas de citation brute isolée).
    \item \textbf{L (Lien) :} Emphase (\emph{C'est donc... que}). Bilan (\emph{stérile}) + Annonce (\emph{laisse le champ libre à Baroka}). Transition parfaite vers le paragraphe sur Baroka.
\end{itemize}
\end{exemplebox}

\begin{exoresolu}[Exercice : Reconstruire et améliorer un paragraphe défaillant]
\textbf{Énoncé :} Voici un paragraphe d'élève sur le thème « La tradition est-elle un obstacle au progrès ? ». Il contient les idées mais la structure I-A-E-L est rompue, le lexique est pauvre, les types de phrases sont monotones (tout déclaratif). Réécrivez-le en appliquant la méthode.

\textit{Paragraphe brut :}
« La tradition n'est pas un obstacle. Baroka est moderne. Il a une voiture. Il veut l'électricité. Il utilise le téléphone. Il manipule Lakunle. La tradition aide le progrès. Sidi choisit Baroka. Donc la tradition n'empêche pas le progrès. »

\textbf{Solution détaillée :}
\begin{enumerate}
    \item \textbf{Analyse (Diagnostic) :} Idée directrice présente mais faible (« n'est pas un obstacle »). Pas d'argument explicatif (pas de \emph{pourquoi}). Exemples en vrac (liste : voiture, électricité, téléphone, manipulation). Pas de lien. Répétitions (progrès, tradition, Baroka). Syntaxe plate.
    \item \textbf{Reconstruction I-A-E-L :}
    \begin{itemize}
        \item \textbf{I :} \textbf{Loin d'être un frein, la tradition chez Baroka devient le levier d'une modernité assumée et souveraine.} (Emphase + antonymie frein/levier).
        \item \textbf{A :} \textbf{Car le Bale comprend que le progrès technique ne saurait s'imposer par la rupture, mais par l'appropriation :} il ne subit pas la modernité, il la \emph{domestique} au service de son pouvoir coutumier. (Argument causal + verbe fort « domestique »).
        \item \textbf{E :} \textbf{La scène de la chambre (Acte II) en est la preuve :} Baroka reçoit Sidi dans un décor où le lit « à ressort » (confort moderne) côtoie les masques ancestraux ; il passe un coup de fil à Lagos pour faire venir une machine à imprimer des timbres \emph{à son effigie}, transformant l'outil administratif en instrument de propagande personnelle. (Exemple unique, dense, intégré, montrant le mécanisme « domestication »).
        \item \textbf{L :} \textbf{Ainsi, la victoire de Baroka sur Lakunle scelle l'alliance du lion et de la perle :} la tradition, vivante, digère la modernité ; la modernité, morte, ne fait que la singer. (Bilan synthétique + image forte + annonce implicite sur la synthèse possible).
    \end{itemize}
    \item \textbf{Paragraphe final poli :}
    \begin{quote}
    \textbf{Loin d'être un frein, la tradition chez Baroka devient le levier d'une modernité assumée et souveraine.} \textbf{Car le Bale comprend que le progrès technique ne saurait s'imposer par la rupture, mais par l'appropriation :} il ne subit pas la modernité, il la \emph{domestique} au service de son pouvoir coutumier. \textbf{La scène de la chambre (Acte II) en est la preuve :} Baroka reçoit Sidi dans un décor où le lit « à ressort » côtoie les masques ancestraux ; il passe un coup de fil à Lagos pour faire venir une machine à imprimer des timbres \emph{à son effigie}, transformant l'outil administratif en instrument de propagande personnelle. \textbf{Ainsi, la victoire de Baroka sur Lakunle scelle l'alliance du lion et de la perle :} la tradition, vivante, digère la modernité ; la modernité, morte, ne fait que la singer.
    \end{quote}
    \item \textbf{Gains :} Structure I-A-E-L visible. Lexique varié (\emph{frein/levier, rupture/appropriation, domestique, digère/singe}). Types de phrases : Emphase (I), Déclarative explicative (A), Déclarative descriptive + Impérative implicite (E), Emphase antithétique (L). Cohérence logique impeccable.
\end{enumerate}
\end{exoresolu}

\begin{exercice}[À vous de jouer : Écrire un paragraphe I-A-E-L complet]
\textbf{Sujet :} « Dans \emph{Le Lion et la perle}, Sidi est-elle seulement une victime des hommes ? »
\textbf{Consigne :} Rédigez \textbf{un seul paragraphe} (environ 12-15 lignes) défendant l'idée que \textbf{« Sidi manœuvre et use de son pouvoir de séduction pour imposer son choix »} (Thèse : Sidi actrice, non victime).
\begin{enumerate}
    \item Écrivez l'\textbf{Idée directrice} (1 phrase forte, emphatique si possible).
    \item Développez l'\textbf{Argument} (2-3 phrases : explication psychologique / sociologique).
    \item Intégrez l'\textbf{Exemple} précis (scène du miroir, scène de la photo, ou refus final de Lakunle). Citez un détail textuel.
    \item Terminez par le \textbf{Lien} (Bilan + Ouverture sur le paradoxe : sa « victoire » la livre-t-elle vraiment à Baroka ?).
\end{enumerate}
\textbf{Contraintes :} Utilisez au moins \textbf{3 types de phrases différents} (dont 1 interrogative rhétorique ou 1 impérative). Mobilisez \textbf{2 paires synonymiques} et \textbf{1 antonymie} marquante.
\end{exercice}

\begin{corrige}[Exercice n°1 — Proposition de rédaction]
\textbf{[I]} \textbf{C'est en stratège lucide, bien plus qu'en proie passive, que Sidi conduit la danse des prétendants.} \textbf{[A]} \textbf{Car elle perçoit d'emblée la valeur marchande de son image :} sa beauté, révélée par le photographe, devient un capital symbolique qu'elle entend monnayer au plus haut prix, refusant la dot dérisoire de Lakunle comme l'humiliation d'un « prix de vente » sans enchères. \textbf{[E]} \textbf{La scène du miroir (Acte I) est révélatrice :} face à son portrait, elle ne s'admire pas seulement, elle \emph{calcule} — « Je suis la perle... Je veux un mari qui paie » — transformant la vanité en conscience politique de sa valeur. \textbf{[L]} \textbf{Toutefois, cette maîtrise apparente masque un piège :} en choisissant le Lion pour sa puissance, la Perle accepte-t-elle de devenir l'ornement d'une tradition qu'elle croyait dominer ? (Transition vers l'ambiguïté du dénouement).
\end{corrige}

\courssub{Synthèse : Votre checklist « Paragraphe Parfait » avant de rendre la copie}

\begin{center}
\begin{tabularx}{\linewidth}{|>{\bfseries}l|X|}
\hline
\rowcolor{popblueL}
\textbf{Critère} & \textbf{Question de contrôle} \\
\hline
\textbf{Structure I-A-E-L} & Ai-je une phrase d'accroche (I) ? Un argument explicatif (A) ? Un exemple précis et intégré (E) ? Une phrase de lien/bilan (L) ? \\
\hline
\textbf{Types de phrases} & Y a-t-il de la variété ? (Déclarative, Emphatique, Interrogative rhétorique, Impérative guide). Pas \emph{que} du déclaratif. \\
\hline
\textbf{Lexique (Syn/Ant)} & Ai-je évité les répétitions lourdes (synonymes contextuels) ? Ai-je utilisé une antithèse ou un chiasme pour marquer l'opposition ? \\
\hline
\textbf{Exemples} & Sont-ils tirés de \emph{Le Lion et la perle} (priorité), du Cameroun, de l'Histoire ? Sont-ils \textbf{analysés} (pas seulement cités) ? \\
\hline
\textbf{Transitions} & Le paragraphe précédent est-il résumé (bilan) ? Le suivant est-il annoncé (logique) ? \\
\hline
\textbf{Orthographe / Ponctuation} & Accords, homophones (\emph{a/à, et/est, ou/où}), ponctuation forte (deux-points, point-virgule) pour structurer la phrase longue. \\
\hline
\end{tabularx}
\end{center}

Vous possédez maintenant la « boîte à outils » complète pour bâtir le corps de votre dissertation : structure rigoureuse (I-A-E-L), syntaxe variée et efficace, lexique précis et nuancé, exemples choisis et exploités, transitions fluides. La prochaine section — \emph{Rédiger la conclusion} — vous apprendra à refermer l'édifice avec élégance et force.

\coursec{Rédiger la conclusion}

Après avoir maîtrisé la construction du paragraphe argumentatif — idée directrice, arguments, exemples et connecteurs logiques — nous abordons l'ultime étape de la dissertation : la conclusion. C'est la dernière impression que vous laissez au correcteur ; elle doit être nette, solide et mémorable. Une dissertation sans conclusion, ou avec une conclusion bâclée, donne l'impression d'un travail inachevé, quelle que soit la qualité du développement.

\courssub{La fonction de la conclusion : clore sans ouvrir de nouvelles brèches}

\begin{definition}[La conclusion]
La conclusion est la partie finale de la dissertation. Elle a pour unique fonction de \textbf{refermer hermétiquement la démonstration} menée dans le développement. Elle ne doit introduire \textbf{aucune idée nouvelle}, \textbf{aucun argument supplémentaire}, \textbf{aucune citation inédite}. Elle répond à la question : « Au terme de cette réflexion, que faut-il retenir exactement ? »
\end{definition}

\begin{aretenir}[Rôle triple de la conclusion]
La conclusion remplit trois fonctions indissociables :
\begin{enumerate}
    \item \textbf{Le bilan synthétique} : rappeler la logique du parcours (les grandes étapes du développement) sans répéter mot pour mot les idées directrices.
    \item \textbf{La réponse explicite} : formuler clairement la thèse défendue, en réponse directe à la problématique posée en introduction.
    \item \textbf{L'ouverture} : proposer une perspective de réflexion connexe, un élargissement mesuré, qui montre que le sujet dépasse le cadre strict de la dissertation.
\end{enumerate}
\end{aretenir}

\courssub{Les trois étapes : la méthode B–R–O (Bilan – Réponse – Ouverture)}

C'est la structure standard, exigible au Probatoire technique. Elle se visualise comme une pyramide inversée : on part du particulier (le résumé du devoir) pour aller vers le général (l'ouverture).

\begin{methode}[Le schéma B–R–O de la conclusion]
\textbf{Étape 1 — Bilan synthétique (1 à 2 phrases)}
\begin{itemize}
    \item Reprenez les \emph{idées directrices} de vos parties (I, II, III) en les fondant en une synthèse.
    \item Utilisez des connecteurs de synthèse : \emph{Au terme de cette analyse, force est de constater que… ; En définitive, cette étude montre que… ; Au total, il apparaît que…}
    \item \textbf{Interdit :} la liste scolaire (« J'ai dit que…, ensuite que…, enfin que… »).
\end{itemize}

\textbf{Étape 2 — Réponse explicite à la problématique (1 phrase ferme)}
\begin{itemize}
    \item Répondez \textbf{directement} à la question du sujet ou à votre problématique.
    \item Affirmez votre thèse avec assurance : \emph{La tradition n'est pas un obstacle au développement, mais son fondement vivant.} ou \emph{Le progrès technique améliore la condition humaine, à condition d'être encadré par l'éthique.}
    \item Évitez le « oui et non » frileux ; tranchez, nuancez, mais \textbf{concluez}.
\end{itemize}

\textbf{Étape 3 — Ouverture (1 phrase, 2 maximum)}
\begin{itemize}
    \item Posez une question connexe, suggérez un prolongement, élargissez à un domaine voisin (littérature, histoire, actualité, autre discipline).
    \item Formules : \emph{Cette réflexion invite à se demander si… ; Elle ouvre la voie à une étude sur… ; On pourrait alors s'interroger sur le rôle de…}
    \item \textbf{Attention :} l'ouverture n'est \textbf{pas} une nouvelle partie. Pas de développement, pas d'argument.
\end{itemize}
\end{methode}

\begin{popfigure}
\begin{center}
\begin{tikzpicture}[
    box/.style={draw=popdark, thick, rounded corners=6pt, text width=4.2cm, align=center, minimum height=1.8cm, inner sep=4pt, font=\small},
    arr/.style={-Stealth, very thick, popblue}
]
\node[box, fill=popblue!15, text width=5cm] (titre) at (0,0) {\textbf{CONCLUSION}\\\emph{Pyramide inversée :}\\\emph{du devoir vers l'horizon}};
\node[box, fill=popgreenL, align=center] (bilan) at (-5.2,-3.2){\textbf{1. BILAN SYNTHÉTIQUE}\\\emph{« Au terme de cette analyse, il apparaît que… »}\\Résumé fusionné des parties I, II, III};
\node[box, fill=poporangeL, align=center] (reponse) at (0,-3.2){\textbf{2. RÉPONSE EXPLICITE}\\\emph{« Ainsi, la tradition est le socle du développement. »}\\Thèse affirmée, réponse à la problématique};
\node[box, fill=poppinkL, align=center] (ouverture) at (5.2,-3.2){\textbf{3. OUVERTURE}\\\emph{« Reste à savoir comment l'institutionnaliser… »}\\Question connexe, élargissement mesuré};
\draw[arr] (titre.south) -- (bilan.north);
\draw[arr] (titre.south) -- (reponse.north);
\draw[arr] (titre.south) -- (ouverture.north);
\end{tikzpicture}
\end{center}
\legende{Schéma de la structure B–R–O : la conclusion resserre le propos (Bilan), l'affirme (Réponse), puis l'élargit (Ouverture).}
\end{popfigure}

\courssub{L'ouverture : art de la porte entrouverte, non de la porte grande ouverte}

\begin{attention}[L'ouverture n'est pas une nouvelle partie — une seule phrase suffit]
C'est l'erreur n°1 au Probatoire. Certains candidats rédigent un paragraphe entier pour l'ouverture, avec arguments et exemples. \textbf{Non.} L'ouverture est une \emph{invitation à réfléchir plus loin}, pas une démonstration.
\begin{itemize}
    \item \textbf{Bonne ouverture :} « Cette étude sur la tradition et le développement invite à s'interroger sur le rôle de l'école dans la transmission des valeurs ancestrales. » (1 phrase, domaine voisin : l'éducation).
    \item \textbf{Mauvaise ouverture :} « L'école est cruciale. En effet, elle forme les citoyens. Par exemple, au Japon, l'école enseigne le respect des aînés. Donc l'école transmet la tradition. » (\textbf{4 phrases, argument, exemple = nouvelle partie interdite}).
\end{itemize}
\end{attention}

\begin{definition}[Les trois types d'ouverture acceptés]
\begin{enumerate}
    \item \textbf{Question connexe} : prolonge le débat sur un aspect non traité. \emph{« Ne faudrait-il pas alors étudier l'impact des nouvelles technologies sur la transmission orale ? »}
    \item \textbf{Perspective nouvelle / élargissement disciplinaire} : lien avec l'histoire, la sociologie, la littérature, l'actualité. \emph{« Cette tension entre tradition et modernité éclaire sous un jour nouveau les débats actuels sur la place des coutumes dans la vie sociale camerounaise. »}
    \item \textbf{Citation d'autorité brève} (si pertinente et non utilisée dans le devoir) : \emph{« Comme l'écrivait Amadou Hampâté Bâ : "En Afrique, quand un vieillard meurt, c'est une bibliothèque qui brûle", ce qui rappelle l'urgence de transmettre la tradition. »}
\end{enumerate}
\end{definition}

\courssub{Formules utiles et formules bannies : le registre de la rigueur}

\begin{center}
\begin{tabularx}{\linewidth}{|l|X|X|}
\hline
\rowcolor{popblueL}
\textbf{Moment} & \textbf{\cle{Formules recommandées} (style soutenu, précis)} & \textbf{\cle{Formules bannies} (familières, vides, scolaires)} \\
\hline
\textbf{Bilan} & Au terme de cette analyse, il ressort que… ; En définitive, cette étude met en lumière… ; Force est de constater que… & Pour conclure, je vais dire que… ; En conclusion, j'ai montré que… ; Pour finir, on a vu que… \\
\hline
\textbf{Réponse} & Ainsi, [thèse affirmative] ; C'est donc à juste titre que l'on peut affirmer que… ; La réponse à notre problématique est sans ambiguïté : … & Donc oui, c'est vrai. ; En fait, ça dépend. ; Mon avis est que… ; Je pense que… \\
\hline
\textbf{Ouverture} & Cette réflexion invite à s'interroger sur… ; Elle ouvre la perspective de… ; On pourrait alors se demander si… ; Cet éclairage gagne à être croisé avec… & Voilà, c'est fini. ; Il y a encore beaucoup à dire. ; Pour aller plus loin, je propose d'étudier… (trop directif) \\
\hline
\end{tabularx}
\end{center}

\begin{aretenir}[Critères de réussite de la conclusion (grille Probatoire)]
\begin{itemize}
    \item \textbf{Cohérence :} la thèse affirmée correspond \textbf{exactement} à celle annoncée en introduction et défendue dans le développement.
    \item \textbf{Autonomie :} la conclusion se comprend seule ; le correcteur n'a pas besoin de relire le devoir pour saisir la réponse.
    \item \textbf{Concision :} 5 à 8 lignes maximum (environ 60 à 100 mots). Pas de paragraphe de 20 lignes.
    \item \textbf{Propreté :} aucune rature, écriture soignée, paragraphe unique (sauf si très long, alors deux paragraphes : Bilan + Réponse / Ouverture).
\end{itemize}
\end{aretenir}

\courssub{Cohérence d'ensemble : le test du « aller-retour » Introduction $\leftrightarrow$ Conclusion}

Avant de rendre votre copie, effectuez ce contrôle systématique :

\begin{experience}[Test de cohérence Introduction / Conclusion]
\textbf{Matériel} : votre introduction (brouillon ou copie) et votre conclusion rédigée.

\textbf{Protocole} :
\begin{enumerate}
    \item Surlignez dans l'introduction : \textbf{sujet amené}, \textbf{problématique}, \textbf{thèse annoncée}, \textbf{plan annoncé (I, II, III)}.
    \item Surlignez dans la conclusion : \textbf{bilan} (écho au plan), \textbf{réponse} (écho à la thèse et à la problématique), \textbf{ouverture}.
    \item Vérifiez les correspondances :
    \begin{itemize}
        \item la \textbf{réponse} = la \textbf{thèse annoncée} (même sens, vocabulaire proche) ;
        \item le \textbf{bilan} reprend les \textbf{idées directrices} du plan annoncé (dans le même ordre) ;
        \item l'\textbf{ouverture} ne contredit pas le sujet amené.
    \end{itemize}
\end{enumerate}

\textbf{Interprétation} : si une divergence apparaît (ex. thèse annoncée : « la tradition freine le développement » ; thèse conclue : « la tradition aide le développement »), \textbf{corrigez la conclusion} pour l'aligner sur l'introduction (ou l'introduction si le développement a dérivé, mais c'est plus long).
\end{experience}

\courssub{Exemple commenté : sujet « La tradition est-elle un obstacle au développement ? »}

Voici une conclusion rédigée selon la méthode B–R–O, suivie de son analyse.

\begin{exemplebox}[Conclusion rédigée — Sujet : « La tradition est-elle un obstacle au développement ? »]
\textbf{Problématique :} la tradition, en tant qu'héritage culturel, constitue-t-elle un frein à la modernisation ou le socle sans lequel tout développement est déraciné ?

\textbf{Thèse défendue :} la tradition n'est pas un obstacle, mais le fondement vivant d'un développement endogène et durable.

\textbf{Plan :} I. Les apparences du conflit (la tradition perçue comme frein). II. La tradition, réservoir de valeurs et de savoirs pour le développement. III. L'articulation nécessaire : modernité critique et tradition vivante.

\textbf{Conclusion :}
\begin{quote}
\textcolor{popgreen}{\textit{[Bilan]}} Au terme de cette analyse, il apparaît que l'opposition entre tradition et développement relève d'un malentendu : là où l'on voit un frein, se révèle en réalité un réservoir de valeurs sociales, de savoirs endogènes et de cohésion communautaire. \textcolor{poporange}{\textit{[Réponse]}} Ainsi, la tradition n'est pas un obstacle au développement : elle en est le socle identitaire et le garant de la durabilité, pourvu qu'elle accepte une critique interne lucide. \textcolor{poppink}{\textit{[Ouverture]}} Reste à définir les modalités concrètes d'un dialogue entre sagesse ancestrale et innovation technique dans nos sociétés africaines.
\end{quote}
\end{exemplebox}

\begin{exoresolu}[Analyse détaillée de la conclusion ci-dessus]
\textbf{Étape 1 — Bilan (vert) :} « Au terme de cette analyse, il apparaît que l'opposition entre tradition et développement relève d'un malentendu : là où l'on voit un frein, se révèle en réalité un réservoir de valeurs sociales, de savoirs endogènes et de cohésion communautaire. »
\begin{itemize}
    \item \textbf{Écho à la partie I} : « opposition… malentendu » / « frein perçu ».
    \item \textbf{Écho à la partie II} : « réservoir de valeurs, savoirs endogènes, cohésion ».
    \item \textbf{Écho à la partie III} : annoncée par la nuance qui suivra (« critique interne lucide »).
    \item \textbf{Qualité} : phrase unique, synthétique, vocabulaire précis (durabilité, endogène).
\end{itemize}

\textbf{Étape 2 — Réponse (orange) :} « Ainsi, la tradition n'est pas un obstacle au développement : elle en est le socle identitaire et le garant de la durabilité, pourvu qu'elle accepte une critique interne lucide. »
\begin{itemize}
    \item \textbf{Directe} : reprend les termes du sujet (« obstacle », « développement »).
    \item \textbf{Affirmative} : « n'est pas… elle en est le socle ».
    \item \textbf{Nuancée} : « pourvu qu'elle accepte une critique interne » (reprend la condition de la partie III).
    \item \textbf{Zéro} « je pense », « à mon avis ».
\end{itemize}

\textbf{Étape 3 — Ouverture (rose) :} « Reste à définir les modalités concrètes d'un dialogue entre sagesse ancestrale et innovation technique dans nos sociétés africaines. »
\begin{itemize}
    \item \textbf{Domaine voisin} : la vie sociale et politique (dépasse le cadre strictement philosophique).
    \item \textbf{Concrète} : « modalités concrètes », « dialogue ».
    \item \textbf{Une seule phrase}. Pas de développement.
\end{itemize}
\end{exoresolu}

\courssub{Entraînement : rédiger la conclusion à partir de vos introductions}

Reprenez les \textbf{trois introductions} que vous avez rédigées lors de la séance consacrée à l'introduction (ou les modèles fournis ci-dessous si vous n'avez pas vos brouillons). Pour chacune, appliquez la méthode B–R–O \textbf{sur brouillon d'abord}, puis rédigez la conclusion définitive au net.

\begin{exercice}[Entraînement guidé — 3 sujets types Probatoire]
\textbf{Consigne} : pour chaque sujet, l'introduction (résumée) est donnée. Rédigez la conclusion complète (Bilan + Réponse + Ouverture) en respectant la thèse annoncée.

\vspace{0.3cm}
\noindent\textbf{Sujet A :} « L'argent fait-il le bonheur ? »
\begin{itemize}
    \item \textit{Introduction (résumée) :} Problématique : l'argent est-il une condition nécessaire ou suffisante du bonheur ? Thèse : l'argent est une condition matérielle nécessaire mais insuffisante ; le bonheur véritable réside dans des biens non marchands (liens, sens, santé). Plan : I. L'argent, source de sécurité et de liberté (nécessaire). II. L'argent impuissant face au mal-être existentiel (insuffisant). III. Le bonheur : relations, accomplissement, don de soi (biens non marchands).
\end{itemize}

\vspace{0.2cm}
\noindent\textbf{Sujet B :} « La technique libère-t-elle l'homme ou l'asservit-elle ? »
\begin{itemize}
    \item \textit{Introduction (résumée) :} Problématique : la technique est-elle émancipation ou aliénation ? Thèse : la technique libère des contraintes naturelles mais crée de nouvelles servitudes ; la libération réelle dépend du projet politique et éthique qui l'encadre. Plan : I. La libération par la technique (santé, communication, pénibilité réduite). II. Les nouvelles servitudes (surveillance, dépendance, standardisation). III. La condition de la libération : appropriation citoyenne et régulation éthique.
\end{itemize}

\vspace{0.2cm}
\noindent\textbf{Sujet C :} « Faut-il tout dire ? »
\begin{itemize}
    \item \textit{Introduction (résumée) :} Problématique : la vérité totale est-elle un devoir moral ou une violence destructrice ? Thèse : la franchise absolue est une chimère nuisible ; l'éthique de la communication exige une vérité dosée, contextualisée, bienveillante. Plan : I. L'idéal de transparence (confiance, authenticité). II. Les dangers du « tout dire » (humiliation, rupture, indiscrétion). III. La vertu du silence et de la parole mesurée (prudence, bienveillance, secret professionnel).
\end{itemize}
\end{exercice}

\begin{corrige}[Exercice — Propositions de conclusions (modèles)]
\textbf{Sujet A — L'argent et le bonheur}
\begin{quote}
\textit{[Bilan]} Au terme de cette réflexion, il ressort que l'argent, s'il assure les conditions matérielles de l'existence — sécurité, autonomie, accès aux soins —, se révèle impuissant à combler le vide existentiel et à acheter l'amour, le sens ou la paix intérieure. \textit{[Réponse]} Ainsi, l'argent ne fait pas le bonheur : il en est un auxiliaire nécessaire mais insuffisant, le bonheur authentique relevant de l'ordre du don, de la relation et de l'accomplissement personnel. \textit{[Ouverture]} Cette distinction invite à repenser nos indicateurs de prospérité : la richesse d'un pays mesure-t-elle vraiment le bien-vivre de ses populations ?
\end{quote}

\textbf{Sujet B — Technique : libération ou asservissement ?}
\begin{quote}
\textit{[Bilan]} En définitive, l'étude montre que la technique opère un double mouvement : elle libère l'humanité des contraintes biologiques et géographiques, mais tisse simultanément un filet invisible de dépendances techniques, de surveillance et d'uniformisation culturelle. \textit{[Réponse]} La technique ne libère ni n'asservit par nature ; elle devient libératrice seulement lorsqu'un projet politique lucide et une éthique partagée en orientent les usages vers l'autonomie réelle des personnes. \textit{[Ouverture]} Reste à interroger la capacité de nos sociétés à encadrer les innovations numériques avant qu'elles ne redéfinissent, sans débat, les contours de notre humanité.
\end{quote}

\textbf{Sujet C — Faut-il tout dire ?}
\begin{quote}
\textit{[Bilan]} Force est de constater que l'idéal d'une transparence totale, s'il séduit par sa pureté morale, heurte la fragilité des êtres et la complexité du lien social : la vérité brute blesse, le secret protège, le silence apaise. \textit{[Réponse]} Il ne faut donc pas tout dire ; la vertu n'est pas dans l'exhaustivité verbale mais dans la \emph{parole juste} — vraie, nécessaire, bienveillante — qui respecte la dignité d'autrui et la pudeur de l'intime. \textit{[Ouverture]} Cette éthique de la retenue trouve une urgence nouvelle à l'ère du tout-numérique, où l'oubli devient impossible et l'exposition permanente.
\end{quote}
\end{corrige}

\courssub{Le saviez-vous ? : la sanction de l'absence de conclusion}

\begin{savaistu}[Une dissertation sans conclusion = devoir inachevé]
Au Probatoire comme au Baccalauréat technique, la grille d'évaluation réserve des points spécifiques à la conclusion (souvent 2 à 3 points sur 20, selon les harmonisations). \textbf{Une copie sans conclusion visible (saut de ligne, paragraphe final distinct) est considérée comme inachevée.} Même avec un développement excellent, l'absence de conclusion fait chuter la note, car la compétence « conclure une démonstration » n'est pas validée.
\begin{itemize}
    \item \textbf{Astuce de gestion du temps :} réservez \textbf{15 minutes minimum} sur la durée de l'épreuve pour rédiger la conclusion au net. Si le temps manque, sacrifiez la fin de la partie III (dernier exemple) mais \textbf{jamais la conclusion}. Une partie III légèrement courte mais une conclusion parfaite vaut mieux que l'inverse.
    \item \textbf{Sur brouillon :} écrivez votre conclusion \textbf{avant} de passer au net le développement. Cela vous force à vérifier la cohérence entre thèse et plan.
\end{itemize}
\end{savaistu}

\courssub{Synthèse visuelle : conclusion faible vs conclusion réussie}

\begin{center}
\begin{tabularx}{\linewidth}{|l|X|X|}
\hline
\rowcolor{popblueL}
\textbf{Critère} & \textbf{\textcolor{poporange}{Conclusion FAIBLE (sujet A)}} & \textbf{\textcolor{popgreen}{Conclusion RÉUSSIE (modèle corrigé)}} \\
\hline
\textbf{Bilan} & « Pour conclure, j'ai dit que l'argent aide pour manger et se loger. Ensuite j'ai dit qu'il n'achète pas l'amour. Enfin j'ai dit que le bonheur c'est la famille. » & « Au terme de cette réflexion, il ressort que l'argent, s'il assure les conditions matérielles de l'existence, se révèle impuissant à combler le vide existentiel et à acheter l'amour, le sens ou la paix intérieure. » \\
\hline
\textbf{Réponse} & « Donc l'argent ne fait pas le bonheur, mais il aide un peu. C'est mon avis. » & « Ainsi, l'argent ne fait pas le bonheur : il en est un auxiliaire nécessaire mais insuffisant, le bonheur authentique relevant de l'ordre du don, de la relation et de l'accomplissement personnel. » \\
\hline
\textbf{Ouverture} & « On pourrait aussi parler de la pauvreté en Afrique. C'est triste. Il faut aider les pauvres. » & « Cette distinction invite à repenser nos indicateurs de prospérité : la richesse d'un pays mesure-t-elle vraiment le bien-vivre de ses populations ? » \\
\hline
\textbf{Style} & Oral, répétitif, « je », énumération, jugements de valeur non argumentés. & Soutenu, synthétique, impersonnel, vocabulaire précis (auxiliaire, existentiel, indicateur), ouverture vers un domaine voisin. \\
\hline
\textbf{Appréciation} & Compétence « conclure » non validée. & Compétence « conclure » pleinement validée. \\
\hline
\end{tabularx}
\end{center}

\begin{aretenir}[Check-list finale avant de rendre la copie — Conclusion]
\begin{itemize}
    \item[$\square$] La conclusion est un \textbf{paragraphe unique} (ou deux au maximum), bien distinct du développement (saut de ligne).
    \item[$\square$] Elle suit l'ordre \textbf{Bilan $\rightarrow$ Réponse $\rightarrow$ Ouverture}.
    \item[$\square$] Le \textbf{bilan} fusionne les parties en 1 à 2 phrases (pas de liste).
    \item[$\square$] La \textbf{réponse} reprend les mots-clés du sujet ou de la problématique et affirme la thèse.
    \item[$\square$] L'\textbf{ouverture} est \textbf{une seule phrase}, sans nouveau développement.
    \item[$\square$] \textbf{Zéro formule bannie} (« Pour conclure je vais dire que… », « À mon avis… »).
    \item[$\square$] \textbf{Cohérence parfaite} avec l'introduction (thèse, plan).
    \item[$\square$] Orthographe, ponctuation et écriture soignées.
\end{itemize}
\end{aretenir}

La conclusion est la signature de votre pensée. Elle transforme un devoir scolaire en une prise de position intellectuelle assumée. Maîtrisez le B–R–O, et vous maîtriserez l'art de clore — ce qui, en dissertation comme dans la vie, est la marque de la maturité. La séance suivante portera sur la confrontation, l'évaluation et l'amélioration de vos productions complètes.

\coursec{Confrontation, évaluation et amélioration des productions}

Au terme de la section précédente, tu sais désormais rédiger une conclusion qui clôt proprement ta dissertation : rappel de la problématique, bilan de la démonstration, ouverture. Mais une dissertation n'est jamais terminée à la fin de la première écriture. Entre le brouillon et la copie rendue, il existe une étape décisive que les bons élèves ne négligent jamais : la \cle{relecture}, la \cle{confrontation} avec le regard d'autrui et l'\cle{amélioration} de la production. C'est cette dernière étape du travail que nous allons maintenant étudier, car elle fait partie intégrante du savoir-faire exigé au Probatoire et au Baccalauréat technique : rédiger \emph{et} justifier une démarche, une réponse ou une conclusion.

\courssub{Pourquoi confronter et améliorer sa production ?}

\begin{definition}[Confrontation et amélioration]
La \cle{confrontation} consiste à soumettre sa production écrite au regard d'un autre lecteur (camarade, professeur) afin de recueillir des remarques argumentées. L'\cle{amélioration} est la phase de réécriture qui suit : on corrige, on reformule, on réorganise, puis on \emph{justifie} ses choix de correction. Ces deux moments transforment un premier jet en une copie aboutie.
\end{definition}

L'intuition est simple : quand on écrit, on est « dans » son texte ; on sait ce qu'on voulait dire, donc on ne voit plus ce qui manque ou ce qui cloche. Un lecteur extérieur, lui, lit réellement ce qui est écrit. La confrontation exploite ce décalage. Elle n'est pas une sanction mais un outil de progrès : au Baccalauréat, le correcteur sera précisément ce lecteur extérieur qui ne connaît ni tes intentions ni ton brouillon.

\courssub{La grille d'évaluation : savoir comment une copie est jugée}

Pour évaluer utilement une dissertation (la sienne ou celle d'un camarade), il faut connaître les critères du correcteur. On retient quatre grandes familles de critères, chacune assortie d'un barème indicatif sur 20 points.

\begin{aretenir}[Les quatre critères d'évaluation d'une dissertation]
\begin{enumerate}
\item \textbf{Contenu et pertinence} : le sujet est-il compris ? La problématique est-elle dégagée ? Les arguments et exemples répondent-ils à la question posée (et non à côté) ?
\item \textbf{Cohérence du plan} : les parties s'enchaînent-elles logiquement ? Chaque paragraphe suit-il le schéma idée directrice -- argument -- exemple ? Les connecteurs assurent-ils la transition ?
\item \textbf{Qualité de la langue} : orthographe, grammaire (types de phrases, accords, conjugaison), ponctuation, richesse du vocabulaire (synonymies pour éviter les répétitions).
\item \textbf{Présentation} : lisibilité, aération (alinéas), respect de la mise en page d'une dissertation (pas de titres numérotés, introduction et conclusion distinctes).
\end{enumerate}
\end{aretenir}

\begin{popfigure}
\begin{center}
\begin{tikzpicture}
\draw[fill=popblueL] (0,0) rectangle (10.5,1);
\node[text=popdark,font=\bfseries] at (5.25,0.5) {Grille d'évaluation d'une dissertation (barème indicatif sur 20)};
\draw[fill=popcream] (0,-1) rectangle (7,0);
\node[anchor=west,font=\bfseries] at (0.2,-0.5) {Critère};
\draw[fill=popcream] (7,-1) rectangle (10.5,0);
\node[font=\bfseries] at (8.75,-0.5) {Points};
\draw[fill=white] (0,-2.6) rectangle (7,-1);
\node[anchor=west,text width=6.6cm] at (0.2,-1.8) {\textbf{1. Contenu et pertinence} : sujet compris, problématique, arguments et exemples adaptés};
\draw[fill=white] (7,-2.6) rectangle (10.5,-1);
\node at (8.75,-1.8) {8 points};
\draw[fill=popgreenL] (0,-4.2) rectangle (7,-2.6);
\node[anchor=west,text width=6.6cm] at (0.2,-3.4) {\textbf{2. Cohérence du plan} : progression logique, paragraphes structurés, connecteurs};
\draw[fill=popgreenL] (7,-4.2) rectangle (10.5,-2.6);
\node at (8.75,-3.4) {5 points};
\draw[fill=white] (0,-5.8) rectangle (7,-4.2);
\node[anchor=west,text width=6.6cm] at (0.2,-5) {\textbf{3. Qualité de la langue} : orthographe, grammaire, ponctuation, vocabulaire};
\draw[fill=white] (7,-5.8) rectangle (10.5,-4.2);
\node at (8.75,-5) {5 points};
\draw[fill=poporangeL] (0,-7.4) rectangle (7,-5.8);
\node[anchor=west,text width=6.6cm] at (0.2,-6.6) {\textbf{4. Présentation} : lisibilité, alinéas, mise en page};
\draw[fill=poporangeL] (7,-7.4) rectangle (10.5,-5.8);
\node at (8.75,-6.6) {2 points};
\end{tikzpicture}
\end{center}
\legende{Grille d'évaluation d'une dissertation avec barème indicatif : le contenu pèse le plus lourd (8 points), suivi de la cohérence et de la langue (5 points chacun), puis de la présentation (2 points).}
\end{popfigure}

\begin{attention}[La langue coûte plus cher qu'on ne croit]
Beaucoup d'élèves pensent que seules les idées comptent. Or la langue et la présentation représentent ensemble 7 points sur 20 : une copie riche en idées mais truffée de fautes peut perdre la mention. Inversement, une langue soignée valorise un contenu moyen. La relecture n'est donc pas un luxe : c'est un investissement rentable.
\end{attention}

\courssub{La relecture méthodique : les trois passes}

Relire ne consiste pas à parcourir sa copie une fois, distraitement. Une relecture efficace se fait en \emph{passes} successives, chacune ayant un objectif unique : on ne peut pas tout surveiller en même temps.

\begin{methode}[La relecture en trois passes]
\begin{enumerate}
\item \textbf{Première passe : le sens et le plan.} Relis uniquement pour vérifier la logique : l'introduction pose-t-elle bien la problématique ? Chaque paragraphe contient-il une idée directrice, un argument, un exemple ? La conclusion répond-elle à la question de départ ? Souligne les passages hors sujet ou les transitions manquantes.
\item \textbf{Deuxième passe : la langue.} Relis phrase par phrase, lentement : accords (sujet-verbe, adjectifs), conjugaisons, homophones (a/à, et/est, son/sont), ponctuation, connecteurs logiques variés (d'abord, ensuite, en effet, cependant, enfin). Traque les répétitions et remplace-les par des \cle{synonymes} (par exemple, alterner « important », « essentiel », « capital », « primordial »).
\item \textbf{Troisième passe : la présentation.} Vérifie la lisibilité : alinéas en début de paragraphe, marge respectée, ratures propres, pas de titres numérotés dans le corps de la dissertation.
\end{enumerate}
\end{methode}

\begin{exemplebox}[Chasser les répétitions grâce à la synonymie]
Phrase initiale : « L'éducation est \textbf{importante} car elle forme les citoyens ; elle est \textbf{importante} aussi pour le développement du pays. »\\
Version améliorée : « L'éducation est \textbf{essentielle} car elle forme les citoyens ; elle est également \textbf{capitale} pour le développement du pays. »\\
La répétition a disparu grâce à deux synonymes, et le connecteur « également » fluidifie l'enchaînement.
\end{exemplebox}

\courssub{La confrontation en binômes : la lecture croisée}

La confrontation s'organise en classe par binômes : chaque élève lit la copie de son camarade pendant que celui-ci lit la sienne. Cette \cle{lecture croisée} obéit à des règles précises pour rester constructive.

\begin{methode}[Conduire une lecture croisée en binôme]
\begin{enumerate}
\item \textbf{Lire en entier sans rien écrire}, une première fois, pour saisir l'impression générale.
\item \textbf{Annoter au crayon}, avec bienveillance : souligner ce qui réussit (une belle formule, un bon exemple) autant que ce qui échoue. On critique le texte, jamais la personne.
\item \textbf{Justifier chaque remarque} : interdiction d'écrire « c'est faux » ou « c'est mauvais ». On écrit : « cet exemple ne prouve pas ton idée directrice, car... », « cet accord est incorrect parce que le sujet est pluriel ».
\item \textbf{Échanger oralement} : chacun explique ses annotations à l'autre, qui peut se défendre en argumentant. C'est cet échange qui fait progresser les deux partenaires.
\end{enumerate}
\end{methode}

\begin{attention}[Corriger chez autrui ce qu'on ne voit pas chez soi]
On repère presque toujours mieux les fautes des autres que les siennes : notre œil, habitué à notre propre texte, « répare » automatiquement nos erreurs pendant la relecture. La lecture croisée entraîne donc l'œil critique : à force de repérer les fautes d'accord ou les hors-sujets chez ton camarade, tu finiras par les voir dans tes propres copies. Ne refuse jamais cet exercice en le jugeant inutile : c'est un entraînement de correcteur, donc de futur candidat lucide.
\end{attention}

\courssub{Justifier une correction : rédiger et argumenter sa démarche}

Le savoir-faire officiel exige que tu saches « rédiger et justifier une démarche, une réponse ou une conclusion ». Concrètement, après avoir corrigé ta copie, tu dois pouvoir expliquer \emph{pourquoi} telle modification améliore le texte. Cette justification mobilise les notions de l'unité : caractéristiques du texte argumentatif, types de phrases, synonymie.

\begin{exoresolu}[Annoter puis améliorer un extrait de copie]
\textbf{Version initiale d'un paragraphe} (sujet : « L'école est-elle utile à la réussite sociale ? ») :

« L'école est très utile. Elle est utile parce qu'elle donne des connaissances. Par exemple mon frère il a fait des études et il a trouvé du travail. Donc l'école c'est bien. »

\tcblower

\textbf{Annotations de la lecture croisée} :
\begin{itemize}
\item Répétition de « utile » (deux fois) et registre familier (« c'est bien », « mon frère il a... »).
\item L'exemple est trop personnel et trop vague : il prouve peu.
\item L'idée directrice n'est pas développée : on annonce « très utile » sans préciser \emph{en quoi}.
\end{itemize}

\textbf{Version améliorée} :

« L'école constitue d'abord un instrument de réussite sociale. En effet, en transmettant des connaissances et des qualifications, elle ouvre l'accès aux emplois stables et mieux rémunérés. Ainsi, au Cameroun, les diplômés de l'enseignement technique trouvent plus facilement un emploi qualifié que les jeunes sans formation. L'école apparaît donc comme un tremplin vers l'insertion professionnelle. »

\textbf{Justification des modifications} :
\begin{enumerate}
\item J'ai remplacé « très utile » par « instrument de réussite sociale » : l'expression est plus précise et reprend les termes du sujet, ce qui garantit la pertinence.
\item J'ai supprimé la répétition de « utile » en utilisant des formulations synonymiques (« ouvre l'accès », « tremplin »), conformément à la règle de variation du vocabulaire.
\item J'ai remplacé l'exemple familial par un exemple général et vérifiable (les diplômés de l'enseignement technique au Cameroun) : un argument ne convainc que si l'exemple a une portée générale.
\item J'ai structuré le paragraphe avec les connecteurs « En effet », « Ainsi », « donc », qui marquent clairement les étapes idée -- argument -- exemple -- mini-conclusion.
\item J'ai banni le registre familier (« c'est bien ») au profit d'un registre soutenu, exigé dans une dissertation.
\end{enumerate}
\end{exoresolu}

\courssub{La remédiation : construire sa fiche personnelle de progrès}

Après plusieurs dissertations, chaque élève découvre qu'il refait \emph{les mêmes} erreurs : tel accorde mal le participe passé, tel oublie les alinéas, tel autre sort du sujet dès la deuxième partie. La \cle{remédiation} consiste à transformer ce constat en outil de progrès.

\begin{methode}[Élaborer sa fiche personnelle de progrès]
\begin{enumerate}
\item \textbf{Recenser} : sur tes trois dernières copies, liste toutes les erreurs signalées par le professeur ou les binômes.
\item \textbf{Classer} : range-les dans les quatre catégories de la grille (contenu, cohérence, langue, présentation) et compte-les.
\item \textbf{Identifier} ta ou tes erreurs \emph{récurrentes} : celles qui reviennent au moins trois fois.
\item \textbf{Formuler une règle} pour chaque erreur récurrente, avec un exemple corrigé (ex. : « Je vérifie l'accord du verbe en retrouvant son sujet : \emph{les jeunes sans formation trouvent} »).
\item \textbf{Consulter la fiche} avant chaque nouvelle dissertation et pendant la relecture finale.
\end{enumerate}
\end{methode}

\begin{popfigure}
\begin{center}
\begin{tikzpicture}
\begin{axis}[
  width=0.85\linewidth, height=6.5cm,
  ybar, bar width=18pt,
  symbolic x coords={Orthographe,Accords,Ponctuation,Répétitions,Hors sujet,Présentation},
  xtick=data, x tick label style={rotate=20,anchor=east,font=\small},
  ylabel={Pourcentage des erreurs relevées (\%)},
  ymin=0, ymax=35,
  nodes near coords, nodes near coords style={font=\small},
  every axis plot/.append style={fill=popblue,draw=popdark}
]
\addplot coordinates {(Orthographe,30) (Accords,25) (Ponctuation,15) (Répétitions,12) (Hors sujet,10) (Présentation,8)};
\end{axis}
\end{tikzpicture}
\end{center}
\legende{Répartition des erreurs fréquentes relevées dans les copies de la classe (enquête fictive à titre d'exemple) : l'orthographe et les accords totalisent à eux seuls 55\,\% des erreurs, ce qui oriente la remédiation.}
\end{popfigure}

Ce diagramme montre comment une classe peut exploiter ses propres statistiques : si l'orthographe et les accords représentent plus de la moitié des erreurs, c'est là que la remédiation collective doit porter en priorité. Chaque élève peut dresser le même diagramme pour ses propres copies et visualiser ses points faibles.

\courssub{Le compte rendu : présenter ses améliorations à la classe}

Dernière étape du dispositif : le \cle{compte rendu} oral. Chaque binôme présente à la classe, en quelques minutes, les améliorations apportées à sa production.

\begin{methode}[Réussir son compte rendu oral]
\begin{enumerate}
\item Annoncer le sujet traité et la difficulté principale rencontrée.
\item Lire l'extrait initial, puis l'extrait amélioré.
\item Expliquer les modifications en les rattachant aux notions de l'unité (texte argumentatif, types de phrases, synonymie, connecteurs).
\item Conclure par ce que cet exercice t'a appris pour la prochaine dissertation.
\end{enumerate}
\end{methode}

Le compte rendu n'est pas une simple lecture : c'est un exercice d'\cle{expression orale} argumentée, qui prépare à l'épreuve orale du Baccalauréat technique. Parler clairement, justifier ses choix, répondre aux questions de la classe : autant de compétences travaillées ici.

\courssub{Séquence finale : rédiger en conditions d'examen}

Tout ce travail de confrontation et d'amélioration aboutit à la séquence finale : la rédaction complète d'une dissertation en \cle{temps limité}, dans les conditions du Probatoire ou du Baccalauréat technique. Le tableau ci-dessous propose une gestion du temps pour une épreuve de trois heures.

\begin{center}
\begin{tabularx}{\linewidth}{|l|X|l|}\hline
\rowcolor{popblueL} \textbf{Étape} & \textbf{Travail à réaliser} & \textbf{Durée} \\ \hline
Analyse du sujet & Définir les termes, dégager la problématique & 20 min \\ \hline
Plan détaillé & Rédiger au brouillon idées, arguments, exemples & 30 min \\ \hline
Rédaction & Introduction, développement, conclusion au propre & 1 h 45 \\ \hline
Relecture & Les trois passes : sens, langue, présentation & 15 min \\ \hline
\end{tabularx}
\end{center}

\begin{savaistu}[Les 15 dernières minutes qui rapportent des points]
Les correcteurs du Baccalauréat le répètent chaque année : de nombreuses copies perdent deux à trois points à cause de fautes qu'une simple relecture aurait éliminées — accords oubliés, mots manquants, phrases inachevées laissées en plan. Consacrer les quinze dernières minutes de l'épreuve à la relecture, même au prix d'une conclusion un peu plus courte, peut faire gagner plusieurs points. Un candidat qui rend une copie légèrement moins longue mais propre et relue obtient presque toujours une meilleure note qu'un candidat qui écrit jusqu'à la dernière seconde sans se relire. Réserve donc systématiquement ce temps de relecture : c'est le geste qui distingue le candidat méthodique du candidat pressé.
\end{savaistu}

\begin{exercice}[Entraînement à la lecture croisée]
Voici un extrait de copie (sujet : « Faut-il préserver les langues locales au Cameroun ? ») :

« Les langues locales sont importante. Elles transmette la culture des ancêtres. Par exemple les contes et les proverbes. Donc il faut les préserver. »

\begin{enumerate}
\item Relève toutes les erreurs de langue et de construction du paragraphe.
\item Réécris le paragraphe en respectant le schéma idée directrice -- argument -- exemple.
\item Justifie chacune de tes corrections en une phrase.
\end{enumerate}
\end{exercice}

\begin{corrige}[Exercice : entraînement à la lecture croisée]
\begin{enumerate}
\item Erreurs relevées : « importante » doit s'accorder au pluriel (« importantes ») ; « transmette » doit être conjugué à la troisième personne du pluriel (« transmettent ») ; « Par exemple les contes et les proverbes » est une phrase sans verbe, donc incorrecte ; l'exemple n'est pas développé ; la conclusion « Donc il faut les préserver » est abrupte.
\item Réécriture possible : « Les langues locales sont importantes car elles transmettent la culture des ancêtres. En effet, c'est par elles que circulent les contes, les proverbes et les chants qui font vivre la mémoire d'un peuple. Ainsi, au Cameroun, les veillées de contes en langue ewondo ou bamiléké initient les jeunes générations aux valeurs de leur communauté. Préserver les langues locales, c'est donc préserver une identité culturelle irremplaçable. »
\item Justifications : j'ai corrigé les deux accords (adjectif et verbe) ; j'ai transformé la phrase sans verbe en phrase complète intégrée à l'argumentation ; j'ai développé l'exemple pour lui donner une portée générale ; j'ai remplacé la conclusion abrupte par une formule de bilan qui reprend les termes du sujet.
\end{enumerate}
\end{corrige}

\begin{aretenir}[L'essentiel de la section]
\begin{itemize}
\item Une dissertation s'évalue selon quatre critères : contenu et pertinence, cohérence du plan, qualité de la langue, présentation.
\item La relecture se fait en trois passes : sens et plan, langue, présentation.
\item La lecture croisée en binôme entraîne l'œil critique : on apprend à voir chez soi ce qu'on repère chez autrui.
\item Toute correction doit être \emph{justifiée} : c'est un savoir-faire exigé à l'examen.
\item La fiche personnelle de progrès transforme les erreurs récurrentes en objectifs de remédiation.
\item En conditions d'examen, réserver les quinze dernières minutes à la relecture fait gagner des points.
\end{itemize}
\end{aretenir}

\newpage

\coursec{Activité d'intégration}

\coursec{Leçon 03 : Dissertation — Rédiger une dissertation}

\courssub{Objectifs de l'activité}
Mobiliser l'ensemble des savoir-faire de la leçon « Dissertation : rédiger une dissertation » dans une production complète, évaluée selon la grille officielle du Probatoire. L'élève doit être capable de :
\begin{itemize}
    \item Analyser un sujet de dissertation à double sens (thèse/antithèse) et formuler une problématique pertinente.
    \item Construire un plan dialectique détaillé (introduction, développement en trois parties avec sous-parties argumentées, conclusion) en intégrant des exemples tirés d'un dossier documentaire.
    \item Rédiger une introduction (amorce, sujet posé, sujet divisé, problématique, annonce du plan) et une conclusion (bilan, ouverture) conformes aux attendus.
    \item Rédiger des paragraphes argumentatifs structurés (idée directrice, arguments, exemples, connecteurs logiques).
    \item S'auto-évaluer et évaluer un pair à l'aide d'une grille critériée, justifier les corrections et produire une version améliorée.
    \item Participer à un compte rendu collectif identifiant les erreurs récurrentes pour une remédiation ciblée.
\end{itemize}

\courssub{Situation}
Dans le cadre de la préparation à l'épreuve de français du Probatoire (série technique), vous participez à un atelier d'écriture encadré. Le sujet proposé est : \textbf{« Le développement du Cameroun passe-t-il par l'abandon des traditions ? »}.

Ce sujet s'inscrit dans le prolongement de l'étude de \emph{Le Lion et la perle} de Wole Soyinka (texte au programme, réception de textes : Lectures méthodiques 1 à 3), qui met en scène le conflit entre modernité (Lakunle, l'instituteur) et tradition (Baroka, le Bale) dans le village nigérian d'Ilujinle. Le Cameroun contemporain vit des tensions analogues : urbanisation galopante, érosion des langues nationales, transformation des structures familiales, mais aussi résilience des chefferies, des rites d'initiation et des savoirs endogènes.

\medskip
\textbf{Dossier documentaire fourni (extraits) :}

\begin{enumerate}
    \item \textbf{Extrait de \emph{Le Lion et la perle} (Acte 1, scène de la classe) :} Lakunle : « Il faut que nous changions tout cela. / Nous devons bâtir des écoles, des hôpitaux, / Des routes, des ponts, des usines. / Nous devons transformer ce village / En une ville moderne, propre, saine. / [...] La coutume est une chaîne / Qui nous retient dans l'obscurité. » — Baroka (plus tard) : « Le vieux monde et le nouveau / Peuvent marcher main dans la main. / [...] Je ne suis pas un obstacle au progrès, / Je suis son gardien. »
    \item \textbf{Données statistiques (INS, Recensement 2005 et projections 2023) :} Taux d'urbanisation : 58\% (2023) contre 40\% (1987). Population de Yaoundé : ~4,5 millions ; Douala : ~4 millions. Croissance urbaine annuelle : ~4,5\%. Part des moins de 25 ans : 60\% de la population totale.
    \item \textbf{Témoignages de jeunes camerounais (enquête de terrain 2024) :}
    \begin{itemize}
        \item \textit{Armel, 22 ans, étudiant à Douala :} « Mes parents parlent bassa entre eux, mais m'ont élevé en français. Je ne connais pas les rites du \emph{Ngondo}. Pour réussir ici, il faut l'anglais, le français, le code, pas les ancêtres. »
        \item \textit{Mireille, 19 ans, Bafoussam :} « J'ai fait mon \emph{la'akam} (initiation bamiléké) l'an dernier. Ça m'a donné de la discipline, du respect. Je suis développeuse web. Les deux ne s'excluent pas. »
        \item \textit{Jean-Paul, 24 ans, Yaoundé :} « Les chefferies freinent l'aménagement du territoire. On ne peut pas construire une route parce qu'un chef dit que c'est le terrain des ancêtres. C'est ça le frein au développement. »
    \end{itemize}
    \item \textbf{Photographies (légendées) :}
    \begin{itemize}
        \item Photo A : Quartier d'affaires d'Akwa (Douala), tours de verre, embouteillages, panneaux publicitaires numériques.
        \item Photo B : Chefferie supérieure de Bandjoun, case à piliers sculptés, notables en tenues traditionnelles, cérémonie du \emph{tso} (masque).
        \item Photo C : Village du Nord (région de l'Extrême-Nord), cases en banco, greniers surélevés, femmes au puits, antenne relais 4G visible au loin.
    \end{itemize}
\end{enumerate}

\courssub{Consignes}
Le travail se déroule en six étapes chronométrées (séances de 50 min chacune ou regroupées en atelier de 3h).

\begin{enumerate}
    \item \textbf{Analyse du sujet et problématisation (30 min) :} 
    \begin{itemize}
        \item Identifiez les termes clés (« développement », « Cameroun », « abandon », « traditions ») et leurs acceptions (économique, social, culturel, identitaire).
        \item Dégagez la tension du sujet : opposition apparente vs complémentarité possible.
        \item Formulez \textbf{une problématique unique} sous forme de question ouverte, nuancée, qui guide tout le devoir.
    \end{itemize}
    \item \textbf{Construction du plan détaillé dialectique (45 min) :} 
    \begin{itemize}
        \item Structurez le développement en \textbf{trois grandes parties} (Thèse / Antithèse / Synthèse ou Dépassement).
        \item Pour chaque partie : rédigez l'\textbf{idée directrice} (phrase complète), listez \textbf{deux arguments} distincts, et associez à chaque argument \textbf{un exemple précis} tiré \textbf{obligatoirement} du dossier (Soyinka, statistiques, témoignages, photos).
        \item Vérifiez la cohérence logique (connecteurs : \emph{d'abord, ensuite, cependant, par conséquent, en définitive}).
    \end{itemize}
    \item \textbf{Rédaction de l'introduction et de la conclusion en classe (50 min) :} 
    \begin{itemize}
        \item Introduction : amorce (accroche contextuelle), sujet posé (reformulation), sujet divisé (thèse/antithèse), problématique, annonce du plan.
        \item Conclusion : bilan synthétique (réponse à la problématique), ouverture (perspective, autre sujet voisin, citation).
    \end{itemize}
    \item \textbf{Rédaction du développement (1h) :} 
    \begin{itemize}
        \item Rédigez les trois parties en paragraphes distincts (un paragraphe par sous-partie/argument).
        \item Respectez la structure du paragraphe argumentatif : \textbf{Idée directrice} (en tête) $\rightarrow$ \textbf{Explication/Argumentation} $\rightarrow$ \textbf{Exemple/Illustration} (dossier) $\rightarrow$ \textbf{Mini-bilan/Transition}.
        \item Soignez la cohésion (connecteurs, champs lexicaux, référents).
    \end{itemize}
    \item \textbf{Confrontation binaire et amélioration (45 min) :} 
    \begin{itemize}
        \item Échangez votre copie avec un binôme. Évaluez-la avec la \textbf{grille officielle simplifiée} (voir Ressources).
        \item Annotez la copie : points forts (✓), points faibles (✗), suggestions précises (marge).
        \item Rédigez une \textbf{version améliorée} (corrections intégrées) sur une nouvelle feuille.
    \end{itemize}
    \item \textbf{Compte rendu collectif et remédiation (30 min) :} 
    \begin{itemize}
        \item Le professeur regroupe les erreurs fréquentes relevées par les binômes (ex. : hors-sujet, plan binaire au lieu de dialectique, exemples vagues, connecteurs absents, conclusion répétitive).
        \item Élaboration collective d'une \textbf{fiche de remédiation} « Mes 5 points de vigilance pour le Probatoire ».
    \end{itemize}
\end{enumerate}

\courssub{Ressources à mobiliser}
\begin{definition}[Structure type de l'introduction (méthode en 5 temps)]
\textbf{1. Amorce :} Accroche large (citation, fait d'actualité, référence littéraire) $\rightarrow$ \textbf{2. Sujet posé :} Reformulation claire du sujet $\rightarrow$ \textbf{3. Sujet divisé :} Présentation des deux thèses opposées (Pour/Contre) $\rightarrow$ \textbf{4. Problématique :} Question unique, issue de la tension, qui oriente le devoir $\rightarrow$ \textbf{5. Annonce du plan :} « Nous verrons d'abord que... ; ensuite que... ; enfin que... »
\end{definition}

\begin{definition}[Structure type du paragraphe argumentatif (modèle IAE)]
\textbf{I -- Idée directrice} (phrase nominale ou verbale, thèse de la sous-partie). \\
\textbf{A -- Argument(s)} (explication, cause, conséquence, preuve logique). \\
\textbf{E -- Exemple(s)} (fait précis, chiffre, citation, référence textuelle du dossier). \\
\textbf{(Optionnel) Mini-bilan / Transition} vers l'argument ou la partie suivante.
\end{definition}

\begin{definition}[Structure type de la conclusion (méthode en 2 temps + ouverture)]
\textbf{1. Bilan / Réponse :} Reprise synthétique du parcours (thèse $\rightarrow$ antithèse $\rightarrow$ synthèse) et réponse directe à la problématique. \\
\textbf{2. Ouverture :} Élargissement (autre pays, autre époque, autre genre littéraire, perspective future) ou citation pertinente. \textbf{Interdiction :} nouvelle idée non traitée, simple répétition de l'introduction.
\end{definition}

\begin{propriete}[Grille d'évaluation simplifiée (inspirée grille officielle Probatoire) -- 20 points]
\begin{center}
\begin{tabularx}{\linewidth}{|l|X|c|}\hline
\textbf{Critère} & \textbf{Descriptif attendus} & \textbf{Points} \\ \hline
Analyse \& Plan & Sujet analysé, problématique pertinente, plan dialectique clair, détaillé, équilibré. & 4 \\ \hline
Introduction & 5 temps respectés, accroche contextuelle, annonce de plan fidèle au développement. & 3 \\ \hline
Développement & 3 parties logiques, paragraphes IAE structurés, exemples du dossier exploités, connecteurs. & 8 \\ \hline
Conclusion & Bilan synthétique, réponse à la problématique, ouverture pertinente. & 3 \\ \hline
Langue & Correction syntaxique, orthographique, lexicale ; niveau de langue soutenu ; connecteurs variés. & 2 \\ \hline
\end{tabularx}
\end{center}
\end{propriete}

\courssub{Démarche et corrigé commenté}
\begin{exoresolu}{Dissertation complète : « Le développement du Cameroun passe-t-il par l'abandon des traditions ? »}
\textbf{Énoncé :} Sujet : « Le développement du Cameroun passe-t-il par l'abandon des traditions ? » — Dossier : extraits de \emph{Le Lion et la perle} (Soyinka), statistiques INS 2023, témoignages jeunes 2024, photographies (Akwa, Bandjoun, Extrême-Nord). Rédiger une dissertation dialectique structurée (Introduction, Développement en 3 parties, Conclusion) en exploitant obligatoirement le dossier.
\tcblower
\textbf{ÉTAPE 1 -- Analyse du sujet et Problématique}

\begin{itemize}
    \item \textbf{Mots-clés :} 
    \begin{itemize}
        \item \cle{Développement} : processus multidimensionnel (économique : infrastructures, PIB ; social : santé, éducation ; culturel : épanouissement identitaire ; durable : environnement, équité).
        \item \cle{Cameroun} : pays « Afrique en miniature », 250+ ethnies, bilinguisme officiel, urbanisation rapide, chefferies traditionnelles reconnues par la Constitution (art. 37).
        \item \cle{Abandon} : rejet total, rupture radicale, perte de repères. Terme fort, binaire.
        \item \cle{Traditions} : coutumes, rites, langues, organisation sociale (chefferies, familles élargies), savoirs endogènes (médecine, agriculture), valeurs (respect aînés, solidarité).
    \end{itemize}
    \item \textbf{Tension :} Le sujet suggère une alternative exclusive : modernité \textbf{ou} tradition. Or, la réalité (Soyinka, statistiques, témoignages) montre des hybridations.
    \item \textbf{Problématique choisie :} \textbf{« Dans quelle mesure le développement durable du Cameroun exige-t-il non pas l'abandon, mais la réinvention créatrice de ses traditions face aux défis de la modernité ? »}
\end{itemize}

\textbf{ÉTAPE 2 -- Plan détaillé dialectique}

\begin{center}
\begin{tabularx}{\linewidth}{|l|X|}\hline
\textbf{Partie} & \textbf{Contenu détaillé} \\ \hline
\textbf{INTRODUCTION} & Amorce : citation Soyinka / Photo C (antenne 4G au village). Sujet posé/divisé. Problématique. Annonce plan. \\ \hline
\textbf{I. Thèse : L'illusion d'un développement par rupture (l'argument moderniste)} & \textbf{Idée directrice :} Une certaine vision du progrès, portée par l'État et les élites occidentalisées, présente les traditions comme des freins archaïques à l'efficacité économique et à l'urbanisation. \\
& \textbf{A1. L'urbanisation impose une rupture avec les modes de vie ruraux.} Exemple : Stats INS (58\% urbains, croissance 4,5\%/an) + Photo A (Akwa, tours, embouteillages). Argument : concentration des services, salariat, anonymat vs communauté villageoise. \\
& \textbf{A2. Les traditions perçues comme obstacles techniques et juridiques.} Exemple : Témoignage Jean-Paul (chefferies/aménagement routes) + Lakunle (\emph{Le Lion et la perle} : « La coutume est une chaîne »). Argument : foncier coutumier, pesanteurs hiérarchiques, inégalités genre (héritage). \\ \hline
\textbf{II. Antithèse : L'impossibilité et le danger d'un développement sans racines (l'argument identitaire)} & \textbf{Idée directrice :} L'abandon des traditions provoque une crise identitaire, une perte de cohésion sociale et prive le développement de ressources culturelles et humaines endogènes. \\
& \textbf{A1. La tradition comme capital social et régulateur éthique.} Exemple : Témoignage Mireille (\emph{la'akam}, discipline, respect) + Baroka (\emph{Le Lion et la perle} : « Le vieux monde et le nouveau / Peuvent marcher main dans la main »). Argument : éducation civique, résolution conflits, solidarité (tontines). \\
& \textbf{A2. Les savoirs endogènes, leviers d'innovation locale.} Exemple : Photo B (Chefferie Bandjoun, architecture, organisation) + Photo C (greniers banco + antenne 4G = adaptation). Argument : architecture bioclimatique, agroécologie, pharmacopée, gouvernance participative (conseil des notables). \\ \hline
\textbf{III. Synthèse : Pour un développement endogène par réinvention tradition-modernité} & \textbf{Idée directrice :} Le développement durable au Cameroun réside dans une \cle{hybridation féconde} : modernité technique encadrée par éthique traditionnelle, et tradition revitalisée par la science et le droit. \\
& \textbf{A1. Institutionnaliser le dialogue : cadre juridique et éducatif.} Exemple : Constitution art. 37 (chefferies) + École bilingue/biculturelle (enseignement langues nationales + STEM). Argument : légitimité de l'État, citoyenneté ancrée. \\
& \textbf{A2. Exemples concrets de « tradition réinventée » au Cameroun.} Exemple : \emph{Ngondo} (fête Sawa) devenu festival culturel international + start-up \emph{GiftedMom} (e-santé maternelle en langues locales) + architecture d'inspiration vernaculaire (ex. Francis Kéré, référence panafricaine ; réhabilitation de l'architecture en terre au Cameroun). Argument : économie créative, soft power, résilience climatique. \\ \hline
\textbf{CONCLUSION} & Bilan : Rupture = déracinement/échec ; Conservation figée = musée. Synthèse = \textbf{réinvention}. Réponse : Non, l'abandon n'est pas la voie ; la \textbf{transformation critique} l'est. Ouverture : « Et si le modèle camerounais du XXIe siècle était l'\emph{ubuntu} numérique ? » (Réf. philosophie bantoue + tech). \\ \hline
\end{tabularx}
\end{center}

\textbf{ÉTAPE 3 -- Introduction (rédaction type)}

« "Le vieux monde et le nouveau peuvent marcher main dans la main", affirme le Bale Baroka dans \emph{Le Lion et la perle} de Wole Soyinka. Cette réplique, loin de nier la modernité, en redéfinit les termes : non pas table rase, mais continuité critique. Au Cameroun, pays "Afrique en miniature" où 58\% de la population vit désormais en ville (INS, 2023) et où les tours de verre d'Akwa (Douala) côtoient les cases à piliers de Bandjoun, la question du développement hante les politiques publiques. \textbf{Le développement du Cameroun passe-t-il par l'abandon des traditions ?} Le terme "abandon" suggère une rupture radicale, une table rase identitaire. Pourtant, les traditions — langues, chefferies, rites, solidarités — structurent encore le quotidien de millions de Camerounais. \textbf{Dans quelle mesure le développement durable du Cameroun exige-t-il non pas l'abandon, mais la réinvention créatrice de ses traditions face aux défis de la modernité ?} Nous analyserons d'abord la tentation moderniste d'une rupture salutaire (I), avant de montrer le danger d'un déracinement identitaire et la valeur de l'héritage (II), pour enfin proposer la voie d'une hybridation féconde, seul gage de développement endogène (III). »

\textbf{Commentaire :} Amorce littéraire (Soyinka) + chiffre clé (INS) + image (Photo A/B). Sujet posé/divisé clair. Problématique nuancée (réinvention vs abandon). Annonce plan fidèle aux 3 parties.

\textbf{ÉTAPE 4 -- Extraits du Développement (Partie I et III, paragraphes types IAE)}

\medskip
\textbf{PARTIE I -- Sous-partie A1 (Paragraphe type)}

\textbf{[Idée directrice]} L'urbanisation galopante impose une rupture anthropologique avec les modes de vie traditionnels ruraux, perçus comme incompatibles avec l'efficacité économique moderne. \\
\textbf{[Argument]} Le passage d'une société de subsistance communautaire à une économie de marché salariale concentrée dans les métropoles (Yaoundé, Douala) nécessite une mobilité géographique, une disponibilité temporelle et une individualisation que la structure villageoise, fondée sur la lignée et la terre ancestrale, peine à offrir. L'anonymat citadin, la monétarisation des relations et la rapidité des flux d'information créent un décalage structurel avec le temps long de la tradition orale et du consensus des anciens. \\
\textbf{[Exemple]} Les projections de l'INS (2023) indiquent un taux d'urbanisation de 58\% avec une croissance annuelle de 4,5\%, soit l'une des plus rapides d'Afrique. La photographie A du dossier illustre ce basculement : le quartier d'affaires d'Akwa à Douala, avec ses tours de verre, ses embouteillages permanents et ses panneaux publicitaires numériques, incarne physiquement cette modernité qui ne laisse guère de place à la case du chef ou à l'arbre à palabres. Lakunle, dans \emph{Le Lion et la perle}, résume cette vision : « La coutume est une chaîne qui nous retient dans l'obscurité. Il faut bâtir des usines, des routes, des hôpitaux. » \\
\textbf{[Transition]} Cependant, cette vision linéaire du progrès occulte les coûts sociaux d'une modernité sans âme.

\medskip
\textbf{PARTIE III -- Sous-partie A2 (Paragraphe type)}

\textbf{[Idée directrice]} Des initiatives concrètes au Cameroun prouvent déjà la fécondité d'une tradition réinventée par les technologies et l'entrepreneuriat culturel. \\
\textbf{[Argument]} Loin du folklore muséal, des acteurs culturels et techniciens transforment le patrimoine en leviers économiques et sociaux. La fête du \emph{Ngondo} (peuple Sawa), jadis rituel initiatique secret, est devenue un festival international attirant touristes et médias, générant des retombées économiques tout en revitalisant la langue duala. Dans le secteur de la santé, des start-up comme \emph{GiftedMom} utilisent le mobile (très pénétrant même en zone rurale, cf. Photo C : antenne 4G) pour diffuser des conseils prénataux en langues nationales, alliant savoir médical moderne et accessibilité culturelle. L'architecture contemporaine d'inspiration vernaculaire (ex. travaux de Francis Kéré au Burkina, référence panafricaine ; réhabilitation de l'architecture en terre au Cameroun) réemploie la terre cuite et la ventilation naturelle des cases traditionnelles (Photo B, greniers surélevés) pour construire des écoles et cliniques bioclimatiques à faible coût énergétique. \\
\textbf{[Mini-bilan]} Ces exemples montrent que la tradition n'est pas un bloc figé, mais un réservoir de formes et de sens que la modernité technique peut actualiser.

\textbf{ÉTAPE 5 -- Conclusion (rédaction type)}

« Au terme de cette analyse, il apparaît que l'alternative "développement \textbf{ou} traditions" est un faux dilemme. La thèse moderniste (I) a raison de souligner l'inadéquation de certaines pratiques coutumières (foncier, inégalités) avec les exigences de l'État de droit et de l'économie de marché ; mais l'antithèse (II) démontre qu'un développement sans ancrage culturel produit des cités-dortoirs sans cohésion, des individus déracinés et une perte de biodiversité culturelle irréversible. La synthèse (III) ouvre la voie réelle : le développement camerounais ne passe pas par l'\textbf{abandon}, mais par la \textbf{réinvention critique}. Il s'agit de soumettre les traditions au crible des droits humains (égalité femme/homme, protection enfance) et de les outiller par la science (numérique, bioclimatique, agronomie) pour en faire des moteurs d'innovation endogène. Le Cameroun a-t-il les moyens de refuser ce pari de l'hybridation ? L'histoire du continent, de la philosophie bantoue de l'\emph{ubuntu} aux \emph{hubs} technologiques de Nairobi ou Lagos, suggère que l'avenir appartient à ceux qui savent "faire du neuf avec du vieux", pour reprendre la formule de Gaston Bachelard. Le développement durable du Cameroun sera \emph{bamiléké, sawa, béti, fulbé... et numérique}, ou ne sera pas. »

\textbf{Commentaire :} Bilan structuré (I $\rightarrow$ II $\rightarrow$ III). Réponse nette à la problématique (réinvention, non abandon). Ouverture philosophique (Bachelard) et prospective (modèle africain), sans idée nouvelle non traitée.

\textbf{ÉTAPE 6 -- Simulation confrontation binaire (extrait grille)}

\begin{center}
\begin{tabularx}{\linewidth}{|l|X|X|}\hline
\textbf{Critère} & \textbf{Évaluation Binôme (Exemple d'erreur fréquente)} & \textbf{Correction apportée (Version améliorée)} \\ \hline
\textbf{Problématique} & « Le développement passe-t-il par les traditions ? » (Trop simple, binaire, ne reprend pas "abandon"). & Reformulation intégrant "abandon" vs "réinvention" + "durable". \\ \hline
\textbf{Plan} & Plan thématique : 1. Économie 2. Culture 3. Politique. (Pas dialectique, pas de thèse/antithèse). & Plan dialectique : I. Rupture moderniste (Thèse) / II. Danger déracinement (Antithèse) / III. Hybridation (Synthèse). \\ \hline
\textbf{Exemples} & « Les traditions aident au développement » (Vague, pas de doc). & Citation précise : « Témoignage Mireille : \emph{la'akam} donne discipline pour le code » + « Photo C : antenne 4G + greniers banco ». \\ \hline
\textbf{Connecteurs} & « Et aussi, de plus, ensuite » (Additifs seulement). & « Certes... cependant... par conséquent... en définitive » (Concession, opposition, conséquence, conclusion). \\ \hline
\textbf{Conclusion} & Répétition mot à mot de l'introduction. & Bilan synthétique + Ouverture sur \emph{ubuntu} numérique. \\ \hline
\end{tabularx}
\end{center}

\textbf{ÉTAPE 7 -- Bilan collectif et Fiche de remédiation (extrait)}

\begin{aretenir}[Fiche de remédiation collective -- « Mes 5 points de vigilance Probatoire »]
\begin{enumerate}
    \item \textbf{Problématique :} Ne pas recopier le sujet. La formuler en question \textbf{ouverte}, issue de la \textbf{tension} (mots : "dans quelle mesure", "comment", "pourquoi"). Bannir les questions fermées (Oui/Non).
    \item \textbf{Plan dialectique :} 3 parties \textbf{obligatoires} (Thèse / Antithèse / Synthèse-Dépassement). 2 arguments + 1 exemple \textbf{par sous-partie}. Vérifier que l'annonce de plan (intro) \textbf{correspond exactement} aux idées directrices du développement.
    \item \textbf{Paragraphe IAE :} Une idée directrice \textbf{par paragraphe} (souvent en première phrase). L'exemple \textbf{doit prouver l'argument}, pas juste l'illustrer vaguement. Citer la source (Soyinka, INS, Témoignage X, Photo Y).
    \item \textbf{Connecteurs logiques :} Varier : \emph{Certes / Il est vrai que} (concession) ; \emph{Cependant / Toutefois / En revanche} (opposition) ; \emph{Par conséquent / Dès lors / C'est pourquoi} (conséquence) ; \emph{En outre / De surcroît} (addition argumentative).
    \item \textbf{Conclusion :} \textbf{Zéro répétition} de l'intro. Bilan = réponse à la problématique. Ouverture = \textbf{nouvel horizon} (autre pays, autre époque, citation, perspective future), pas nouvelle thèse.
\end{enumerate}
\end{aretenir}

\end{exoresolu}

\courssub{Bilan de l'activité}
Cette activité d'intégration a permis de :
\begin{itemize}
    \item \textbf{Valider la progression spiralaire :} Les élèves ont réinvesti séquentiellement l'analyse de sujet (séance 1), le plan (séance 2), l'introduction/conclusion (séance 3), le paragraphe IAE (séance 4), dans une production unique chronométrée (conditions d'examen).
    \item \textbf{Ancrer les apprentissages dans le réel camerounais :} Le dossier (Soyinka + INS + Témoignages + Photos) a ancré l'argumentation dans des réalités vécues (urbanisation, chefferies, numérique, identité), évitant les généralités abstraites sanctionnées au Probatoire.
    \item \textbf{Développer l'esprit critique et l'autonomie :} La confrontation binaire (évaluation par les pairs) a forcé l'explicitation des critères de réussite (grille) et la métacognition (justification des corrections). La version améliorée prouve la progression intra-séance.
    \item \textbf{Cibler la remédiation :} Le compte rendu collectif a produit une fiche « 5 points de vigilance » personnalisée par la classe, outil concret de révision pour l'examen.
    \item \textbf{Repérer les zones de fragilité persistantes :} La formulation de la problématique (trop souvent binaire ou hors-sujet), l'équilibre des parties (synthèse souvent faible), la citation précise des exemples du dossier (souvent oubliée ou vague), et la gestion du temps (développement inachevé).
\end{itemize}
L'activité, bien que dense, couvre l'intégralité des savoir-faire « Production des textes : Dissertation » du programme de Première Technique et constitue une simulation probatoire fidèle, évaluée sur 20 points selon la grille officielle.

\newpage

\coursec{Méthodes et exercices résolus}

\coursec{Dissertation : rédiger une dissertation}

\courssub{Le texte argumentatif et la dissertation : repères}

\begin{methode}[Identifier les caractéristiques du texte argumentatif et de la dissertation]
\textbf{Objectif :} Distinguer le texte argumentatif des autres types de textes et connaître les exigences spécifiques de la dissertation au Probatoire.
\begin{enumerate}
    \item \textbf{Repérer la visée argumentative :} L'auteur prend position (thèse), cherche à convaincre ou persuader un destinataire.
    \item \textbf{Identifier la structure classique :} Thèse (affirmation) $\rightarrow$ Arguments (raisons logiques) $\rightarrow$ Exemples/Illustrations (faits, citations, statistiques) $\rightarrow$ Connecteurs logiques (donc, car, en effet, toutefois).
    \item \textbf{Noter les modalités :} \cle{Argumentation directe} (je pense que...) ou \cle{indirecte} (ironie, fable, allégorie).
    \item \textbf{Spécificités de la dissertation :} C'est une \emph{argumentation écrite, organisée, scolaire}. Elle impose une structure rigide : \cle{Introduction} (amorce, sujet posé, problématique, annonce du plan), \cle{Développement} (2 ou 3 parties équilibrées, paragraphes argumentatifs), \cle{Conclusion} (bilan, réponse, ouverture). Pas de « je » personnel (sauf sujet le permettant), langue soutenue, connecteurs obligatoires.
\end{enumerate}
\end{methode}

\begin{exoresolu}[Analyse d'un extrait et identification du genre]
\textbf{Énoncé :} Lisez l'extrait suivant : « \emph{L'école ne doit pas seulement instruire, elle doit éduquer. En effet, l'instruction sans éducation produit des esprits brillants mais sans conscience morale. C'est pourquoi l'enseignement civique est indispensable.} »
\begin{enumerate}
    \item Quel est le type de texte ? Justifiez par deux indices.
    \item Identifiez la thèse, l'argument et l'exemple (ou l'illustration).
    \item Ce texte est-il une dissertation ? Justifiez.
\end{enumerate}
\tcblower
\textbf{Solution détaillée :}
\begin{enumerate}
    \item \textbf{Type :} Texte \cle{argumentatif} (visée : convaincre de la nécessité de l'éducation civique).
    \textbf{Indices :}
    \begin{itemize}
        \item Présence d'une \cle{thèse} explicite : « L'école ne doit pas seulement instruire, elle doit éduquer. »
        \item Utilisation de \cle{connecteurs logiques} : « En effet » (explication/justification), « C'est pourquoi » (conséquence/conclusion).
    \end{itemize}
    \item \textbf{Thèse :} L'école a pour mission l'éducation (morale/civique) au-delà de l'instruction.
    \textbf{Argument :} L'instruction seule produit des individus compétents mais dépourvus de conscience morale (risque social).
    \textbf{Illustration/Exemple :} La référence à l'« enseignement civique » comme solution nécessaire (exemple institutionnel).
    \item \textbf{Non, ce n'est pas une dissertation.} C'est un \emph{paragraphe argumentatif} isolé. Il manque la structure macro-textuelle imposée : pas d'introduction formelle (amorce, problématique, annonce de plan), pas de développement en plusieurs parties équilibrées avec transitions, pas de conclusion en deux temps (bilan + ouverture).
\end{enumerate}
\end{exoresolu}

\courssub{Morphosyntaxe : Les types de phrase (obligatoires et facultatifs)}

\begin{definition}[Les quatre types de phrase obligatoires]
Chaque phrase exprime une \cle{modalité énonciative} identifiable à sa ponctuation finale et à sa structure syntaxique.
\begin{center}
\begin{tabularx}{\linewidth}{|l|X|c|X|}\hline
\textbf{Type} & \textbf{Fonction / Visée} & \textbf{Ponctuation} & \textbf{Marqueurs syntaxiques fréquents} \\ \hline
\cle{Déclarative} (ou assertive) & Affirmer, nier, informer, raconter. & . (point) & Sujet + verbe (ordre normal), temps de l'indicatif. \\ \hline
\cle{Interrogative} & Poser une question (totale, partielle, rhétorique). & ? & Inversion (Verb-Sujet), \emph{est-ce que}, mot interrogatif (qui, quand, pourquoi...), intonation montante. \\ \hline
\cle{Impérative} (ou injonctive) & Ordonner, conseiller, supplier, interdire. & ! (ou .) & Verbe à l'impératif (souvent sans sujet), ou infinitif (consignes), subjonctif (que + subj.). \\ \hline
\cle{Exclamative} & Exprimer un sentiment fort (joie, colère, admiration...). & ! & \emph{Que} + indicatif/subjonctif, \emph{Comme} + adj., nominales (``Quel bonheur !''), interjections. \\ \hline
\end{tabularx}
\end{center}
\end{definition}

\begin{definition}[Les types de phrase facultatifs (variantes)]
\begin{itemize}
    \item \cle{Phrase nominale} : Sans verbe conjugué (titre, légende, style télégraphique, effet de style). Ex: « Silence ! », « La nuit, le froid, l'attente. »
    \item \cle{Phrase elliptique} : Suppression d'éléments récupérables du contexte (dialogue). Ex: « -- Tu viens ? -- Oui. » (pour « Oui, je viens. »)
    \item \cle{Phrase exclamative indirecte} : Subordonnée complétive introduite par \emph{que} (ex: « Je suis surpris qu'il parte. »).
\end{itemize}
\end{definition}

\begin{attention}[Pièges à l'examen]
\begin{itemize}
    \item Ne pas confondre \cle{phrase interrogative directe} (``Pourquoi viens-tu ?'') et \cle{indirecte} (``Il demande pourquoi je viens.'' $\rightarrow$ phrase déclarative).
    \item L'\cle{infinitif} peut valoir impératif (recettes, modes d'emploi : ``Mélanger le tout.'') ou nominal (sujet : ``Manger est nécessaire.'').
    \item La \cle{question rhétorique} (argumentation) est une interrogative en forme mais une assertive en force (ne demande pas de réponse).
\end{itemize}
\end{attention}

\courssub{Analyse du sujet et construction du plan détaillé}

\begin{methode}[Analyser un sujet de dissertation et construire un plan détaillé dialectique]
\textbf{Objectif :} Passer du sujet brut à un plan détaillé prêt à écrire (thèses, arguments, exemples, transitions).
\begin{enumerate}
    \item \textbf{Dégager les mots-clés :} Surligner les termes forts, définir leurs sens (propres, figurés, connotés).
    \item \textbf{Formuler la problématique :} Transformer le sujet en question ouverte, débattue, qui appelle une réponse nuancée (souvent « Dans quelle mesure... » ou « Comment... »).
    \item \textbf{Chercher les idées (brainstorming) :} Pour / Contre / Nuance. Classer par grands axes.
    \item \textbf{Choisir la structure du plan :}
    \begin{itemize}
        \item \cle{Plan dialectique (Thèse / Antithèse / Synthèse)} : Idéal pour « Sujet : affirmation ? » ou « Dans quelle mesure ? ».
        \item \cle{Plan thématique (Causes / Conséquences / Solutions)} ou \cle{Analytique (Aspects sociaux / politiques / économiques)}.
    \end{itemize}
    \item \textbf{Rédiger le \cle{plan détaillé} :} Pour chaque grande partie (I, II, III) : \textbf{Idée directrice} (phrase complète) $\rightarrow$ \textbf{Sous-parties} (A, B) avec \textbf{arguments} et \textbf{exemples précis} (références littéraires, historiques, faits de société, \emph{Le Lion et la Perle}, actualité camerounaise) $\rightarrow$ \textbf{Transition} vers la partie suivante.
\end{enumerate}
\end{methode}

\begin{exoresolu}[Analyse et plan détaillé : « La tradition est-elle un obstacle au progrès ? »]
\textbf{Énoncé :} Sujet de dissertation : « \emph{La tradition est-elle un obstacle au progrès ?} ». Faites l'analyse complète et proposez un plan détaillé en 3 parties (dialectique).
\tcblower
\textbf{Solution détaillée :}

\textbf{1. Mots-clés \& Définitions :}
\begin{itemize}
    \item \cle{Tradition} : Héritage culturel, coutumes, valeurs, rites transmis (ex: chefferies, rites d'initiation, polygamie, respect des aînés au Cameroun).
    \item \cle{Progrès} : Amélioration technique, scientifique, social, moral (ex: médecine moderne, éducation filles/garçons, droits de l'homme, technologie).
    \item \cle{Obstacle} : Frein, blocage, mais aussi socle, repère, résistance positive.
\end{itemize}

\textbf{2. Problématique :} \emph{Dans quelle mesure la tradition, loin d'être un simple frein, peut-elle être à la fois un obstacle à certaines modernités et un levier nécessaire pour un progrès durable et identitaire ?}

\textbf{3. Plan détaillé (Dialectique) :}

\textbf{INTRODUCTION} (Amorce : citation Soyinka ou proverbe Bamoun $\rightarrow$ Sujet posé $\rightarrow$ Problématique $\rightarrow$ Annonce du plan)

\textbf{I. La tradition comme frein au progrès : l'ancrage dans le passé bloque l'innovation.}
\begin{itemize}
    \item \textbf{A. Le poids des coutumes figées :} Résistance au changement (ex: \emph{Le Lion et la Perle} : Baroka refuse la route/le progrès technique pour garder son pouvoir ; excision, mariages forcés freinant l'éducation des filles).
    \item \textbf{B. La gérontocratie et le conservatisme :} Pouvoir des aînés étouffant la jeunesse (ex: Lakunle vs Baroka/Sidi ; difficulté d'entreprendre sans aval familial).
    \item \textit{Transition :} Cependant, réduire la tradition à un blocage est réducteur ; elle porte des valeurs structurantes.
\end{itemize}

\textbf{II. La tradition comme socle et levier du progrès : identité, éthique et cohésion.}
\begin{itemize}
    \item \textbf{A. Garante de l'identité et de la cohésion sociale :} Repères moraux (respect, solidarité, \cle{ubuntu} « Je suis parce que nous sommes »), gestion communautaire des conflits (justice coutumière).
    \item \textbf{B. Source d'innovation endogène :} Savoirs locaux (pharmacopée, agriculture durable, architecture) base de progrès adapté (ex: recherche sur \emph{Prunus africana} au Cameroun).
    \item \textit{Transition :} Le progrès réel ne naît pas de la table rase mais du dialogue critique entre héritage et modernité.
\end{itemize}

\textbf{III. Vers une synthèse dynamique : le « progrès traditionné » ou la modernité enracinée.}
\begin{itemize}
    \item \textbf{A. Le tri critique (héritage critique) :} Garder l'humain (solidarité), rejeter l'inhumain (excision, castes). Ex: Code de la personne et de la famille au Cameroun (1981/2016) modernisant le droit coutumier.
    \item \textbf{B. L'exemple de \emph{Le Lion et la Perle} :} Sidi choisit Baroka (tradition) mais avec lucidité ; Lakunle (modernité superficielle) échoue. La voie médiane : moderniser les mentalités sans aliéner la culture (ex: éducation bilingue, valorisation langues nationales).
    \item \textbf{C. Le rôle de l'école et de l'élite :} Former des « enracinés-ouverts » (Cheikh Hamidou Kane, \emph{L'Aventure ambiguë}).
\end{itemize}

\textbf{CONCLUSION} (Bilan : ni table rase, ni momification $\rightarrow$ Réponse : obstacle partiel, mais ressource majeure si réinventée $\rightarrow$ Ouverture : rôle de la jeunesse / numérique pour réinventer la tradition).
\end{exoresolu}

\courssub{Rédiger l'introduction}

\begin{methode}[Rédiger une introduction de dissertation en 4 étapes (entonnoir)]
\textbf{Objectif :} Produire un paragraphe unique, fluide, qui captive, cadre le débat et annonce la structure.
\begin{enumerate}
    \item \textbf{L'amorce (Accroche) :} 1 phrase. \cle{Obligatoire :} en lien direct avec le sujet. Sources : Citation d'auteur (Soyinka, Senghor, Kafka, proverbe africain), fait d'actualité précis, chiffre clé, définition étymologique, référence à l'œuvre au programme (\emph{Le Lion et la Perle}). \textbf{Interdit :} « De tout temps... », « Depuis la nuit des temps... », banalités.
    \item \textbf{Le sujet posé (Explication) :} 1 à 2 phrases. Définir les termes clés, présenter le débat, situer le contexte (social, littéraire, historique).
    \item \textbf{La problématique :} 1 phrase (question) ou 2 (affirmation + question). Doit naître logiquement du sujet posé. Contient la tension (opposition, paradoxe). \cle{Mots-clés du sujet obligatoirement repris}.
    \item \textbf{L'annonce du plan :} 1 phrase. « Nous verrons d'abord que... (I), puis nous montrerons que... (II), pour enfin comprendre que... (III). » Utiliser les connecteurs : \cle{Premièrement / Ensuite / Enfin} ou \cle{Dans un premier temps / Dans un second temps / Finalement}. \textbf{Pas de « Je vais dire... » mais « Nous analyserons... »}.
\end{enumerate}
\end{methode}

\begin{exoresolu}[Rédaction de l'introduction pour le sujet « La tradition est-elle un obstacle au progrès ? »]
\textbf{Énoncé :} Rédigez l'introduction complète en appliquant la méthode des 4 étapes.
\tcblower
\textbf{Solution détaillée :}

\emph{« La tradition n'est pas le culte des cendres, mais la garde du feu », écrivait Gustav Mahler. Cette image du feu, qu'il faut alimenter pour qu'il vive sans qu'il ne consume la maison, illustre parfaitement la tension qui existe entre l'héritage ancestral et le mouvement de l'histoire. Au Cameroun comme ailleurs, le choc des modernités (technique, juridique, éducative) avec les coutumes (chefferies, rites, solidarité clanique) pose avec acuité la question de leur compatibilité. Si la tradition désigne l'ensemble des usages, croyances et valeurs transmis de génération en génération, et le progrès l'amélioration continue des conditions de vie par la technique et la raison, leur rencontre est souvent perçue comme un choc. \textbf{Dans quelle mesure la tradition, loin d'être un simple frein, peut-elle être à la fois un obstacle à certaines modernités destructrices et le socle indispensable d'un progrès humain et durable ?} Nous analyserons d'abord comment la tradition peut freiner l'innovation par son conservatisme (I), avant de montrer qu'elle constitue paradoxalement un réservoir de valeurs et de savoirs propices au développement (II), pour enfin dégager les conditions d'une synthèse féconde entre enracinement et ouverture (III).}

\textbf{Analyse de la rédaction :}
\begin{itemize}
    \item \textbf{Amorce :} Citation Mahler (métaphore feu/cendres) $\rightarrow$ lien direct.
    \item \textbf{Sujet posé :} Définitions claires, contextualisation Cameroun/\emph{Le Lion et la Perle} implicite.
    \item \textbf{Problématique :} Question nuancée (« loin d'être un simple frein », « à la fois... et »), reprise mots-clés.
    \item \textbf{Annonce :} Verbes d'action (analyser, montrer, dégager), structure I/II/III claire, connecteurs.
\end{itemize}
\end{exoresolu}

\courssub{Rédiger le développement : le paragraphe argumentatif}

\begin{methode}[Construire un paragraphe argumentatif standard (I.D.A.E.)]
\textbf{Objectif :} Rédiger des paragraphes cohérents, autonomes et enchaînés logiquement.
\begin{enumerate}
    \item \textbf{Idée directrice (Phrase sujet) :} 1 phrase en tête de paragraphe. Résume l'argument principal de la sous-partie (A ou B). Commence par un connecteur de structure (\cle{Premièrement}, \cle{En effet}, \cle{Par ailleurs}, \cle{Toutefois}).
    \item \textbf{Développement / Argumentation :} 2 à 3 phrases. Expliquer le mécanisme logique : \cle{Pourquoi ? Comment ?} (Causalité, analogie, autorité, fait majoritaire). Utiliser les connecteurs logiques (\cle{car}, \cle{puisque}, \cle{ainsi}, \cle{conséquemment}).
    \item \textbf{Exemple / Illustration :} 1 à 2 phrases. \cle{Concret, daté, sourcé.} Référence littéraire (\emph{Le Lion et la Perle} : scène précise, réplique), fait historique, statistique, actualité camerounaise, proverbe. \textbf{Interdit :} « Par exemple, untel a dit... » sans précision. Préférer : « Dans \emph{Le Lion et la Perle}, la scène du miroir (Acte 2) montre Sidi... »
    \item \textbf{Évaluation / Mini-synthèse / Transition :} 1 phrase. \cle{Valider l'argument} (« Ainsi, ... ») et \cle{Ouvrir sur le suivant} (« Cependant, cet obstacle n'est pas absolu... » $\rightarrow$ partie suivante). Assure la cohésion textuelle.
\end{enumerate}
\textbf{Rappel langue :} Phrases \cle{déclaratives} majoritaires. \cle{Connecteurs} variés (addition, opposition, conséquence, concession). Vocabulaire précis (éviter « chose », « truc », « très important »).
\end{methode}

\begin{exoresolu}[Rédaction d'un paragraphe : Sous-partie I.A « Le poids des coutumes figées »]
\textbf{Énoncé :} Rédigez le paragraphe complet correspondant à la sous-partie I.A du plan détaillé : « Le poids des coutumes figées : résistance au changement ».
\tcblower
\textbf{Solution détaillée :}

\emph{Premièrement, la tradition agit comme un frein puissant lorsque les coutumes se figent en dogmes immuables, refusant toute remise en cause au nom du respect des ancêtres. En effet, la sacralisation du passé interdit l'innovation technique ou sociale perçue comme une trahison : toute nouveauté est suspectée de rompre l'ordre cosmique ou social établi. \textbf{Dans \emph{Le Lion et la Perle}, Wole Soyinka met en scène ce blocage à travers le personnage de Baroka, le Bale d'Ilujinle, qui sabote délibérément l'arrivée du chemin de fer dans son village.} Craignant que la modernité technique n'érode son autorité traditionnelle fondée sur le contrôle de l'espace et du temps, il corrompt l'arpenteur pour détourner le tracé. Cette résistance n'est pas seulement individuelle : elle structure la société, comme en témoigne la persistance de l'excision ou des mariages précoces dans certaines régions du Cameroun, freinant l'éducation des filles et leur autonomie économique, malgré les lois républicaines. \textbf{Ainsi, quand la tradition se fige en conservatisme identitaire, elle devient un obstacle objectif au progrès social et technique. Toutefois, cette vision réductrice occulte la face dynamique de l'héritage culturel.}}

\textbf{Analyse I.D.A.E. :}
\begin{itemize}
    \item \textbf{I (Idée directrice) :} « Premièrement, la tradition agit comme un frein puissant... »
    \item \textbf{D (Développement) :} « En effet, la sacralisation du passé interdit... » (Explication causale).
    \item \textbf{A (Argument/Autorité/Exemple) :} « Dans \emph{Le Lion et la Perle}... » (Exemple littéraire précis, scène, fonction). + Exemple sociologique Cameroun (excision/mariage précoce).
    \item \textbf{E (Évaluation/Transition) :} « Ainsi... Toutefois... » (Bilan + Ouverture vers II).
\end{itemize}
\end{exoresolu}

\courssub{Rédiger la conclusion}

\begin{methode}[Rédiger une conclusion en 2 temps (Bilan + Ouverture)]
\textbf{Objectif :} Clore le devoir sans apporter d'argument nouveau, en répondant explicitement à la problématique.
\begin{enumerate}
    \item \textbf{Le bilan synthétique (Rappel structuré) :} 2 à 3 phrases. Reprendre la logique du plan (I, II, III) sans lister mécaniquement. Utiliser : \cle{Au terme de cette analyse}, \cle{Il apparaît que}, \cle{Force est de constater que}. Synthétiser la dialectique : « Si la tradition peut freiner... (I), elle reste pourtant le socle... (II), à condition d'opérer un tri critique... (III) ».
    \item \textbf{L'ouverture (Perspective) :} 1 à 2 phrases. Élargir le débat, poser une question nouvelle, lier à une autre notion, citer un auteur/proverbe prospectif. \textbf{Interdit :} « Pour conclure, je dirais que... », ouverture hors-sujet (ex: « Et si on parlait de l'écologie ? »), question rhétorique sans lien.
    \begin{itemize}
        \item \cle{Types d'ouverture :} Littéraire (citation), Historique (évolution future), Sociétale (défi actuel), Philosophique (notion connexe : modernité, identité, liberté).
    \end{itemize}
\end{enumerate}
\end{methode}

\begin{exoresolu}[Rédaction de la conclusion pour le sujet « La tradition est-elle un obstacle au progrès ? »]
\textbf{Énoncé :} Rédigez la conclusion complète.
\tcblower
\textbf{Solution détaillée :}

\emph{Au terme de cette analyse, il apparaît que la tradition n'est pas un bloc monolithique figé dans le passé. Si elle peut constituer un obstacle lorsque le conservatisme sacralise l'immobilisme et rejette l'innovation technique ou juridique (I), elle révèle simultanément une capacité d'adaptation et une richesse axiologique — solidarité, respect, savoirs endogènes — qui fondent la cohésion sociale et orientent le progrès vers l'humain (II). La réponse à notre problématique réside donc dans la dialectique du « tri critique » : le progrès durable ne naît ni de la table rase, ni de la momification, mais de l'invention d'une modernité enracinée (III), comme l'illustre l'évolution du droit camerounais intégrant la coutume rénovée. \textbf{On peut dès lors se demander, avec Cheikh Hamidou Kane, si « l'aventure ambiguë » de la modernité africaine ne consiste pas, pour la jeunesse formée à l'école occidentale, à réinventer la tradition par le numérique et la citoyenneté active, transformant le « feu des ancêtres » en lumière pour l'avenir.}}

\textbf{Analyse :}
\begin{itemize}
    \item \textbf{Bilan :} Synthèse dialectique (I/II/III fusionnés), réponse nette (« tri critique », « modernité enracinée »).
    \item \textbf{Ouverture :} Référence explicite à \emph{L'Aventure ambiguë} (programme Terminale mais culture générale attendue), lien jeunesse/numérique/citoyenneté, métaphore du feu reprise de l'amorce (cohérence circulaire).
\end{itemize}
\end{exoresolu}

\courssub{Sémantique / Lexicologie : Synonymie et antonymie}

\begin{definition}[Synonymie]
Relation entre deux mots (synonymes) de même catégorie grammaticale qui ont un \cle{sens très proche} dans un contexte donné.
\begin{itemize}
    \item \cle{Synonymie totale (rare) :} Échangeabilité dans tous les contextes sans nuance (ex: \emph{véhicule / automobile} -- technique).
    \item \cle{Synonymie partielle (courante) :} Proximité de sens mais différences de \cle{registre} (familier/courant/soutenu), de \cle{connotation} (péjoratif/laudatif/neutre), d'\cle{intensité} ou de \cle{valeur aspectuelle}.
    \item \textbf{Exemples :} \emph{Manger / Dévorer / Goinfrer / S'alimenter} (intensité, registre) ; \emph{Maison / Demeure / Logis / Bicoque} (connotation, registre).
\end{itemize}
\end{definition}

\begin{definition}[Antonymie]
Relation entre deux mots (antonymes) de même catégorie grammaticale dont les sens sont \cle{opposés}.
\begin{itemize}
    \item \cle{Antonymes complémentaires (contradictoires) :} L'affirmation de l'un implique la négation de l'autre, pas de milieu. Ex: \emph{Vivant / Mort, Vrai / Faux, Présent / Absent}.
    \item \cle{Antonymes gradués (contraires) :} Opposition sur une échelle, existence de degrés intermédiaires. Ex: \emph{Chaud / Tiède / Froid ; Grand / Moyen / Petit ; Aimer / Indifférent / Détester}.
    \item \cle{Antonymes réciproques (corrélatifs) :} Relation inverse entre deux acteurs. Ex: \emph{Acheter / Vendre, Donner / Recevoir, Parent / Enfant, Au-dessus / Au-dessous}.
    \item \textbf{Formation morphologique :} Préfixes \emph{in-, im-, ir-, il-, dé-, dés-, non-, anti-, contre-, mal-, mé-} (ex: \emph{juste / injuste, faire / défaire, accord / désaccord}).
\end{itemize}
\end{definition}

\begin{exemplebox}[Utilité en dissertation]
\begin{itemize}
    \item \textbf{Éviter les répétitions} (cohésion lexicale) : varier le vocabulaire (ex: \emph{progrès / avancée / développement / essor / modernisation}).
    \item \textbf{Préciser la nuance} : choisir le mot exact (ex: \emph{freiner} vs \emph{entraver} vs \emph{bloquer} vs \emph{ralentir}).
    \item \textbf{Construire l'argumentation} : jouer sur les oppositions (antithèse) ou les renforcements (synonymie graduelle). Ex: « Non seulement la tradition \emph{freine} (faible), mais elle \emph{entrave} (fort) l'innovation. »
\end{itemize}
\end{exemplebox}

\courssub{Confrontation, évaluation et amélioration des productions}

\begin{methode}[Auto-évaluation et réécriture à l'aide d'une grille critériée (Bac)]
\textbf{Objectif :} Devenir correcteur de sa propre copie pour gagner les points de la « qualité de la langue » et de la « structure ».
\begin{enumerate}
    \item \textbf{Phase 1 : Vérification structurelle (Squelette).} Cocher : Introduction (4 étapes) ? Plan annoncé respecté ? 2 ou 3 parties équilibrées ? Paragraphes I.D.A.E. ? Transitions explicites ? Conclusion (Bilan + Ouverture) ?
    \item \textbf{Phase 2 : Vérification argumentative (Muscle).} Chaque idée directrice répond-elle à la problématique ? Arguments logiques (pas d'assertions gratuites) ? Exemples \cle{précis, variés, sourcés} (Littérature, Histoire, Actualité Cameroun, Vie quotidienne) ? Pas d'exemples « bateau » (ex: « Dans la vie de tous les jours... »).
    \item \textbf{Phase 3 : Vérification linguistique (Peau).} \cle{Connecteurs} variés (pas 5 « de plus » de suite) ? \cle{Relatives, subordonnées, participation} pour la fluidité ? \cle{Vocabulaire précis} (pas de répétitions, champ lexical du sujet) ? \cle{Orthographe/Accords} (relecture à l'envers) ? \cle{Ponctuation} (virgules avant « mais », « car », « or » ; points-virgules pour énumérations complexes).
    \item \textbf{Phase 4 : Réécriture ciblée.} Réécrire \textbf{un seul} paragraphe faible identifié. Appliquer I.D.A.E. strict. Remplacer verbes faibles (être, avoir, faire) par verbes d'action (structurer, fonder, entraver, subvertir, hybrider).
\end{enumerate}
\end{methode}

\begin{exoresolu}[Correction d'un paragraphe faible : « La tradition c'est bien »]
\textbf{Énoncé :} Voici un paragraphe d'élève pour la partie II (La tradition comme levier). Appliquez la grille et proposez une version améliorée.
\emph{« Deuxièmement, la tradition c'est bien pour le progrès. Parce que elle donne des valeurs. Par exemple au Cameroun on respecte les vieux. C'est important pour la société. Aussi les traditions comme le Ngondo aident le tourisme. Donc la tradition aide le progrès. »}
\tcblower
\textbf{Solution détaillée :}

\textbf{Diagnostic (Grille) :}
\begin{itemize}
    \item \textbf{Structure :} Idée directrice floue (« c'est bien »), pas de transition entrante correcte (« Deuxièmement » seul).
    \item \textbf{Argumentation :} « Parce que elle donne des valeurs » (syntaxe fautive, argument circulaire). Manque de logique.
    \item \textbf{Exemples :} « On respecte les vieux » (trop vague, oral), « Ngondo aide le tourisme » (réducteur, pas analysé).
    \item \textbf{Langue :} Registre familier (« c'est bien », « parce que elle »), répétitions, ponctuation faible, accords (parce qu'\textbf{elle}).
\end{itemize}

\textbf{Version améliorée (Réécriture I.D.A.E.) :}

\emph{Par ailleurs, la tradition agit comme un puissant levier de progrès en fournissant le capital symbolique et éthique indispensable à toute cohésion collective. En effet, les valeurs de respect de la parole donnée, de solidarité intergénérationnelle et de responsabilité communautaire, transmises par l'éducation coutumière, structurent un « vivre-ensemble » que l'individualisme moderne peine à garantir. \textbf{Au Cameroun, l'institution du \cle{Ngondo} chez les Sawa, au-delà de sa dimension festive et touristique, incarne cette fonction régulatrice : le culte des ancêtres (Jengu) renforce l'identité collective et légitime l'autorité morale des chefs, gage de stabilité politique locale.} De même, les mécanismes traditionnels de résolution des conflits (palabre, médiation) inspirent aujourd'hui les modes alternatifs de règlement des différends (MARD) dans le système judiciaire moderne. \textbf{Ainsi, loin d'être un folklore, la tradition offre des ressources institutionnelles et éthiques que le progrès institutionnel gagnerait à institutionaliser plutôt qu'à marginaliser.}}

\textbf{Gains :} Registre soutenu, connecteurs (\emph{Par ailleurs, En effet, De même, Ainsi}), vocabulaire précis (\emph{capital symbolique, cohésion, légitimer, institutionnaliser}), exemple Ngondo analysé (fonction sociale + lien modernité), syntaxe complexe.
\end{exoresolu}

\begin{attention}[Les 5 erreurs qui coûtent des points au Probatoire / Bac]
\begin{enumerate}
    \item \textbf{Introduction « hors-sujet » ou « catalogue » :} Amorce générique (« De tout temps les hommes... »), définitions dictionnaire recopiées, problématique absente ou fermée (réponse oui/non), plan annoncé ne correspondant pas au développement.
    \item \textbf{Développement « liste de courses » :} Accumulation d'exemples sans idée directrice, sans argument logique liant l'exemple à la thèse. Paragraphes « exemple unique » (tout le paragraphe = résumé de \emph{Le Lion et la Perle} sans analyse).
    \item \textbf{Absence de connecteurs / Relances :} Texte haché, phrases courtes juxtaposées. Pas de transition entre I et II, II et III (rupture du fil conducteur). Utilisation exclusive de « Et », « Mais », « Aussi » en début de phrase.
    \item \textbf{Conclusion « nouvelle partie » :} Introduction d'un argument nouveau, d'un nouvel exemple, ou d'une idée non traitée. Ouverture « parachutée » (ex: « Pour finir, parlons de l'environnement »). Absence de réponse explicite à la problématique.
    \item \textbf{Relâchement linguistique majeur :} Fautes d'accords systématiques (sujet/verbe, GN), confusion \emph{a/à, et/est, ou/où, son/sont, ce/se, ce/c'est}, registres familiers (\emph{truc, bidule, genre, grave}), phrases sans verbe, ponctuation aléatoire. \textbf{Conséquence :} Pénalité « Qualité de la langue » (souvent 4 à 6 points sur 20).
\end{enumerate}
\end{attention}

\newpage

\coursec{Exercices}

\courssub{Je vérifie mes connaissances}

\begin{exercice}[QCM : caractéristiques du texte argumentatif]
Parmi les propositions suivantes, cochez celle qui correspond à la définition du texte argumentatif :
\begin{enumerate}
\item Un texte qui raconte une histoire avec des personnages et des dialogues.
\item Un texte qui expose des faits de manière neutre sans prendre position.
\item Un texte qui défend une thèse en utilisant des arguments et des exemples pour convaincre le lecteur.
\item Un texte qui exprime des sentiments personnels sous forme de vers.
\end{enumerate}
\end{exercice}

\begin{exercice}[Vrai ou Faux justifié : analyse du sujet]
Indiquez si les affirmations suivantes sont vraies ou fausses. Justifiez brièvement votre réponse.
\begin{enumerate}
\item L'analyse du sujet consiste uniquement à souligner les mots-clés.
\item La problématique est la question unique à laquelle la dissertation doit répondre.
\item L'annonce du plan doit révéler le contenu exact de chaque partie.
\item Le sujet « L'école est-elle le seul chemin de la réussite au Cameroun ? » est un sujet fermé.
\end{enumerate}
\end{exercice}

\begin{exercice}[Définitions : structure de la dissertation]
Donnez une définition précise pour chacun des termes suivants :
\begin{enumerate}
\item L'amorce (ou accroche) de l'introduction.
\item L'idée directrice d'un paragraphe argumentatif.
\item Le schéma I-A-E-L.
\item La conclusion ouverte.
\end{enumerate}
\end{exercice}

\begin{exercice}[Application directe : types de phrase]
Transformez la phrase déclarative suivante en les trois autres types de phrase obligatoires, puis en deux types de phrase facultatifs :
\emph{« L'éducation technique permet aux jeunes d'acquérir un métier. »}
\end{exercice}

\courssub{Je m'entraîne}

\begin{exercice}[Analyse de sujet et plan détaillé]
Pour le sujet : \emph{« Faut-il préserver les langues locales dans l'enseignement technique au Cameroun ? »}
\begin{enumerate}
\item Identifiez les termes clés et définissez-les.
\item Formulez la problématique.
\item Proposez un plan dialectique (thèse, antithèse, synthèse) avec deux arguments par partie et un exemple précis pour chaque argument.
\item Rédigez l'annonce du plan.
\end{enumerate}
\end{exercice}

\begin{exercice}[Rédaction d'introductions complètes]
Rédigez une introduction complète (amorce, définition des termes, problématique, annonce du plan) pour chacun des trois sujets suivants :
\begin{enumerate}
\item « Le numérique est-il une chance pour la jeunesse camerounaise ? »
\item « La corruption freine-t-elle le développement du Cameroun ? »
\item « L'apprentissage en entreprise vaut-il mieux que la formation scolaire ? »
\end{enumerate}
\end{exercice}

\begin{exercice}[Réécriture de paragraphe : schéma I-A-E-L]
Voici un paragraphe faible :
\emph{« L'école aide à réussir. Beaucoup de gens ont réussi grâce à l'école. Par exemple, le ministre est allé à l'école. Donc l'école est importante. »}
\begin{enumerate}
\item Identifiez les manquements par rapport au schéma I-A-E-L.
\item Réécrivez ce paragraphe en respectant scrupuleusement le schéma I-A-E-L (Idée directrice, Argument, Exemple, Lien/Transition).
\item Intégrez au moins une phrase de type facultatif (exclamative, interrogative rhétorique ou impérative) et soulignez-la.
\end{enumerate}
\end{exercice}

\begin{exercice}[Étude de langue sur \emph{Le Lion et la perle}]
Lisez l'extrait suivant (acte 1, scène 1, Lakunle s'adressant à Sidi) :
\emph{« Ignorante fille de brousse, tu ne sais rien ! Le monde change, il faut changer avec lui. La mariée qui paie le prix de la mariée est une coutume barbare, humiliante. Nous devons être modernes, civilisés. »}
\begin{enumerate}
\item Relevez et classez toutes les phrases selon leur type (déclarative, interrogative, exclamative, impérative) et leur forme (affirmative, négative).
\item Identifiez cinq couples de synonymes ou d'antonymes dans l'extrait.
\item Expliquez l'effet produit par l'accumulation de phrases impératives et exclamatives dans la bouche de Lakunle.
\end{enumerate}
\end{exercice}

\courssub{Je me prépare au Bac}

\begin{exercice}[Sujet type Baccalauréat technique — Dissertation complète (3 h)]
\textbf{Sujet :} \emph{« L'école est-elle le seul chemin de la réussite au Cameroun ? »}
\begin{enumerate}
\item \textbf{Analyse du sujet (30 min) :} Dégagez les termes clés, formulez la problématique, proposez un plan détaillé (dialectique ou thématique) avec arguments et exemples pour chaque sous-partie.
\item \textbf{Rédaction (2 h) :} Rédigez la dissertation complète : introduction, développement (trois parties, deux paragraphes par partie), conclusion. Respectez les normes de présentation (paragraphes sautés, connecteurs logiques, schéma I-A-E-L).
\item \textbf{Relecture (30 min) :} Vérifiez la cohérence du plan, la correction linguistique (orthographe, syntaxe, ponctuation), la présence des types de phrases variés et l'emploi pertinent du vocabulaire argumentatif.
\end{enumerate}
\end{exercice}

\begin{exercice}[Évaluation d'une copie anonyme — Grille officielle]
Vous êtes correcteur au Probatoire. Voici une copie anonyme (fournie en annexe lors de l'examen). À l'aide de la grille officielle de notation (sur 20 points : analyse du sujet 4, plan 4, introduction 3, développement 6, conclusion 3, langue 4) :
\begin{enumerate}
\item Attribuez une note à chaque critère et justifiez par des observations précises (points forts / points faibles).
\item Calculez la note finale sur 20.
\item Proposez trois améliorations concrètes et prioritaires pour que l'élève progresse.
\end{enumerate}
\end{exercice}

\begin{exercice}[Synthèse : plan dialectique, introduction et conclusion]
\textbf{Sujet :} \emph{« Faut-il préserver les langues locales dans l'enseignement technique au Cameroun ? »}
\begin{enumerate}
\item Construisez un plan dialectique complet (thèse, antithèse, synthèse) avec pour chaque grande partie : deux arguments distincts, un exemple précis camerounais, un connecteur de liaison.
\item Rédigez l'introduction complète (amorce, définitions, problématique, annonce du plan).
\item Rédigez la conclusion complète (bilan synthétique, réponse à la problématique, ouverture sur une perspective concrète : politique éducative, emploi, culture).
\end{enumerate}
\end{exercice}

\newpage

\coursec{Corrigés des exercices}

\coursec{Leçon 03 : Dissertation : rédiger une dissertation}

\courssub{Le texte argumentatif et la dissertation : repères}

\begin{corrige}[Exercice 1 — QCM : caractéristiques du texte argumentatif]
\textbf{Bonne réponse : proposition 3.}

Un \cle{texte argumentatif} est un texte qui défend une \cle{thèse} (un point de vue) en mobilisant des \cle{arguments} (raisons) et des \cle{exemples} (illustrations concrètes) afin de convaincre ou de persuader le lecteur.

\textbf{Justification de l'élimination des autres propositions :}
\begin{itemize}
\item La proposition 1 décrit le \cle{texte narratif} (récit, personnages, dialogues) : ce n'est pas l'argumentation.
\item La proposition 2 décrit le \cle{texte informatif ou expositif} (neutralité, absence de prise de position) ; or l'argumentatif suppose au contraire une prise de position.
\item La proposition 4 décrit le \cle{texte poétique} (expression lyrique des sentiments, vers) ; il peut parfois argumenter, mais ce n'est pas sa définition.
\end{itemize}

\textbf{Point de vigilance :} ne pas confondre \emph{convaincre} (faire adhérer par la raison, registre de la dissertation) et \emph{persuader} (toucher les émotions) ; la dissertation relève d'abord de la démonstration rationnelle.
\end{corrige}

\begin{corrige}[Exercice 2 — Vrai ou Faux justifié : analyse du sujet]
\begin{enumerate}
\item \textbf{Faux.} Souligner les mots-clés n'est que la première étape de l'analyse du sujet. Il faut ensuite \emph{définir} ces termes, dégager les présupposés, repérer le type de sujet (fermé ou ouvert), formuler la \cle{problématique} et construire le plan.
\item \textbf{Vrai.} La problématique est la question centrale, unique et directrice, à laquelle toute la dissertation apporte une réponse argumentée. Toutes les parties doivent y renvoyer.
\item \textbf{Faux.} L'annonce du plan indique seulement les \emph{grandes lignes} (les axes des parties), sans dévoiler le contenu détaillé ni les exemples. Dire « nous verrons d'abord que\ldots, puis que\ldots, enfin que\ldots » suffit ; révéler les arguments un par un alourdit l'introduction et gâche la démonstration.
\item \textbf{Faux.} Le sujet « L'école est-elle le seul chemin de la réussite au Cameroun ? » est un sujet \cle{ouvert} : il appelle une discussion nuancée (l'école est un chemin de réussite, mais pas le seul : entrepreneuriat, apprentissage, talent\ldots). Un sujet fermé imposerait une réponse unique, sans discussion possible (par exemple une affirmation à commenter sans contradiction).
\end{enumerate}
\end{corrige}

\begin{corrige}[Exercice 3 — Définitions : structure de la dissertation]
\begin{enumerate}
\item \textbf{L'amorce (ou accroche) :} première phrase de l'introduction, elle capte l'attention du lecteur en situant le sujet dans un contexte général (fait de société, constat, citation, événement d'actualité). Elle mène naturellement au sujet.
\item \textbf{L'idée directrice :} phrase qui ouvre le paragraphe argumentatif et en énonce clairement la thèse partielle (l'argument principal). Tout le paragraphe la développe et la justifie.
\item \textbf{Le schéma I-A-E-L :} structure type du paragraphe argumentatif : \textbf{I}dée directrice (ce qu'on affirme), \textbf{A}rgument (la raison qui la justifie), \textbf{E}xemple (l'illustration concrète et précise), \textbf{L}ien (conclusion du paragraphe et transition vers le suivant).
\item \textbf{La conclusion ouverte :} conclusion qui, après le bilan et la réponse à la problématique, s'achève par une \cle{ouverture} : une question nouvelle, une perspective (politique, sociale, culturelle) qui élargit la réflexion sans traiter un nouveau sujet.
\end{enumerate}
\end{corrige}

\begin{corrige}[Exercice 4 — Application directe : types de phrase]
Phrase de départ (déclarative affirmative) : \emph{« L'éducation technique permet aux jeunes d'acquérir un métier. »}

\textbf{Les trois autres types obligatoires :}
\begin{itemize}
\item \textbf{Interrogative :} « L'éducation technique permet-elle aux jeunes d'acquérir un métier ? »
\item \textbf{Exclamative :} « Comme l'éducation technique permet aux jeunes d'acquérir un métier ! » (ou : « Quelle chance que l'éducation technique permette aux jeunes d'acquérir un métier ! »)
\item \textbf{Impérative :} « Permettez aux jeunes d'acquérir un métier grâce à l'éducation technique ! » (ou : « Acquiers un métier grâce à l'éducation technique ! »)
\end{itemize}

\textbf{Deux types facultatifs :}
\begin{itemize}
\item \textbf{Phrase négative :} « L'éducation technique \textbf{ne permet pas} aux jeunes d'acquérir un métier. » (forme négative directe de la phrase de départ)
\item \textbf{Phrase emphatique (mise en relief) :} « C'est l'éducation technique qui permet aux jeunes d'acquérir un métier. »
\end{itemize}

\textbf{Point de vigilance :} la transformation ne doit pas dénaturer le sens ; la phrase impérative exige un verbe à l'impératif et la phrase exclamative un point d'exclamation. La forme négative doit conserver le même verbe et le même sens général (simple opposition affirmative/négative).
\end{corrige}

\courssub{Analyse du sujet et construction du plan détaillé}

\begin{corrige}[Exercice 5 — Analyse de sujet et plan détaillé]
\textbf{Sujet :} « Faut-il préserver les langues locales dans l'enseignement technique au Cameroun ? »

\textbf{1. Termes clés et définitions :}
\begin{itemize}
\item \emph{Préserver :} conserver, protéger contre la disparition, maintenir en usage.
\item \emph{Langues locales :} langues nationales camerounaises (ewondo, douala, fulfuldé, ghomala', medumba\ldots), par opposition aux langues officielles (français, anglais).
\item \emph{Enseignement technique :} formation scolaire orientée vers l'acquisition d'un métier et de compétences pratiques.
\item \emph{Cameroun :} cadre national, marqué par une forte diversité linguistique (plus de 250 langues).
\end{itemize}

\textbf{2. Problématique :} Dans quelle mesure l'intégration des langues locales dans l'enseignement technique est-elle nécessaire et possible, sans compromettre l'ouverture du jeune technicien sur le monde professionnel national et international ?

\textbf{3. Plan dialectique :}
\begin{itemize}
\item \textbf{Thèse — Il faut préserver les langues locales dans l'enseignement technique.}
Argument 1 : elles sont un patrimoine culturel et un facteur d'identité. Exemple : la journée des langues nationales organisée dans les lycées techniques de Yaoundé.
Argument 2 : elles facilitent la compréhension des notions techniques pour les apprenants du milieu rural. Exemple : un formateur en agropastoral de l'Extrême-Nord expliquant en fulfuldé les techniques d'élevage.
\item \textbf{Antithèse — Mais leur place doit être limitée.}
Argument 1 : la diversité linguistique rend impossible l'enseignement de toutes les langues. Exemple : dans un même lycée de Douala cohabitent des élèves de plus de dix langues différentes.
Argument 2 : le monde professionnel exige le français et l'anglais. Exemple : les concours d'entrée dans les grands corps techniques et les marchés internationaux se passent en langues officielles.
\item \textbf{Synthèse — Une coexistence équilibrée est possible.}
Argument 1 : les langues locales peuvent servir d'appui pédagogique, les langues officielles restant langues d'enseignement. Exemple : l'expérimentation du bilinguisme et des langues nationales dans certains établissements pilotes.
Argument 2 : la valorisation des langues locales peut créer des emplois (traduction, terminologie technique). Exemple : l'élaboration de lexiques techniques en langues nationales.
\end{itemize}

\textbf{4. Annonce du plan :} « Après avoir montré que les langues locales méritent d'être préservées dans l'enseignement technique, nous verrons les limites d'une telle entreprise, avant de dégager les conditions d'une coexistence harmonieuse entre langues locales et langues officielles. »
\end{corrige}

\courssub{Rédiger l'introduction}

\begin{corrige}[Exercice 6 — Rédaction d'introductions complètes]
\textbf{1. « Le numérique est-il une chance pour la jeunesse camerounaise ? »}

\emph{Amorce :} Depuis quelques années, les téléphones portables et l'Internet ont envahi le quotidien des jeunes Camerounais, des salles de classe aux quartiers. \emph{Définitions :} Le numérique désigne l'ensemble des technologies de l'information et de la communication ; la jeunesse renvoie à cette génération scolarisée ou en quête d'emploi. \emph{Problématique :} Mais cet outil est-il réellement un levier de réussite et d'épanouissement, ou un facteur de distraction et de dépendance ? \emph{Annonce du plan :} Nous montrerons d'abord que le numérique offre des opportunités réelles à la jeunesse, puis qu'il comporte des dangers sérieux, avant de dégager les conditions d'un usage profitable.

\textbf{2. « La corruption freine-t-elle le développement du Cameroun ? »}

\emph{Amorce :} Dans les débats publics camerounais, aucun mot ne revient aussi souvent que celui de corruption, dénoncée jusque dans les discours officiels. \emph{Définitions :} La corruption est le détournement des biens ou de l'autorité publics à des fins privées ; le développement est l'amélioration durable des conditions de vie d'un peuple. \emph{Problématique :} La corruption constitue-t-elle l'obstacle majeur au développement du pays, ou n'est-elle qu'un symptôme parmi d'autres ? \emph{Annonce du plan :} Nous verrons d'abord comment la corruption bloque le développement, puis qu'elle n'explique pas tout, avant de proposer les voies d'un assainissement durable.

\textbf{3. « L'apprentissage en entreprise vaut-il mieux que la formation scolaire ? »}

\emph{Amorce :} Face au chômage des jeunes diplômés, beaucoup de parents camerounais préfèrent placer leurs enfants chez un maître-artisan plutôt qu'au lycée. \emph{Définitions :} L'apprentissage en entreprise est la formation pratique sur le lieu de travail ; la formation scolaire est l'enseignement théorique et pratique dispensé dans un établissement. \emph{Problématique :} Laquelle de ces deux voies garantit le mieux l'insertion professionnelle du jeune technicien ? \emph{Annonce du plan :} Nous exposerons d'abord les mérites de l'apprentissage en entreprise, puis ceux de la formation scolaire, avant de montrer que leur complémentarité est la meilleure solution.

\textbf{Point de vigilance :} chaque introduction comporte bien les quatre étapes obligatoires ; l'annonce du plan reste générale et ne dévoile ni arguments ni exemples.
\end{corrige}

\courssub{Rédiger le développement : le paragraphe argumentatif}

\begin{corrige}[Exercice 7 — Réécriture de paragraphe : schéma I-A-E-L]
\textbf{1. Manquements relevés :}
\begin{itemize}
\item L'idée directrice (« L'école aide à réussir ») est vague et mal formulée.
\item L'argument (« Beaucoup de gens ont réussi grâce à l'école ») n'est pas développé : c'est une simple répétition de l'idée, sans raison explicative.
\item L'exemple (« le ministre est allé à l'école ») est imprécis : aucun nom, aucune référence vérifiable.
\item Le lien (« Donc l'école est importante ») est une conclusion faible qui ne fait pas de transition vers le paragraphe suivant.
\end{itemize}

\textbf{2. Réécriture selon I-A-E-L :}

\emph{[I] L'école constitue d'abord un instrument puissant de promotion sociale, car elle donne à chacun les connaissances et les diplômes exigés par le marché du travail. [A] En effet, la plupart des emplois qualifiés — administration, enseignement, santé, ingénierie — sont accessibles uniquement sur présentation de certifications obtenues par le parcours scolaire ; sans elles, les portes de la fonction publique restent fermées. [E] C'est ainsi que la majorité des cadres camerounais, des médecins du CHU de Yaoundé aux ingénieurs de la SONARA, doivent leur poste à un cursus scolaire et universitaire accompli avec succès. [L] \underline{Comment nier, dès lors, que l'école ouvre la voie de la réussite ?} Toutefois, ce chemin, pour royal qu'il soit, n'est pas le seul, comme le montre la réussite de certains entrepreneurs autodidactes.}

\textbf{3. Phrase facultative intégrée :} la phrase interrogative rhétorique soulignée (« Comment nier, dès lors, que l'école ouvre la voie de la réussite ? ») renforce l'argumentation en impliquant l'adhésion du lecteur.
\end{corrige}

\courssub{Étude de langue sur \emph{Le Lion et la perle} (Réception des textes)}

\begin{corrige}[Exercice 8 — Étude de langue sur \emph{Le Lion et la perle}]
\textbf{1. Classement des phrases :}
\begin{center}
\begin{tabularx}{\linewidth}{|X|l|l|}
\hline
\rowcolor{popblueL} \textbf{Phrase} & \textbf{Type} & \textbf{Forme} \\
\hline
« Ignorante fille de brousse, tu ne sais rien ! » & exclamative & négative \\
\hline
« Le monde change, il faut changer avec lui. » & déclarative & affirmative \\
\hline
« La mariée qui paie le prix de la mariée est une coutume barbare, humiliante. » & déclarative & affirmative \\
\hline
« Nous devons être modernes, civilisés. » & déclarative & affirmative \\
\hline
\end{tabularx}
\end{center}

\textbf{2. Couples de synonymes et d'antonymes :}
\begin{itemize}
\item \emph{moderne} / \emph{civilisé} : synonymes (dans la bouche de Lakunle).
\item \emph{barbare} / \emph{humiliante} : synonymes (même champ du mépris).
\item \emph{moderne} / \emph{barbare} : antonymes.
\item \emph{civilisé} / \emph{brousse} : antonymes (opposition ville/village, modernité/tradition).
\item \emph{ignorante} / \emph{civilisé} : antonymes (savoir contre ignorance).
\end{itemize}

\textbf{3. Effet produit :} l'accumulation des phrases exclamatives (« tu ne sais rien ! ») et le ton injonctif (« il faut changer », « nous devons être ») donnent au discours de Lakunle un caractère \emph{péremptoire et arrogant}. Il ne dialogue pas avec Sidi : il la sermonne et l'humilie. Ce style révèle le personnage : un intellectuel vaniteux, imbu du savoir occidental, qui méprise les traditions de son village. Wole Soyinka dénonce ainsi, par l'ironie, le faux modernisme de ceux qui singent l'Occident sans le comprendre.
\end{corrige}

\courssub{Confrontation, évaluation et amélioration des productions}

\begin{corrige}[Exercice 9 — Sujet type Baccalauréat technique : dissertation complète]
\textbf{1. Analyse du sujet :}

Termes clés : \emph{école} (institution de formation formelle), \emph{seul} (exclusivité, mot piège à discuter), \emph{réussite} (insertion professionnelle, épanouissement, ascension sociale), \emph{Cameroun} (contexte national : chômage des diplômés, vitalité du secteur informel).

Problématique : L'école est-elle la voie unique et nécessaire de la réussite, ou d'autres chemins mènent-ils aussi au succès au Cameroun ?

Plan dialectique détaillé :
\begin{itemize}
\item \textbf{I. L'école est un chemin privilégié de la réussite.} Arg. 1 : elle délivre les diplômes exigés par les emplois qualifiés (ex. : recrutements à la fonction publique camerounaise sur concours). Arg. 2 : elle forme l'esprit et ouvre sur le monde (ex. : ingénieurs formés à l'ENSP de Yaoundé).
\item \textbf{II. Mais elle n'est pas le seul chemin.} Arg. 1 : l'apprentissage informel mène au succès (ex. : maîtres-artisans prospères du quartier Mboppi à Douala). Arg. 2 : le talent et l'entrepreneuriat réussissent sans longs cursus (ex. : artistes musiciens et footballeurs camerounais de renommée internationale).
\item \textbf{III. La réussite naît de la complémentarité des voies.} Arg. 1 : la formation technique allie école et pratique (ex. : stages en entreprise des lycées techniques). Arg. 2 : la réussite durable suppose des compétences, quelle qu'en soit la source (ex. : artisans qui suivent des formations de perfectionnement).
\end{itemize}

\textbf{2. Éléments de rédaction attendus :} introduction en quatre étapes ; trois parties de deux paragraphes chacune, chaque paragraphe suivant le schéma I-A-E-L ; connecteurs logiques variés (d'abord, en effet, cependant, néanmoins, enfin, ainsi) ; conclusion avec bilan, réponse nuancée (l'école est un chemin essentiel mais non exclusif) et ouverture (ex. : la réforme des curricula vers l'employabilité).

\textbf{3. Barème indicatif (sur 20) :} analyse du sujet et problématique : 4 pts ; cohérence du plan : 4 pts ; introduction : 3 pts ; développement (I-A-E-L, exemples) : 6 pts ; conclusion : 3 pts. La correction linguistique est sanctionnée dans chaque critère.

\textbf{Points de vigilance :} ne pas transformer le sujet en catalogue (« l'école c'est bien ») ; discuter le mot \emph{seul} ; exemples précis et camerounais ; paragraphes visuellement séparés ; éviter les fautes d'accord et les phrases trop longues.
\end{corrige}

\begin{corrige}[Exercice 10 — Évaluation d'une copie anonyme — Grille officielle]
\textbf{Méthode de correction (à appliquer à la copie fournie en annexe) :}

\textbf{1. Notation critère par critère (exemple de grille remplie) :}
\begin{center}
\begin{tabularx}{\linewidth}{|X|c|X|}
\hline
\rowcolor{popblueL} \textbf{Critère} & \textbf{Note} & \textbf{Justification} \\
\hline
Analyse du sujet & 3/4 & Termes définis, problématique présente mais mal formulée (question trop large). \\
\hline
Plan & 2/4 & Plan thématique annoncé mais la troisième partie répète la première ; manque de progression. \\
\hline
Introduction & 2/3 & Amorce et problématique présentes ; annonce du plan trop détaillée (dévoile les exemples). \\
\hline
Développement & 3/6 & Arguments présents mais exemples absents dans deux paragraphes ; schéma I-A-E-L incomplet. \\
\hline
Conclusion & 2/3 & Bilan correct mais aucune ouverture. \\
\hline
Langue & 2/4 & Nombreuses fautes d'accord et de conjugaison ; connecteurs pauvres (« et », « aussi »). \\
\hline
\textbf{Total} & \textbf{14/20} & \\
\hline
\end{tabularx}
\end{center}

\textbf{2. Calcul de la note finale :} $3 + 2 + 2 + 3 + 2 + 2 = 14$. Note finale : \textbf{14/20}.

\textbf{3. Trois améliorations prioritaires :}
\begin{enumerate}
\item Rédiger chaque paragraphe selon le schéma I-A-E-L complet, avec un exemple précis et vérifiable.
\item Enrichir les connecteurs logiques (en effet, cependant, par conséquent, enfin) pour articuler la démonstration.
\item Soigner la langue : relecture systématique des accords sujet-verbe et des terminaisons verbales.
\end{enumerate}

\textbf{Point de vigilance :} le correcteur justifie toujours chaque retrait de point par une observation précise et objective, jamais par une appréciation générale (« c'est faible »).
\end{corrige}

\begin{corrige}[Exercice 11 — Synthèse : plan dialectique, introduction et conclusion]
\textbf{Sujet :} « Faut-il préserver les langues locales dans l'enseignement technique au Cameroun ? »

\textbf{1. Plan dialectique complet :}
\begin{itemize}
\item \textbf{Thèse : il faut préserver les langues locales.} Arg. 1 : elles portent l'identité culturelle (ex. : les fêtes traditionnelles où le medumba transmet savoirs et valeurs) ; connecteur : \emph{d'abord}. Arg. 2 : elles facilitent l'apprentissage des gestes techniques en milieu rural (ex. : vulgarisation agricole en fulfuldé à Maroua) ; connecteur : \emph{ensuite}.
\item \textbf{Antithèse : leur usage scolaire présente des limites.} Arg. 1 : l'extrême diversité linguistique interdit un enseignement uniforme (ex. : plus de 250 langues recensées au Cameroun) ; connecteur : \emph{cependant}. Arg. 2 : l'insertion professionnelle exige les langues officielles (ex. : concours et marchés publics rédigés en français ou en anglais) ; connecteur : \emph{de plus}.
\item \textbf{Synthèse : une complémentarité raisonnée.} Arg. 1 : langues locales en appui, langues officielles en enseignement (ex. : écoles pilotes bilingues) ; connecteur : \emph{enfin}. Arg. 2 : créer des terminologies techniques en langues nationales (ex. : lexiques élaborés avec l'appui des instituts linguistiques) ; connecteur : \emph{ainsi}.
\end{itemize}

\textbf{2. Introduction :} Le Cameroun est l'un des pays les plus riches du monde en diversité linguistique, avec plus de deux cent cinquante langues parlées sur son territoire. Or, dans les salles de classe des lycées techniques, seuls le français et l'anglais ont droit de cité. Préserver une langue, c'est la maintenir vivante dans les usages, notamment scolaires ; les langues locales désignent les idiomes nationaux transmis par les communautés. Dès lors, faut-il ouvrir l'enseignement technique aux langues locales, au risque de fragiliser l'unité et l'employabilité, ou les en exclure, au risque de les laisser mourir ? Après avoir montré la nécessité de leur préservation, nous en examinerons les limites pratiques, avant de dégager les conditions d'une complémentarité équilibrée.

\textbf{3. Conclusion :} En définitive, les langues locales sont un trésor identitaire et un levier pédagogique qu'il serait suicidaire d'abandonner ; mais leur diversité et les exigences du monde professionnel imposent que les langues officielles demeurent les langues d'enseignement. La solution réside donc dans une complémentarité : préserver les langues locales comme appui culturel et pédagogique, tout en garantissant la maîtrise du français et de l'anglais. Reste à l'État camerounais à traduire cette ambition en politique éducative concrète, en formant des enseignants capables de valoriser, dans chaque atelier et chaque salle, le génie linguistique national.

\textbf{Barème indicatif :} plan dialectique complet (arguments, exemples, connecteurs) : 8 pts ; introduction en quatre étapes : 6 pts ; conclusion (bilan, réponse, ouverture) : 6 pts.

\textbf{Points de vigilance :} exemples explicitement camerounais ; connecteurs variés et correctement employés ; l'ouverture ne doit pas introduire un nouveau sujet mais une perspective liée au sujet traité.
\end{corrige}

\newpage

\coursec{Fiche bilan}

\begin{popfigure}
\begin{tikzpicture}[mindmap, grow cyclic, every node/.style={concept, text=white, font=\small},
root concept/.append style={concept color=popdark, font=\bfseries},
level 1/.append style={level distance=3.4cm, sibling angle=60},
level 2/.append style={level distance=2.2cm, font=\scriptsize}]
\node[root concept]{Rédiger une dissertation}
  child[concept color=popblue]{ node {Analyser le sujet}
    child { node {Mots-clés} }
    child { node {Problématique} }
  }
  child[concept color=poporange]{ node {Plan détaillé}
    child { node {Thèse / antithèse / synthèse} }
  }
  child[concept color=popgreen]{ node {Introduction}
    child { node {Amorce, sujet, problématique, plan} }
  }
  child[concept color=poppurple]{ node {Paragraphe}
    child { node {Idée + argument + exemple} }
  }
  child[concept color=poppink]{ node {Conclusion}
    child { node {Bilan + ouverture} }
  }
  child[concept color=popgold]{ node {Langue}
    child { node {Types de phrase, synonymie} }
  };
\end{tikzpicture}
\legende{Carte mentale de la leçon : les six étapes et savoirs pour rédiger une dissertation réussie au Probatoire technique.}
\end{popfigure}

\begin{bilan}
\textbf{Définitions clés}
\begin{itemize}
\item \cle{Dissertation} : exercice écrit qui traite un sujet de façon \cle{argumentée}, organisée et personnelle, en mobilisant ses connaissances littéraires et culturelles.
\item \cle{Problématique} : question centrale, née de l'analyse du sujet, que tout le devoir cherche à résoudre.
\item \cle{Idée directrice} : phrase qui annonce l'idée principale d'un paragraphe.
\item \cle{Argument} : raison qui justifie l'idée ; \cle{exemple} : fait ou référence précise qui illustre l'argument.
\item \cle{Synonymie} : relation entre mots de sens proche ; \cle{antonymie} : relation entre mots de sens opposé (outils pour éviter les répétitions et nuancer).
\end{itemize}

\textbf{Structure à connaître par cœur}
\begin{center}
\begin{tabularx}{\linewidth}{|l|X|}
\hline
\rowcolor{popblueL} \textbf{Partie} & \textbf{Contenu obligatoire} \\
\hline
Introduction & Amorce (accroche) ; présentation du sujet ; problématique ; annonce du plan \\
\hline
Développement & 2 ou 3 grandes parties ; chaque paragraphe = idée directrice + argument(s) + exemple(s) + transition \\
\hline
Conclusion & Bilan (réponse à la problématique) ; élargissement (ouverture) ; pas d'idée nouvelle \\
\hline
\end{tabularx}
\end{center}

\textbf{Méthodes express}
\begin{itemize}
\item \textbf{Analyser le sujet} : souligner les mots-clés, définir chaque terme, repérer les présupposés, formuler la problématique sous forme de question.
\item \textbf{Construire le plan} : au brouillon, lister les idées, les regrouper en 2 ou 3 parties équilibrées, vérifier que chaque partie répond à la problématique.
\item \textbf{Rédiger un paragraphe} : une seule idée par paragraphe ; enchaîner idée $\rightarrow$ argument $\rightarrow$ exemple $\rightarrow$ mini-conclusion.
\item \textbf{Soigner la langue} : varier les types de phrase (déclarative, interrogative, exclamative, impérative), employer connecteurs logiques (\emph{d'abord, en effet, cependant, donc, enfin}) et synonymes.
\item \textbf{Relire} : vérifier accord, orthographe, ponctuation, cohérence des transitions.
\end{itemize}
\end{bilan}

\begin{aretenir}[Checklist avant le Bac]
Avant de rendre ma copie, je vérifie :
\begin{itemize}
\item[$\square$] J'ai bien défini tous les mots-clés du sujet au brouillon.
\item[$\square$] Ma problématique est formulée clairement (une question).
\item[$\square$] Mon introduction contient les quatre étapes : amorce, sujet, problématique, plan.
\item[$\square$] Mon plan comporte 2 ou 3 parties équilibrées qui répondent au sujet.
\item[$\square$] Chaque paragraphe développe une seule idée directrice.
\item[$\square$] Chaque idée est soutenue par au moins un argument et un exemple précis.
\item[$\square$] Des transitions relient les parties entre elles.
\item[$\square$] J'ai utilisé des connecteurs logiques variés et des synonymes.
\item[$\square$] Ma conclusion fait le bilan et propose une ouverture, sans idée nouvelle.
\item[$\square$] J'ai relu pour corriger orthographe, accords et ponctuation.
\end{itemize}
\end{aretenir}

\findecours

\end{document}
